% Réactions acide/base et dosages — Chimie Tle E — Cours signé OnBuch+
% Programme officiel MINESEC, Terminale C, D, E. Compiler avec : tectonic L09-reactions-acide-base-et-dosages.tex
\newcommand{\DOCMATIERE}{Chimie}
\newcommand{\DOCNIVEAU}{Tle E}
\newcommand{\DOCMODULE}{Module 2 · Acides et bases}
\newcommand{\DOCLECON}{09}
\newcommand{\DOCTITRE}{Réactions acide/base et dosages}
% OnBuch+ — préambule « cours signés » ENRICHI (compilé avec tectonic / XeLaTeX)
% Système visuel « Pop » d'OnBuch+ : fond crème (jamais blanc), bordures encre
% épaisses, ombres « dures » décalées, titres Archivo Black en capitales.
% Version enrichie pour les cours de chimie : chimie (mhchem, chemfig),
% graphes (pgfplots), boîtes pédagogiques supplémentaires (définition,
% propriété, méthode, attention, le savais-tu, exercice résolu, exercice,
% TP, bilan), page de garde refaite et signature OnBuch+ ancrée en bas.
%
% Macros attendues dans main.tex AVANT \input{preamble} :
%   \DOCMATIERE  \DOCNIVEAU  \DOCMODULE  \DOCLECON (numéro)  \DOCTITRE
\documentclass[11pt]{article}

\usepackage{fontspec}
\usepackage[french]{babel}
\usepackage[a4paper,margin=2cm,top=3.4cm,bottom=2.8cm,headheight=1.4cm,footskip=1.3cm]{geometry}
\usepackage{amsmath,amssymb,amsfonts}
\usepackage[table]{xcolor} % [table] : fournit aussi \rowcolors (alternance de lignes)
\usepackage{pagecolor}
\usepackage{tikz}
\usetikzlibrary{arrows.meta,positioning,calc,shapes.geometric,shapes.misc,shapes.arrows,decorations.pathmorphing,decorations.markings,patterns,shadows,mindmap,trees,fit,backgrounds,matrix,intersections}
\usepackage{pgfplots}
\pgfplotsset{compat=1.18}
\usepackage[version=4]{mhchem}
\usepackage{chemfig}
\usepackage{enumitem}
\usepackage{fancyhdr}
% [most] charge toutes les bibliothèques tcolorbox (listings, minted, raster,
% documentation...) et gonfle inutilement la mémoire TeX sur un document long
% (~300 boîtes) : on ne charge que ce qui est réellement utilisé.
\usepackage[skins,breakable]{tcolorbox}
\usepackage{graphicx}
\usepackage{colortbl}
\usepackage{booktabs}
\usepackage{tabularx}
\usepackage{array}
\usepackage{multicol}
\usepackage{siunitx}
% parse-numbers=false : accepte aussi \SI{2,5 \times 10^{-3}}{...}
\sisetup{locale=FR,output-decimal-marker={,},group-minimum-digits=4,parse-numbers=false}
\usepackage{qrcode}
% \degree : commande fréquemment utilisée par les modèles pour un angle ou une
% température en dehors de \SI{}{} (non fournie par les paquets chargés).
\providecommand{\degree}{^{\circ}}
\usepackage{lastpage}
\usepackage{hyperref}
\hypersetup{hidelinks}

% ============================================================
% Palette « Pop » OnBuch+ — mêmes hex que la classe P (plus_ui.dart)
% ============================================================
\definecolor{poporange}{HTML}{F59321}
\definecolor{popdark}{HTML}{E07A0C}
\definecolor{popcream}{HTML}{FFF6E8}
\definecolor{popink}{HTML}{1C1714}
\definecolor{poppurple}{HTML}{7A5AE0}
\definecolor{popgreen}{HTML}{2E9E6A}
\definecolor{poppink}{HTML}{E0457B}
\definecolor{popblue}{HTML}{2F7DE1}
\definecolor{popgold}{HTML}{C98A12}
\definecolor{popmuted}{HTML}{8A7F76}
\definecolor{popline}{HTML}{EADFD2}
\definecolor{popgray}{HTML}{8A7F76}
\colorlet{popgrey}{popgray}
% Alias tolérants (le modèle nomme parfois les couleurs différemment)
\colorlet{popwhite}{white}
\colorlet{popblack}{popink}
\colorlet{popred}{poppink}
\colorlet{popyellow}{popgold}
% Teintes claires pour fonds de tableaux / zones
\colorlet{popblueL}{popblue!10!white}
\colorlet{poporangeL}{poporange!14!white}
\colorlet{popgreenL}{popgreen!12!white}
\colorlet{poppurpleL}{poppurple!12!white}
\colorlet{poppinkL}{poppink!12!white}
\colorlet{popgoldL}{popgold!14!white}

\pagecolor{popcream}

% ============================================================
% Typographie
% ============================================================
\defaultfontfeatures{Ligatures=TeX}
\setmainfont{PlusJakartaSans-Medium.ttf}[Path=../../../fonts/,BoldFont=PlusJakartaSans-Bold.ttf]
% Maths en sans-serif (Fira Math) avec chiffres et lettres droites de Plus
% Jakarta, pour que les formules se fondent dans le texte.
\usepackage[math-style=ISO,bold-style=ISO]{unicode-math}
\setmathfont{FiraMath-Regular.otf}[Path=../../../fonts/]
\setmathfont{PlusJakartaSans-Medium.ttf}[Path=../../../fonts/,range={up/num,up/{Latin,latin},\mathup}]
\usepackage{icomma} % virgule décimale sans espace en mode math
% Caractères absents de la police de texte (sinon carré vide) : rendus par
% leur équivalent mathématique, en texte comme en formule.
\usepackage{newunicodechar}
\newunicodechar{α}{\ensuremath{\alpha}}
\newunicodechar{Α}{\ensuremath{\alpha}}
\newunicodechar{β}{\ensuremath{\beta}}
\newunicodechar{δ}{\ensuremath{\delta}}
\newunicodechar{✓}{\popcheck{popgreen}}
\newfontfamily\popdisplay{ArchivoBlack-Regular.ttf}[Path=../../../fonts/]
\newfontfamily\poplogo{SpaceGrotesk-Bold.ttf}[Path=../../../fonts/]
\newfontfamily\popsemi{PlusJakartaSans-SemiBold.ttf}[Path=../../../fonts/]
\newfontfamily\popxbold{PlusJakartaSans-ExtraBold.ttf}[Path=../../../fonts/]
\newfontfamily\popxboldls{PlusJakartaSans-ExtraBold.ttf}[Path=../../../fonts/,LetterSpace=3.0]

\setlist[enumerate]{leftmargin=1.7em,itemsep=5pt,topsep=4pt}
\setlist[itemize]{leftmargin=1.7em,itemsep=4pt,topsep=4pt,label={\color{poporange}\textbullet}}
\setlength{\parindent}{0pt}
\setlength{\parskip}{5pt}
\renewcommand{\arraystretch}{1.25}

% chemfig : liaisons un peu plus épaisses, encre Pop
\setchemfig{atom sep=2em,bond style={line width=0.9pt,popink},cram width=3pt}

% pgfplots : style Pop par défaut
\pgfplotsset{
  every axis/.append style={axis line style={popink,line width=1pt},tick style={popink},
    grid style={popline,line width=0.5pt},label style={font=\small},tick label style={font=\footnotesize}},
  popcurve/.style={poporange,line width=1.8pt,no markers,smooth},
}

% ============================================================
% Pill / squiggle
% ============================================================
\newcommand{\poppill}[3]{%
  \tikz[baseline=-0.55ex]\node[draw=popink,line width=1.2pt,rounded corners=2.6pt,%
    fill=#2,inner xsep=6.5pt,inner ysep=3pt]{%
    \popxboldls\fontsize{7}{7}\selectfont\color{#3}\MakeUppercase{#1}};%
}
\newcommand{\popsquiggle}[1]{%
  \tikz[baseline=-0.4ex]{
    \draw[#1,line width=2.6pt,line cap=round]
      plot [smooth] coordinates {(0,0.09) (0.34,0.01) (0.68,0.21) (1.02,0.01) (1.36,0.10)};
  }%
}
% Mot-clé surligné (marqueur orange sous le texte)
\newcommand{\cle}[1]{{\popxbold\color{popdark}#1}}

% ============================================================
% En-tête / pied
% ============================================================
\pagestyle{fancy}
\fancyhf{}
\renewcommand{\headrulewidth}{0pt}
\renewcommand{\footrulewidth}{0pt}
\lhead{\raisebox{-0.15ex}{\poplogo\fontsize{13}{13}\selectfont\color{popink}OnBuch\color{poporange}+}}
\rhead{\poppill{\DOCMATIERE\ \textbullet\ \DOCNIVEAU\ \textbullet\ Leçon \DOCLECON}{white}{popink}}
\lfoot{\popxbold\fontsize{7.6}{7.6}\selectfont\color{popmuted}\MakeUppercase{Cours signé OnBuch+}\ \textbullet\ {\popsemi\DOCTITRE}}
\rfoot{\tikz[baseline=-0.6ex]\node[rounded corners=8.5pt,fill=popink,minimum height=17pt,minimum width=17pt,inner xsep=5pt,inner sep=0]{\popxbold\fontsize{8}{8}\selectfont\color{popcream}\thepage\,/\,\pageref*{LastPage}};}

% ============================================================
% Titres
% ============================================================
\newcommand{\chaptitre}[2]{%
  \begin{tcolorbox}[enhanced,colback=poporange,colframe=popink,boxrule=2.6pt,arc=11pt,
      drop shadow=popink,left=15pt,right=15pt,top=12pt,bottom=11pt,before skip=3pt,after skip=18pt]
    \poppill{#2}{popink}{popcream}\par\vspace{9pt}
    {\raggedright\hyphenpenalty=10000\exhyphenpenalty=10000\popdisplay\fontsize{22}{24}\selectfont\color{popcream}\MakeUppercase{#1}\par}
  \end{tcolorbox}
}
% Section numérotée : pastille orange avec le numéro + titre Archivo
\newcounter{popsec}
\newcommand{\coursec}[1]{%
  \stepcounter{popsec}\setcounter{popsub}{0}%
  \par\vspace{16pt}\noindent
  \begin{minipage}{\linewidth}
  \tikz[baseline=-0.7ex]\node[draw=popink,line width=1.6pt,fill=poporange,rounded corners=4pt,minimum size=22pt,inner sep=0]{\popdisplay\fontsize{12}{12}\selectfont\color{popcream}\thepopsec};%
  \hspace{8pt}{\popdisplay\fontsize{15}{17}\selectfont\color{popink}\MakeUppercase{#1}}\par
  \vspace{4pt}\noindent{\color{poporange}\rule{36pt}{3.6pt}}
  \end{minipage}\par\nopagebreak\vspace{8pt}
}
\newcounter{popsub}
\newcommand{\courssub}[1]{%
  \stepcounter{popsub}%
  \par\vspace{9pt}\noindent{\popxbold\fontsize{11.5}{13}\selectfont\color{popink}{\color{poporange}\thepopsec.\thepopsub}\hspace{6pt}#1}\par\nopagebreak\vspace{4pt}
}

% ============================================================
% Boîtes pédagogiques — même recette : fond blanc, bordure épaisse colorée,
% ombre dure encre, badge plein. Titre optionnel [#1].
% ============================================================
\tcbset{popbase/.style={breakable,enhanced,colback=white,boxrule=2.3pt,arc=9pt,
  shadow={3pt}{-3pt}{0pt}{popink},left=14pt,right=14pt,top=11pt,bottom=11pt,before skip=11pt,after skip=15pt,
  fontupper=\popsemi\fontsize{10.6}{14.5}\selectfont\color{popink},
  fontlower=\popsemi\fontsize{10.6}{14.5}\selectfont\color{popink}}}
% Coche dessinée (le glyphe ✓ n'existe pas dans les polices du document)
\newcommand{\popcheck}[1]{\tikz[baseline=-0.5ex]\draw[#1,line width=1.6pt,line cap=round,line join=round](-0.13,0.0)--(-0.04,-0.09)--(0.13,0.1);}
\newcommand{\popboxhead}[3]{\poppill{#1}{#2}{white}\ifx\relax#3\relax\else\hspace{6pt}{\popxbold\fontsize{10.6}{12}\selectfont\color{popink}#3}\fi\par\vspace{7pt}}

\newtcolorbox{aretenir}[1][]{popbase,colframe=popgreen,before upper={\popboxhead{À retenir}{popgreen}{#1}}}
\newtcolorbox{exemplebox}[1][]{popbase,colframe=poporange,before upper={\popboxhead{Exemple}{poporange}{#1}}}
\newtcolorbox{definition}[1][]{popbase,colframe=popblue,before upper={\popboxhead{Définition}{popblue}{#1}}}
\newtcolorbox{propriete}[1][]{popbase,colframe=poppurple,before upper={\popboxhead{Propriété}{poppurple}{#1}}}
\newtcolorbox{methode}[1][]{popbase,colframe=popgold,before upper={\popboxhead{Méthode}{popgold}{#1}}}
\newtcolorbox{attention}[1][]{popbase,colframe=poppink,before upper={\popboxhead{Attention}{poppink}{#1}}}
\newtcolorbox{experience}[1][]{popbase,colframe=popblue,colback=popblueL,before upper={\popboxhead{Expérience}{popblue}{#1}}}
% « Le savais-tu ? » : carte violette pleine (contexte, culture, Cameroun)
\newtcolorbox{savaistu}[1][]{popbase,colback=poppurple,colframe=popink,
  fontupper=\popsemi\fontsize{10.4}{14.2}\selectfont\color{white},
  before upper={\poppill{Le savais-tu\,?}{white}{poppurple}\ifx\relax#1\relax\else\hspace{6pt}{\popxbold\color{white}#1}\fi\par\vspace{7pt}}}
% Exercice résolu : énoncé en haut, \tcblower puis la solution sur fond crème
\newcounter{popexo}\newcounter{popexr}
\newtcolorbox{exoresolu}[1][]{popbase,colframe=poporange,colbacklower=popcream,
  segmentation style={popink,line width=1.2pt,dashed},
  before upper={\stepcounter{popexr}\popboxhead{Exercice résolu \thepopexr}{poporange}{#1}},
  before lower={\poppill{Solution}{popink}{popcream}\par\vspace{6pt}}}
% Exercice (non résolu en place) — la correction est dans la partie Corrigés
\newtcolorbox{exercice}[1][]{popbase,colframe=popink,
  before upper={\stepcounter{popexo}\popboxhead{Exercice \thepopexo}{popink}{#1}}}
% Corrigé
\newtcolorbox{corrige}[1][]{popbase,colframe=popgreen,colback=popgreenL,
  before upper={\popboxhead{Corrigé}{popgreen}{#1}}}
% Objectifs / prérequis / bilan
\newtcolorbox{objectifs}{popbase,colframe=popink,colback=poporangeL,
  before upper={\popboxhead{Objectifs de la leçon}{popink}{}}}
\newtcolorbox{prerequis}{popbase,colframe=popink,colback=popblueL,
  before upper={\popboxhead{Prérequis}{popblue}{}}}
\newtcolorbox{bilan}{popbase,colframe=popink,colback=white,boxrule=2.8pt,
  before upper={\popboxhead{Fiche bilan}{poporange}{L'essentiel à connaître pour l'examen}}}
% Figure encadrée avec légende
\newtcolorbox{popfigure}[1][]{enhanced,breakable=false,colback=white,colframe=popline,boxrule=1.4pt,arc=8pt,
  left=10pt,right=10pt,top=10pt,bottom=8pt,before skip=10pt,after skip=12pt,halign=center,
  lower separated=false,halign lower=center,
  fontlower=\popsemi\fontsize{9}{11.5}\selectfont\color{popmuted},
  #1}
\newcommand{\legende}[1]{\par\vspace{6pt}{\popsemi\fontsize{9}{11.5}\selectfont\color{popmuted}#1}}

% ============================================================
% Signature OnBuch+
% ============================================================
\newcommand{\poplogoinline}[1]{{\poplogo\fontsize{#1}{#1}\selectfont\color{popink}OnBuch\color{poporange}+}}
\newcommand{\signatureonbuch}{%
  \begin{tcolorbox}[enhanced,colback=popink,colframe=popink,boxrule=2.2pt,arc=10pt,
      drop shadow=poporange,left=14pt,right=14pt,top=10pt,bottom=10pt,before skip=0pt,after skip=0pt]
    \tikz[baseline=-0.7ex]\node[circle,fill=popgreen,draw=popcream,line width=1.3pt,minimum size=22pt,inner sep=0]{\popcheck{white}};%
    \hspace{9pt}\begin{minipage}[c]{0.62\linewidth}
      {\popxbold\fontsize{12}{13}\selectfont\color{popcream}Signé\ \textbullet\ {\poplogo OnBuch\color{poporange}+}}\par\vspace{2pt}
      {\raggedright\popsemi\fontsize{8}{10.5}\selectfont\color{popline}\MakeUppercase{Équipe pédagogique OnBuch+}\\\MakeUppercase{Conforme au programme officiel MINESEC}\par}
    \end{minipage}\hfill
    {\poplogo\fontsize{22}{22}\selectfont\color{popcream}OnBuch\color{poporange}+}
  \end{tcolorbox}%
}

% ============================================================
% Page de garde
% #1 titre · #2 matière · #3 niveau · #4 module · #5 n° leçon · #6 durée ·
% #7 description / objectifs (texte ou itemize)
% ============================================================
\newcommand{\pagegarde}[7]{%
  \thispagestyle{empty}
  \newgeometry{a4paper,margin=1.7cm,top=1.6cm,bottom=1.4cm}%
  \begin{tikzpicture}[remember picture,overlay]
    \fill[poporange!18] ([xshift=-3.2cm,yshift=-3.6cm]current page.north east) circle (4.6cm);
    \fill[poppurple!14] ([xshift=2.4cm,yshift=7.6cm]current page.south west) circle (3.8cm);
  \end{tikzpicture}%
  {\poplogo\fontsize{30}{30}\selectfont\color{popink}OnBuch\color{poporange}+}\hfill
  \raisebox{0.6ex}{\poppill{Cours signé \textbullet\ Premium}{popink}{popcream}}\par
  \vspace{1pt}\popsquiggle{poppurple}\par
  \vspace{8pt}
  \begin{tcolorbox}[enhanced,colback=poporange,colframe=popink,boxrule=2.8pt,arc=13pt,
      drop shadow=popink,left=18pt,right=18pt,top=12pt,bottom=13pt,after skip=10pt]
    \poppill{#2 \textbullet\ #3}{popink}{popcream}\hspace{5pt}\poppill{#4}{white}{popink}\par\vspace{8pt}
    {\popdisplay\fontsize{14}{14}\selectfont\color{popink}LEÇON #5}\par\vspace{4pt}
    {\raggedright\hyphenpenalty=10000\exhyphenpenalty=10000\popdisplay\fontsize{25}{27}\selectfont\color{popcream}\MakeUppercase{#1}\par}
    \vspace{8pt}
    {\popxbold\fontsize{9.5}{9.5}\selectfont\color{popink}\MakeUppercase{Durée conseillée : #6}}
  \end{tcolorbox}
  \begin{tcolorbox}[enhanced,colback=white,colframe=popink,boxrule=2.2pt,arc=10pt,
      drop shadow=popink,left=16pt,right=16pt,top=10pt,bottom=10pt,after skip=8pt]
    \poppill{Dans cette leçon}{poppurple}{white}\par\vspace{5pt}
    \popsemi\fontsize{9.3}{12.6}\selectfont\color{popink}#7
  \end{tcolorbox}
  \noindent\begin{minipage}[t]{0.487\linewidth}
    \begin{tcolorbox}[enhanced,colback=popgreen,colframe=popink,boxrule=2.2pt,arc=10pt,
        drop shadow=popink,left=11pt,right=11pt,top=10pt,bottom=10pt,height=3.1cm,valign=center]
      \tcbox[colback=white,colframe=popink,boxrule=1.3pt,arc=5pt,boxsep=0pt,nobeforeafter,
        left=3pt,right=3pt,top=3pt,bottom=3pt,valign=center]{\qrcode[height=1.9cm]{https://play.google.com/store/apps/details?id=com.luvvix.onbuch&hl=fr}}%
      \hspace{8pt}\begin{minipage}[c]{0.5\linewidth}
        {\popdisplay\fontsize{11}{12.5}\selectfont\color{white}\MakeUppercase{Télécharge OnBuch}}\par\vspace{4pt}
        {\raggedright\popsemi\fontsize{8.4}{11}\selectfont\color{white}Gratuit \textbullet\ Google Play\\Tuteur IA, annales, concours, résultats\par}
      \end{minipage}
    \end{tcolorbox}
  \end{minipage}\hfill
  \begin{minipage}[t]{0.487\linewidth}
    \begin{tcolorbox}[enhanced,colback=poppurple,colframe=popink,boxrule=2.2pt,arc=10pt,
        drop shadow=popink,left=11pt,right=11pt,top=10pt,bottom=10pt,height=3.1cm,valign=center]
      \tikz[baseline=-0.7ex]\node[circle,fill=white,draw=popink,line width=1.3pt,minimum size=1.6cm,inner sep=0]{\popdisplay\fontsize{24}{24}\selectfont\color{poppurple}+};%
      \hspace{8pt}\begin{minipage}[c]{0.6\linewidth}
        {\popdisplay\fontsize{11}{12.5}\selectfont\color{white}\MakeUppercase{Passe à OnBuch+}}\par\vspace{4pt}
        {\raggedright\popsemi\fontsize{8.4}{11}\selectfont\color{white}Tous les cours signés, Tuteur IA illimité, fiches PDF — onglet OnBuch+ de l'app\par}
      \end{minipage}
    \end{tcolorbox}
  \end{minipage}
  \vspace{8pt}
  \popcontenu
  \vfill
  \signatureonbuch
  \restoregeometry
  \newpage
}

% Rangée « Ce cours contient » (page de garde)
\newcommand{\poptile}[3]{% #1 couleur #2 gros texte #3 légende
  \begin{tcolorbox}[enhanced,colback=#1,colframe=popink,boxrule=2pt,arc=9pt,shadow={2.5pt}{-2.5pt}{0pt}{popink},
      left=6pt,right=6pt,top=9pt,bottom=9pt,halign=center,valign=center,height=1.75cm,width=\linewidth,nobeforeafter]
    {\popdisplay\fontsize{12.5}{13}\selectfont\color{white}\MakeUppercase{#2}}\par\vspace{3pt}
    {\centering\popsemi\fontsize{7.8}{9.5}\selectfont\color{white}#3\par}
  \end{tcolorbox}}
\newcommand{\popcontenu}{%
  \par\noindent{\popxboldls\fontsize{7.5}{7.5}\selectfont\color{popmuted}CE COURS SIGNÉ CONTIENT}\par\vspace{7pt}
  \noindent\begin{minipage}[t]{0.235\linewidth}\poptile{popblue}{Cours}{complet, illustré, pas à pas}\end{minipage}\hfill
  \begin{minipage}[t]{0.235\linewidth}\poptile{poporange}{Méthodes}{et exercices résolus}\end{minipage}\hfill
  \begin{minipage}[t]{0.235\linewidth}\poptile{poppink}{Exercices}{type examen + corrigés}\end{minipage}\hfill
  \begin{minipage}[t]{0.235\linewidth}\poptile{popgreen}{Fiche bilan}{l'essentiel à retenir}\end{minipage}\par}

% Sommaire
\newtcolorbox{sommaire}{enhanced,colback=white,colframe=popink,boxrule=2.2pt,arc=10pt,
  drop shadow=popink,left=18pt,right=18pt,top=13pt,bottom=13pt,before skip=6pt,after skip=20pt,
  before upper={\poppill{Sommaire}{poppurple}{white}\par\vspace{9pt}\popsemi\fontsize{10.6}{14.5}\selectfont\color{popink}}}

% Signature de fin de cours — poussée tout en bas de la dernière page
\newcommand{\findecours}{%
  \par\vfill\nopagebreak
  \begin{center}\popsquiggle{poporange}\end{center}\vspace{4pt}
  \signatureonbuch
}

%% FIGFIX-BEGIN v4 — correctifs globaux des figures (docs/onbuch-generation/tools/figfix)
\usepackage{adjustbox}\usepackage{etoolbox}
% toute figure TikZ/pgfplots plus large que la ligne est réduite à la largeur disponible
\BeforeBeginEnvironment{tikzpicture}{\begin{adjustbox}{max width=\linewidth}}
\AfterEndEnvironment{tikzpicture}{\end{adjustbox}}
% légendes pgfplots lisibles par-dessus les courbes
\pgfplotsset{every axis/.append style={legend style={fill=white,fill opacity=0.92,text opacity=1,draw=black!15,inner sep=3pt}}}
% étiquettes d'arêtes (syntaxe edge["..."]) avec fond blanc
\tikzset{every edge quotes/.append style={fill=white,inner sep=1.2pt}}
%% FIGFIX-END
\begin{document}

\pagegarde{Réactions acide/base et dosages}{Chimie}{Tle E}{Module 2 · Acides et bases}{09}{5 h (cours + TP)}{Cette leçon étudie les réactions acido-basiques quantitatives et leur exploitation dans les dosages pH-métriques et colorimétriques. Elle fournit les outils théoriques et expérimentaux pour déterminer une concentration inconnue, vérifier l'étiquette d'un produit commercial et choisir judicieusement un indicateur coloré.
\vspace{4pt}
\begin{multicols}{2}\small
\begin{itemize}
\item À la fin de cette leçon, je sais écrire l'équation-bilan de la réaction de dosage entre un acide et une base et justifier ses caractéristiques (totale, rapide, exothermique)
\item À la fin de cette leçon, je sais définir l'équivalence acido-basique et exploiter la relation CA·VA = CB·VBE
\item À la fin de cette leçon, je sais tracer et interpréter une courbe pH = f(VB) pour un dosage acide fort–base forte
\item À la fin de cette leçon, je sais déterminer graphiquement le point d'équivalence par la méthode des tangentes parallèles et par la méthode de la dérivée
\item À la fin de cette leçon, je sais interpréter les courbes de dosage d'un acide faible par une base forte (demi-équivalence, pH = pKa, pHE > 7) et d'un acide fort par une base faible (pHE < 7)
\item À la fin de cette leçon, je sais choisir un indicateur coloré adapté à un dosage donné
\end{itemize}\end{multicols}}

\begin{sommaire}
\begin{multicols}{2}
\begin{enumerate}
\item Réaction acide fort – base forte : étude et équivalence
\item Dosage pH-métrique d'un acide fort par une base forte
\item Dosage d'un acide faible par une base forte
\item Dosage d'un acide fort par une base faible
\item Dosage colorimétrique et choix de l'indicateur
\item Applications : vérification d'étiquettes et degré d'acidité
\item Travaux pratiques
\item Méthodes et exercices résolus
\item Exercices
\item Corrigés des exercices
\item Fiche bilan
\end{enumerate}
\end{multicols}
\end{sommaire}

\begin{objectifs}
\begin{itemize}
\item À la fin de cette leçon, je sais écrire l'équation-bilan de la réaction de dosage entre un acide et une base et justifier ses caractéristiques (totale, rapide, exothermique)
\item À la fin de cette leçon, je sais définir l'équivalence acido-basique et exploiter la relation CA·VA = CB·VBE
\item À la fin de cette leçon, je sais tracer et interpréter une courbe pH = f(VB) pour un dosage acide fort–base forte
\item À la fin de cette leçon, je sais déterminer graphiquement le point d'équivalence par la méthode des tangentes parallèles et par la méthode de la dérivée
\item À la fin de cette leçon, je sais interpréter les courbes de dosage d'un acide faible par une base forte (demi-équivalence, pH = pKa, pHE > 7) et d'un acide fort par une base faible (pHE < 7)
\item À la fin de cette leçon, je sais choisir un indicateur coloré adapté à un dosage donné
\item À la fin de cette leçon, je sais réaliser un dosage pH-métrique et un dosage colorimétrique pour vérifier l'étiquette d'une solution commerciale (vinaigre, détergent)
\item À la fin de cette leçon, je sais calculer une concentration inconnue et le degré d'acidité d'un vinaigre à partir des résultats d'un dosage
\end{itemize}
\end{objectifs}

\begin{prerequis}
\begin{itemize}
\item Couples acide/base, définitions de Brønsted, pKa et Ka (leçons précédentes du module)
\item Définition du pH, relation pH = -log[H3O+] et produit ionique de l'eau Ke
\item Réactions acido-basiques et constante d'équilibre (notion de réaction totale)
\item Quantité de matière, concentration molaire, dilution (Première)
\item Manipulation de la verrerie de précision : pipette jaugée, burette graduée (TP de Première)
\end{itemize}
\end{prerequis}

\newpage

\coursec{Réaction acide fort – base forte : étude et équivalence}

\textbf{Situation d'accroche.}\par
Au marché Mokolo à Yaoundé, Aminatou hésite devant un bidon de vinaigre étiqueté « \SI{8}{\degree} » : ce degré d'acidité suffira-t-il pour sa marinade de poisson braisé ? Son fils, élève en Terminale C, sourit : « Maman, on peut vérifier ça au laboratoire ! » Quelques kilomètres plus loin, à Douala, un technicien d'un laboratoire de contrôle se pose exactement la même question devant un bidon d'acide chlorhydrique destiné au détartrage : la concentration annoncée sur l'étiquette est-elle exacte ? La réponse tient en un mot : le \cle{dosage}. Avec une simple solution de soude de concentration connue, un pH-mètre ou quelques gouttes d'indicateur coloré, on peut « compter » les ions \ce{H3O+} présents dans une solution et remonter à sa concentration. Cette leçon te révèle toute la puissance de cette méthode, en commençant par le cas le plus simple et le plus fondamental : la réaction entre un \cle{acide fort} et une \cle{base forte}.

\textbf{Problématique.} Que se passe-t-il, à l'échelle microscopique et macroscopique, lorsqu'on verse une solution de soude dans une solution d'acide chlorhydrique ? Comment définir et repérer le moment précis — l'\cle{équivalence} — où l'acide est exactement consommé, et comment en déduire la concentration inconnue ?

\courssub{Rappels : acides forts et bases forts dans l'eau}

Avant de faire réagir, rappelons qui sont les acteurs. Un \cle{acide fort} est un acide qui réagit \textbf{totalement} avec l'eau : chaque molécule d'acide cède son proton à une molécule d'eau. C'est le cas du chlorure d'hydrogène \ce{HCl} :
\[ \ce{HCl + H2O -> H3O+ + Cl-} \]
La transformation est totale : dans une solution d'acide chlorhydrique, il ne reste \emph{aucune} molécule \ce{HCl} ; la solution contient uniquement des ions \ce{H3O+} (responsables de l'acidité) et des ions \ce{Cl-} (indifférents sur le plan acido-basique). On écrit donc simplement : solution de \ce{(H3O+ + Cl-)}.

De même, une \cle{base forte} se dissocie totalement dans l'eau. L'hydroxyde de sodium (soude) est un solide ionique qui se dissout intégralement :
\[ \ce{NaOH ->[eau] Na+ + OH-} \]
La solution de soude contient des ions \ce{OH-} (responsables de la basicité) et des ions \ce{Na+} (spectateurs).

\begin{aretenir}[Les deux solutions de départ]
\begin{itemize}
\item Acide chlorhydrique : \ce{(H3O+ + Cl-)} ; l'espèce acide active est l'ion \ce{H3O+}.
\item Soude : \ce{(Na+ + OH-)} ; l'espèce basique active est l'ion \ce{OH-}.
\end{itemize}
Tout dosage acide fort -- base forte est donc, au fond, un combat entre deux seuls protagonistes : \ce{H3O+} contre \ce{OH-}.
\end{aretenir}

\courssub{Étude de la réaction entre un acide fort et une base forte}

\textbf{Intuition.} Verse goutte à goutte de la soude dans un bécher d'acide chlorhydrique. Chaque ion \ce{OH-} qui tombe dans le bécher rencontre immédiatement un ion \ce{H3O+} et lui « arrache » un proton : il naît deux molécules d'eau. Les ions \ce{Na+} et \ce{Cl-}, eux, restent en solution sans participer à la fête : ce sont des \cle{ions spectateurs}.

\begin{experience}[Mélange acide chlorhydrique + soude]
\textbf{Dispositif et protocole.} Dans un bécher contenant \SI{20,0}{mL} d'acide chlorhydrique à \SI{0,10}{mol.L^{-1}}, on verse progressivement une solution de soude à \SI{0,10}{mol.L^{-1}} à l'aide d'une burette. Un thermomètre plonge dans le mélange, un pH-mètre suit le pH.\par
\textbf{Observations.} (1) Le pH augmente au fur et à mesure du versement. (2) La température du mélange s'élève nettement (de quelques degrés). (3) Aucun dégagement gazeux, aucun précipité, aucun changement de couleur : la réaction est invisible à l'œil nu, d'où la nécessité d'un suivi (pH-mètre ou indicateur).\par
\textbf{Interprétation.} Les ions \ce{H3O+} et \ce{OH-} réagissent entre eux pour former de l'eau ; la réaction dégage de la chaleur.
\end{experience}

L'équation-bilan de la réaction s'obtient en écrivant toutes les espèces présentes, puis en éliminant les ions spectateurs :
\[ \ce{(H3O+ + Cl-) + (Na+ + OH-) -> 2 H2O + (Na+ + Cl-)} \]
soit, après simplification, la seule équation chimiquement significative :
\[ \ce{H3O+ + OH- -> 2 H2O} \]

\begin{popfigure}
\begin{center}
\begin{tikzpicture}
\node[draw=popblue, rounded corners, fill=popblueL, inner sep=6pt, align=center] (r) {\ce{H3O+} + \ce{OH-} -> \ce{2 H2O}};
\node[below=0.35cm of r, text=popmuted, align=center, font=\small] {réaction réelle (seuls les ions actifs)};
\node[above=0.5cm of r, align=center] {$\ce{(H3O+ + Cl-)} + \ce{(Na+ + OH-)}$};
\node[above=1.2cm of r, text=poporange, font=\small] {ions spectateurs \ce{Na+} et \ce{Cl-} : inchangés};
\end{tikzpicture}
\end{center}
\legende{Les ions \ce{Na+} et \ce{Cl-} n'interviennent pas dans la réaction : ce sont des spectateurs.}
\end{popfigure}

\begin{propriete}[Caractéristiques de la réaction \ce{H3O+ + OH- -> 2 H2O}]
\begin{enumerate}
\item \textbf{Totale.} Sa constante d'équilibre vaut $K = \dfrac{1}{K_e} = \dfrac{1}{10^{-14}} = 10^{14} \gg 10^{4}$ : la réaction peut être considérée comme totale. Au moins l'un des deux réactifs est entièrement consommé.
\item \textbf{Rapide.} Elle est quasi instantanée : dès que la goutte de soude tombe, la réaction est terminée.
\item \textbf{Exothermique.} Elle dégage de la chaleur : la température du mélange s'élève pendant le dosage.
\item \textbf{Unique.} C'est la seule réaction qui se produit : aucune réaction parasite ne vient fausser le bilan de matière.
\end{enumerate}
\end{propriete}

\begin{popfigure}
\begin{center}
\begin{tikzpicture}[scale=0.9]
% thermomètre
\draw[thick, popdark] (0,0) circle (0.45);
\draw[thick, popdark] (-0.18,0.35) -- (-0.18,3.4) arc (180:0:0.18) -- (0.18,0.35);
\fill[poporange] (0,0) circle (0.32);
\fill[poporange] (-0.09,0) rectangle (0.09,2.6);
\foreach \y in {0.6,1.2,1.8,2.4,3.0} \draw[popdark] (0.18,\y) -- (0.38,\y);
\node[right] at (0.45,0.6) {\small $\theta_0$};
\node[right] at (0.45,2.6) {\small $\theta_f > \theta_0$};
\draw[-{Stealth}, poporange, thick] (1.6,0.6) -- (1.6,2.6) node[midway, right, align=left, font=\small] {élévation de\\température};
\node[align=center, font=\small, text=popmuted] at (0,-1.0) {réaction exothermique\\$\Delta H \approx -\SI{57}{kJ.mol^{-1}}$};
\end{tikzpicture}
\end{center}
\legende{La réaction \ce{H3O+ + OH- -> 2 H2O} dégage de la chaleur : le mélange se réchauffe pendant le dosage.}
\end{popfigure}

\begin{savaistu}[Pourquoi verser l'acide dans l'eau, et jamais l'inverse ?]
La réaction \ce{H3O+ + OH- -> 2 H2O} dégage environ \SI{57}{kJ} par mole d'eau formée. La dissolution d'un acide concentré dans l'eau est elle aussi très exothermique. C'est pourquoi la règle de sécurité impose de verser \textbf{l'acide dans l'eau} (et jamais l'eau dans l'acide) : la grande masse d'eau absorbe la chaleur dégagée et évite les projections bouillantes. Retiens la formule du laboratoire : « \emph{l'acide fait un plongeon dans l'eau} ».
\end{savaistu}

\courssub{La réaction de dosage}

\begin{definition}[Réaction de dosage]
Un \cle{dosage} est une méthode d'analyse qui consiste à déterminer la concentration inconnue d'une espèce en solution (l'\cle{espèce titrée}) en la faisant réagir avec une espèce de concentration connue (l'\cle{espèce titrante}). La réaction utilisée, dite \cle{réaction de dosage}, doit être : \textbf{totale}, \textbf{rapide} et \textbf{unique}.
\end{definition}

La réaction \ce{H3O+ + OH- -> 2 H2O} remplit parfaitement ces trois conditions : c'est la réaction de dosage par excellence. Le dispositif expérimental classique est représenté ci-dessous.

\begin{popfigure}
\begin{center}
\begin{tikzpicture}[scale=0.95]
% burette
\draw[thick, popdark] (-0.25,5.6) -- (-0.25,2.6) -- (-0.06,2.2) -- (-0.06,1.9);
\draw[thick, popdark] (0.25,5.6) -- (0.25,2.6) -- (0.06,2.2) -- (0.06,1.9);
\fill[popblueL] (-0.22,5.5) rectangle (0.22,3.4);
\foreach \y in {2.9,3.3,3.7,4.1,4.5,4.9,5.3} \draw[popdark] (-0.25,\y) -- (-0.10,\y);
\draw[thick, popdark] (-0.25,2.15) -- (0.25,2.15);
\node[right, align=left, font=\small] at (0.45,4.4) {burette graduée\\soude \ce{(Na+ + OH-)}\\à $C_B$ connue};
% goutte
\fill[popblue] (0,1.55) circle (0.07);
% support
\draw[thick, popmuted] (-1.4,0) -- (-1.4,5.8);
\draw[thick, popmuted] (-1.4,5.6) -- (-0.25,5.6);
% bécher
\draw[thick, popdark] (-1.1,0) -- (-1.1,1.5) (1.1,0) -- (1.1,1.5);
\draw[thick, popdark] (-1.1,0) -- (1.1,0);
\fill[popgreenL] (-1.05,0.03) rectangle (1.05,0.95);
\node[below, align=center, font=\small] at (0,-0.15) {bécher : acide \ce{(H3O+ + Cl-)}\\volume $V_A$ prélevé à la pipette};
% barreau aimanté
\fill[popgold] (-0.35,0.10) rectangle (0.35,0.28);
\node[font=\scriptsize] at (0,0.19) {\textcolor{white}{N\;S}};
% agitateur
\draw[thick, popmuted] (-1.6,-0.35) rectangle (1.6,-0.85);
\node[font=\small, text=white] at (0,-0.6) {agitateur magnétique};
% pH-mètre
\draw[thick, poppurple] (2.6,2.2) rectangle (3.8,3.0);
\node[font=\small, text=poppurple] at (3.2,2.6) {pH-mètre};
\draw[thick, poppurple] (2.9,2.2) -- (2.9,0.9);
\draw[thick, poppurple] (3.3,2.2) -- (3.3,0.9);
\fill[poppurpleL] (2.82,0.7) rectangle (2.98,1.0);
\fill[poppurpleL] (3.22,0.7) rectangle (3.38,1.0);
\end{tikzpicture}
\end{center}
\legende{Dispositif de dosage pH-métrique d'un acide fort par une base forte.}
\end{popfigure}

\courssub{L'équivalence acido-basique}

\textbf{Intuition.} Imagine que les ions \ce{H3O+} sont des invités et les ions \ce{OH-} des chaises. Tant qu'il reste des invités debout, la solution est acide. Quand chaque invité a exactement une chaise, personne ne reste debout, aucune chaise n'est libre : c'est l'\cle{équivalence}. Si l'on continue à ajouter des chaises, elles s'accumulent inutilement : la solution devient basique.

\begin{definition}[Équivalence acido-basique]
L'\cle{équivalence} est l'état du mélange où les réactifs ont été introduits et ont réagi dans les \textbf{proportions stœchiométriques} de l'équation de dosage. Pour la réaction \ce{H3O+ + OH- -> 2 H2O}, les coefficients sont tous égaux à 1, donc à l'équivalence :
\[ n(\ce{H3O+})_{\text{initial}} = n(\ce{OH-})_{\text{versé}} \]
\end{definition}

Traduisons cette égalité avec les concentrations et les volumes. On prélève un volume $V_A$ d'acide de concentration inconnue $C_A$ ; on verse un volume $V_{BE}$ de base de concentration connue $C_B$ pour atteindre l'équivalence. Comme $n = C \times V$ :

\begin{aretenir}[Relation à l'équivalence]
\[ \boxed{\; C_A \cdot V_A = C_B \cdot V_{BE} \;} \]
avec $C_A$ et $C_B$ en \si{mol.L^{-1}}, $V_A$ et $V_{BE}$ \textbf{dans la même unité de volume} (les deux en litres ou les deux en millilitres).
\end{aretenir}

\textbf{Que vaut le pH à l'équivalence ?} À l'équivalence, tous les ions \ce{H3O+} et tous les ions \ce{OH-} versés ont disparu pour former de l'eau. Il ne reste en solution que de l'eau et les ions spectateurs \ce{Na+} et \ce{Cl-} : c'est exactement une solution de chlorure de sodium, comme de l'eau salée pour la cuisine ! Or \ce{Na+} et \ce{Cl-} sont indifférents vis-à-vis de l'eau. La solution est donc \textbf{neutre} :

\begin{propriete}[pH à l'équivalence, cas acide fort -- base forte]
Pour le dosage d'un acide fort par une base forte (ou l'inverse), à \SI{25}{\celsius} :
\[ \mathrm{pH}_E = 7 \]
\end{propriete}

\begin{attention}[Erreur fréquente : le mauvais volume]
Dans la relation $C_A \cdot V_A = C_B \cdot V_{BE}$, le volume de base à utiliser est \textbf{uniquement} le volume versé \emph{à l'équivalence}, $V_{BE}$, repéré sur la courbe ou au changement de couleur de l'indicateur. Utiliser le volume total versé pendant toute l'expérience (par exemple \SI{25}{mL} si l'on a continué à verser après l'équivalence) fausse complètement le calcul. Autre piège : $V_A$ est le volume \emph{prélevé à la pipette jaugée} dans le bécher, pas le volume de la solution mère ni le volume total du mélange.
\end{attention}

\courssub{Exploitation quantitative : détermination d'une concentration inconnue}

\begin{methode}[Exploiter la relation d'équivalence $C_A \cdot V_A = C_B \cdot V_{BE}$]
\begin{enumerate}
\item \textbf{Identifier} l'espèce titrée (concentration inconnue, dans le bécher) et l'espèce titrante (concentration connue, dans la burette).
\item \textbf{Écrire} l'équation de la réaction de dosage et vérifier la stœchiométrie (ici 1 pour 1).
\item \textbf{Relever} le volume à l'équivalence $V_{BE}$ (sur la courbe pH-métrique ou au virage de l'indicateur) et le volume prélevé $V_A$.
\item \textbf{Convertir} les volumes dans la même unité (par exemple en litres) \textbf{si elles sont différentes} ; si les deux volumes sont exprimés dans la même unité (mL ou L), le rapport est sans dimension et la conversion est inutile.
\item \textbf{Isoler} l'inconnue : $C_A = \dfrac{C_B \cdot V_{BE}}{V_A}$, puis calculer et donner le résultat avec un nombre cohérent de chiffres significatifs et son unité.
\end{enumerate}
\end{methode}

\begin{exoresolu}[Dosage de l'acide chlorhydrique par la soude]
Un technicien de laboratoire à Douala prélève $V_A = \SI{20,0}{mL}$ d'une solution commerciale d'acide chlorhydrique diluée. Il dose ce prélèvement par une solution de soude de concentration $C_B = \SI{0,10}{mol.L^{-1}}$. Le suivi pH-métrique montre un saut de pH centré à $\mathrm{pH}_E = 7$ pour un volume versé $V_{BE} = \SI{15,0}{mL}$. Déterminer la concentration $C_A$ de la solution d'acide chlorhydrique.
\tcblower
\textbf{Étape 1 — Équation de la réaction de dosage.}
\[ \ce{H3O+ + OH- -> 2 H2O} \quad \text{(totale, rapide, exothermique, stœchiométrie 1/1)} \]
\textbf{Étape 2 — Condition d'équivalence.} À l'équivalence, les réactifs ont été mélangés dans les proportions stœchiométriques :
\[ n(\ce{H3O+})_{\text{initial}} = n(\ce{OH-})_{\text{versé}} \quad \Longrightarrow \quad C_A \cdot V_A = C_B \cdot V_{BE} \]
\textbf{Étape 3 — Application numérique.} Les deux volumes sont en millilitres (même unité, inutile de convertir) :
\[ C_A = \frac{C_B \cdot V_{BE}}{V_A} = \frac{0{,}10 \times 15{,}0}{20{,}0} = \SI{7,5e-2}{mol.L^{-1}} \]
\textbf{Étape 4 — Conclusion.} La solution d'acide chlorhydrique diluée a pour concentration $C_A = \SI{7,5e-2}{mol.L^{-1}}$ (deux chiffres significatifs, limités par $C_B = \SI{0,10}{mol.L^{-1}}$). Vérification de cohérence : $V_{BE} < V_A$, donc $C_A < C_B$ : l'acide est plus dilué que la soude, ce qui est logique puisqu'il a fallu moins de \SI{15,0}{mL} de soude pour neutraliser \SI{20,0}{mL} d'acide.
\end{exoresolu}

\begin{exemplebox}[Et si l'on inverse les rôles ?]
On peut tout aussi bien doser une base forte par un acide fort : la soude (volume $V_B$, concentration $C_B$ inconnue) est placée dans le bécher et l'acide chlorhydrique de concentration $C_A$ connue dans la burette. À l'équivalence, on a toujours $n(\ce{OH-})_{\text{initial}} = n(\ce{H3O+})_{\text{versé}}$, soit $C_B \cdot V_B = C_A \cdot V_{AE}$. Par exemple, si \SI{10,0}{mL} de soude sont neutralisés par \SI{12,5}{mL} d'acide à \SI{0,080}{mol.L^{-1}}, alors $C_B = \dfrac{0{,}080 \times 12{,}5}{10{,}0} = \SI{0,10}{mol.L^{-1}}$.
\end{exemplebox}

\begin{attention}[Précautions expérimentales et pièges de calcul]
\begin{itemize}
\item Le prélèvement $V_A$ se fait à la \textbf{pipette jaugée} (verrerie de précision) ; le volume $V_{BE}$ se lit sur la \textbf{burette graduée} au dixième de millilitre près.
\item Homogénéiser le mélange (agitation magnétique) avant chaque lecture de pH.
\item Ne pas oublier les \textbf{facteurs de dilution} si la solution commerciale a été diluée avant le dosage : la concentration calculée est celle de la solution \emph{diluée}, il faut remonter à la solution mère.
\item Respecter les chiffres significatifs : le nombre de chiffres significatifs du résultat est limité par la donnée la moins précise (ici $C_B = \SI{0,10}{mol.L^{-1}}$ impose deux chiffres significatifs).
\end{itemize}
\end{attention}

\begin{exercice}[Vérification d'un détartrant (acide chlorhydrique)]
Pour vérifier l'étiquette d'un détartrant, on prélève $V_A = \SI{10,0}{mL}$ d'une solution d'acide chlorhydrique et on la dose par de la soude à $C_B = \SI{0,050}{mol.L^{-1}}$. L'équivalence est obtenue pour $V_{BE} = \SI{18,4}{mL}$, avec $\mathrm{pH}_E = 7$.\par
1. Écrire l'équation-bilan de la réaction de dosage et citer ses trois caractéristiques qui en font une réaction de dosage.\par
2. Définir l'équivalence et écrire la relation entre les quantités de matière.\par
3. Calculer la concentration $C_A$ de la solution d'acide.
\end{exercice}

\begin{corrige}[Exercice 1]
1. $\ce{H3O+ + OH- -> 2 H2O}$ ; la réaction est totale ($K = 10^{14}$), rapide et unique : elle convient donc pour un dosage.\par
2. L'équivalence est l'état du mélange où les réactifs ont été introduits dans les proportions stœchiométriques : $n(\ce{H3O+})_{\text{initial}} = n(\ce{OH-})_{\text{versé}}$, soit $C_A \cdot V_A = C_B \cdot V_{BE}$.\par
3. $C_A = \dfrac{C_B \cdot V_{BE}}{V_A} = \dfrac{0{,}050 \times 18{,}4}{10{,}0} = \SI{9,2e-2}{mol.L^{-1}}$ (deux chiffres significatifs).
\end{corrige}

\begin{aretenir}[L'essentiel de la section]
\begin{itemize}
\item Acide fort et base forte sont totalement dissociés dans l'eau : la réaction de dosage se réduit à \ce{H3O+ + OH- -> 2 H2O}, les ions \ce{Na+} et \ce{Cl-} étant spectateurs.
\item Cette réaction est \textbf{totale} ($K = 10^{14}$), \textbf{rapide} et \textbf{exothermique} : parfaite pour un dosage.
\item À l'\textbf{équivalence}, les réactifs sont en proportions stœchiométriques : $C_A \cdot V_A = C_B \cdot V_{BE}$.
\item Pour un dosage acide fort -- base forte, $\mathrm{pH}_E = 7$ (solution équivalente à du \ce{NaCl} en solution, neutre).
\item La mesure de $V_{BE}$ permet de calculer la concentration inconnue : c'est le principe de la vérification des étiquettes.
\end{itemize}
\end{aretenir}

\coursec{Dosage pH-métrique d'un acide fort par une base forte}

Dans la section précédente, nous avons établi que la réaction entre un acide fort et une base forte, \ce{H3O+ + OH- -> 2 H2O}, est \cle{totale}, \cle{rapide} et \cle{exothermique}, et nous avons défini l'\cle{équivalence acido-basique} par la relation $C_A \cdot V_A = C_B \cdot V_{BE}$. Une question pratique demeure : comment, au laboratoire, repérer avec précision cet instant où les réactifs ont été mélangés dans les proportions stœchiométriques ? C'est toute la question du \cle{dosage}. La méthode la plus rigoureuse est le \cle{dosage pH-métrique} : on suit, à l'aide d'un pH-mètre, l'évolution du pH du mélange au fur et à mesure que l'on verse la solution titrante. La courbe obtenue est une véritable « carte » de la transformation : elle révèle le point d'équivalence, permet de calculer la concentration inconnue et, comme nous le verrons dans les sections suivantes, trahit la force de l'acide dosé.

\courssub{Principe du dosage pH-métrique}

\begin{definition}[Dosage pH-métrique]
Un \cle{dosage pH-métrique} consiste à suivre l'évolution du pH d'une solution à doser (l'\cle{analyte}), placée dans un bécher, en fonction du volume $V_B$ (ou $V_A$) de solution \cle{titrante} versée goutte à goutte depuis une burette graduée. Le pH est mesuré après chaque ajout, à l'aide d'un pH-mètre dont l'électrode de verre plonge dans le mélange maintenu sous agitation magnétique constante. On trace ensuite la courbe $\mathrm{pH} = f(V_B)$.
\end{definition}

\begin{experience}[Dispositif expérimental]
\textbf{Dispositif.} Une burette graduée contenant la solution de soude (base forte) de concentration $C_B$ connue surmonte un bécher contenant un volume $V_A$ précis (prélevé à la pipette jaugée) de solution d'acide chlorhydrique de concentration $C_A$ inconnue. Une électrode de pH-mètre plonge dans le bécher ; un agitateur magnétique homogénéise le mélange en permanence.

\textbf{Protocole.}
\begin{enumerate}
\item Étalonner le pH-mètre avec deux solutions tampons (pH = 4 et pH = 7, voire pH = 10).
\item Introduire $V_A = \SI{20,0}{mL}$ d'acide chlorhydrique dans le bécher, ajouter un peu d'eau distillée pour que l'électrode soit bien immergée.
\item Verser la soude par fractions : mL par mL loin de l'équivalence, goutte à goutte (par exemple \SI{0,2}{mL}) au voisinage du saut de pH.
\item Noter le pH après chaque ajout, une fois la valeur stabilisée.
\item Tracer la courbe $\mathrm{pH} = f(V_B)$.
\end{enumerate}

\textbf{Observation.} Le pH, initialement faible, augmente d'abord très lentement, puis brutalement autour d'un certain volume, avant de se stabiliser à une valeur élevée.
\end{experience}

\begin{attention}[Deux précautions expérimentales essentielles]
\begin{itemize}
\item L'agitation doit être \textbf{constante mais modérée} : sans agitation, le pH mesuré ne reflète pas le pH réel du mélange ; une agitation trop vive fait entrer de l'air et projette la solution sur l'électrode.
\item L'eau distillée ajoutée pour immerger l'électrode \textbf{ne modifie pas le volume équivalent} : la dilution change les concentrations, mais pas les quantités de matière $n(\ce{H3O+})$ initiales. La relation $C_A V_A = C_B V_{BE}$ reste donc valable.
\end{itemize}
\end{attention}

\courssub{Allure de la courbe pH = f(VB) et analyse des quatre zones}

Considérons le dosage de $V_A = \SI{20,0}{mL}$ d'acide chlorhydrique de concentration $C_A = \SI{0,10}{mol.L^{-1}}$ par une solution de soude de concentration $C_B = \SI{0,10}{mol.L^{-1}}$. La courbe $\mathrm{pH} = f(V_B)$ a l'allure d'une marche d'escalier adoucie, en forme de « S » couché.

\begin{popfigure}
\begin{center}
\begin{tikzpicture}
\begin{axis}[
  width=0.95\linewidth, height=7.5cm,
  xlabel={$V_B$ (mL)}, ylabel={pH},
  xmin=0, xmax=40, ymin=0, ymax=14,
  xtick={0,5,10,15,20,25,30,35,40},
  ytick={0,1,2,3,4,5,6,7,8,9,10,11,12,13,14},
  grid=both, grid style={popline!40},
  axis line style={popdark},
  label style={popdark},
]
\addplot[popcurve,domain=0:39.6,samples=400]
  {7 + 6*tanh((x-20)/1.6) - 6*tanh((0-20)/1.6) + 0.02*x};
\addplot[popgold, dashed, domain=0:20] {7};
\addplot[popgold, dashed] coordinates {(20,0) (20,7)};
\addplot[only marks, mark=*, mark size=2.5pt, poporange] coordinates {(20,7)};
\node[poporange, anchor=west, font=\bfseries] at (axis cs:21,6.2) {$E\,(V_{BE}\,;\,\mathrm{pH}_E)$};
\node[popmuted, anchor=east, font=\small] at (axis cs:19,1.6) {acide en excès};
\node[popmuted, anchor=west, font=\small] at (axis cs:23,12.6) {base en excès};
\node[popmuted, anchor=south, font=\small] at (axis cs:20,7.4) {$\mathrm{pH}_E = 7$};
\end{axis}
\end{tikzpicture}
\end{center}
\legende{Courbe $\mathrm{pH} = f(V_B)$ du dosage de \SI{20,0}{mL} de \ce{HCl} à \SI{0,10}{mol.L^{-1}} par \ce{NaOH} à \SI{0,10}{mol.L^{-1}}. Le point d'équivalence $E$ est repéré à $V_{BE} = \SI{20,0}{mL}$ et $\mathrm{pH}_E = 7$.}
\end{popfigure}

Analysons qualitativement les \textbf{quatre zones} de cette courbe, en raisonnant sur les espèces présentes dans le bécher.

\textbf{Zone 1 : point de départ ($V_B = 0$).} Le bécher ne contient que l'acide fort. Le pH est faible : ici $\mathrm{pH} = -\log C_A = -\log(0{,}10) = 1{,}0$. Les ions \ce{H3O+} règnent en maîtres.

\textbf{Zone 2 : avant l'équivalence ($0 < V_B < V_{BE}$), l'acide est en excès.} Chaque ion \ce{OH-} versé est immédiatement consommé par la réaction totale \ce{H3O+ + OH- -> 2 H2O}. Il reste donc des ions \ce{H3O+} en excès : le milieu demeure acide. Le pH augmente, mais \emph{lentement} : la courbe est presque plate. Par exemple, pour $V_B = \SI{10}{mL}$ (la moitié de l'équivalence), il reste $n(\ce{H3O+}) = C_A V_A - C_B V_B = 0{,}10 \times (20 - 10)\times 10^{-3} = \SI{1,0e-3}{mol}$ dans un volume total de \SI{30}{mL}, soit $[\ce{H3O+}] \approx \SI{3,3e-2}{mol.L^{-1}}$ et $\mathrm{pH} \approx 1{,}5$. Même à \SI{19}{mL} versés, le pH ne vaut encore qu'environ $2{,}6$ : l'augmentation reste modeste.

\textbf{Zone 3 : l'équivalence ($V_B = V_{BE}$), le saut de pH.} Au voisinage immédiat de l'équivalence, la moindre goutte de soude fait basculer le milieu : juste avant, il reste un léger excès d'ions \ce{H3O+} ; juste après, il y a un léger excès d'ions \ce{OH-}. Comme le pH est une échelle \emph{logarithmique}, ce basculement se traduit par une variation brutale : le pH passe d'environ 4 à environ 10 pour quelques gouttes versées. C'est le \cle{saut de pH}. À l'équivalence exacte, le bécher ne contient qu'une solution de chlorure de sodium (\ce{Na+} et \ce{Cl-}, ions spectateurs sans action sur l'eau) : la solution est \textbf{neutre}, d'où $\mathrm{pH}_E = 7$ à \SI{25}{\celsius}.

\textbf{Zone 4 : après l'équivalence ($V_B > V_{BE}$), la base est en excès.} Les ions \ce{OH-} versés ne trouvent plus d'ions \ce{H3O+} à consommer : ils s'accumulent. Le pH, élevé, croît de plus en plus lentement et tend vers un \cle{palier} dont la valeur limite est voisine du pH de la soude titrante ($\mathrm{pH} = 14 + \log C_B = 13$ pour $C_B = \SI{0,10}{mol.L^{-1}}$).

\begin{aretenir}[Ce qu'il faut retenir de la courbe]
\begin{itemize}
\item La courbe $\mathrm{pH} = f(V_B)$ d'un dosage acide fort -- base forte est \textbf{croissante} et présente un \textbf{saut de pH} marqué.
\item Avant l'équivalence, l'acide est en excès (pH faible, croissance lente) ; après l'équivalence, la base est en excès (pH élevé, palier).
\item À l'équivalence, pour un dosage \textbf{acide fort -- base forte}, la solution est neutre : $\mathrm{pH}_E = 7$ à \SI{25}{\celsius}.
\end{itemize}
\end{aretenir}

\courssub{Le point d'équivalence E : définition et propriétés}

\begin{definition}[Point d'équivalence E]
Le \cle{point d'équivalence} $E$ est le \cle{point d'inflexion} de la courbe $\mathrm{pH} = f(V_B)$ : c'est le point où la courbe change de concavité, situé \textbf{au milieu du saut de pH}. Ses coordonnées sont $(V_{BE}\,;\,\mathrm{pH}_E)$, où $V_{BE}$ est le volume de base versé à l'équivalence. Pour le dosage d'un acide fort par une base forte, $\mathrm{pH}_E = 7$ à \SI{25}{\celsius}.
\end{definition}

Pourquoi parler de point d'inflexion ? Avant le saut, la courbe s'incurve vers le haut (sa pente augmente) ; après le saut, elle s'incurve vers le bas (sa pente diminue). Le point où la pente est \emph{maximale} est précisément le milieu du saut : c'est $E$. Cette propriété géométrique fournit deux méthodes de détermination graphique : la méthode des tangentes parallèles et la méthode de la dérivée.

\courssub{Méthode des tangentes parallèles}

\begin{methode}[Détermination de E par les tangentes parallèles]
\begin{enumerate}
\item \textbf{Tracer la première tangente} $T_1$ à la courbe dans la partie \emph{avant} le saut de pH (zone où la courbe est presque rectiligne).
\item \textbf{Tracer la deuxième tangente} $T_2$ à la courbe dans la partie \emph{après} le saut, en veillant à ce qu'elle soit \textbf{parallèle} à $T_1$.
\item \textbf{Tracer la droite $\Delta$ équidistante} de $T_1$ et $T_2$, c'est-à-dire la parallèle à ces deux tangentes située à égale distance de chacune.
\item \textbf{Lire le point E} : l'intersection de $\Delta$ avec la courbe est le point d'équivalence $E$. On projette $E$ sur l'axe des abscisses pour lire $V_{BE}$ et sur l'axe des ordonnées pour lire $\mathrm{pH}_E$.
\end{enumerate}
\end{methode}

\begin{popfigure}
\begin{center}
\begin{tikzpicture}[scale=1.05]
% axes
\draw[->, popdark, thick] (0,0) -- (9.2,0) node[below left] {$V_B$ (mL)};
\draw[->, popdark, thick] (0,0) -- (0,6.6) node[below left] {pH};
% courbe (zoom sur le saut, V_B de 15 à 25 mL)
\draw[popblue, very thick, domain=0:9, samples=200]
  plot (\x, {3.0 + 2.55*tanh((\x-4.5)/0.9)});
% tangentes paralleles (pente = 0.5)
\draw[popgreen, thick] (0.2,0.85) -- (2.6,2.05) node[pos=0.1, below, sloped, font=\small] {$T_1$};
\draw[popgreen, thick] (6.4,4.35) -- (8.8,5.55) node[pos=0.9, above, sloped, font=\small] {$T_2$};
% droite equidistante
\draw[poppurple, thick, dashed] (2.6,2.6) -- (6.6,4.65) node[pos=0.95, above, sloped, font=\small] {$\Delta$};
% point E
\fill[poporange] (4.5,3.0) circle (2.6pt) node[above right, poporange, font=\bfseries] {$E$};
% projections
\draw[popgold, dashed] (4.5,3.0) -- (4.5,0) node[below, popdark] {$V_{BE}$};
\draw[popgold, dashed] (4.5,3.0) -- (0,3.0) node[left, popdark] {$\mathrm{pH}_E = 7$};
% graduations indicatives
\foreach \x/\lab in {1.8/17, 4.5/20, 7.2/23}
  \draw[popdark] (\x,0.06) -- (\x,-0.06) node[below, font=\small] {\lab};
\foreach \y/\lab in {1.2/5, 3.0/7, 4.8/9}
  \draw[popdark] (0.06,\y) -- (-0.06,\y) node[left, font=\small] {\lab};
\end{tikzpicture}
\end{center}
\legende{Zoom sur le saut de pH : méthode des tangentes parallèles. Les tangentes $T_1$ et $T_2$ sont parallèles ; la droite $\Delta$, équidistante des deux, coupe la courbe au point d'équivalence $E$.}
\end{popfigure}

\courssub{Méthode de la dérivée (notion)}

La pente de la courbe $\mathrm{pH} = f(V_B)$ en un point est donnée par la \cle{dérivée} $\dfrac{\mathrm{d}\mathrm{pH}}{\mathrm{d}V_B}$. Or nous avons vu que cette pente est \emph{maximale} au point d'équivalence. Par conséquent, si l'on trace la courbe dérivée $\dfrac{\mathrm{d}\mathrm{pH}}{\mathrm{d}V_B} = g(V_B)$, elle présente un \textbf{maximum (pic)} dont l'abscisse est exactement $V_{BE}$.

\begin{popfigure}
\begin{center}
\begin{tikzpicture}
\begin{axis}[
  width=0.95\linewidth, height=6cm,
  xlabel={$V_B$ (mL)}, ylabel={$\dfrac{\mathrm{d}\mathrm{pH}}{\mathrm{d}V_B}$ (mL$^{-1}$)},
  xmin=0, xmax=40, ymin=0, ymax=4.5,
  xtick={0,5,10,15,20,25,30,35,40},
  grid=both, grid style={popline!40},
  axis line style={popdark},
  label style={popdark},
]
\addplot[popcurve,domain=0:40,samples=300] {0.1 + 3.6*exp(-((x-20)^2)/2.6)};
\addplot[popgold, dashed] coordinates {(20,0) (20,3.7)};
\node[poporange, anchor=south, font=\bfseries] at (axis cs:20,3.75) {pic : $V_{BE} = \SI{20,0}{mL}$};
\end{axis}
\end{tikzpicture}
\end{center}
\legende{Courbe dérivée $\dfrac{\mathrm{d}\mathrm{pH}}{\mathrm{d}V_B} = g(V_B)$ : le maximum du pic donne l'abscisse du point d'équivalence sans tracé de tangentes.}
\end{popfigure}

Cette méthode est plus précise que celle des tangentes, car elle ne dépend pas de l'habileté du traceur. En pratique, on calcule pour chaque intervalle entre deux mesures le rapport $\dfrac{\Delta \mathrm{pH}}{\Delta V_B}$ et on le reporte en fonction de $V_B$ : le tableau de valeurs fait immédiatement apparaître le maximum.

\begin{savaistu}[Les pH-mètres modernes et l'ExAO]
Dans les laboratoires équipés, le pH-mètre est relié à un ordinateur (système d'\emph{Expérimentation Assistée par Ordinateur}, ExAO) : le logiciel enregistre le pH en temps réel pendant que la burette verse la solution titrante, trace la courbe $\mathrm{pH} = f(V_B)$ \textbf{et} sa dérivée simultanément, puis repère automatiquement le maximum de la dérivée. Le tracé manuel des tangentes parallèles devient inutile ! Mais attention : au Baccalauréat, c'est bien toi qui dois savoir exploiter une courbe tracée sur papier. La méthode des tangentes reste donc un savoir-faire incontournable.
\end{savaistu}

\courssub{Exploitation de la courbe : calcul de la concentration inconnue}

Une fois $V_{BE}$ lu graphiquement, l'exploitation repose sur la relation à l'équivalence. Il est recommandé de présenter le raisonnement à l'aide d'un \cle{tableau d'avancement} pour bien visualiser les quantités de matière.

\begin{methode}[Exploitation par tableau d'avancement]
\begin{enumerate}
\item Écrire l'équation-bilan de la réaction de dosage : \ce{H3O+ + OH- -> 2 H2O}.
\item Construire le tableau d'avancement (états initial, intermédiaire, final/équivalence).
\item À l'équivalence, l'avancement $x_{\text{max}}$ est tel que l'un des réactifs est entièrement consommé. Pour un dosage, on verse la titrante jusqu'à consommer tout l'analyte : $x_{\text{max}} = n_A(\text{initial}) = n_B(\text{versé à l'équivalence})$.
\item Exprimer les quantités de matière : $n_A = C_A V_A$ et $n_B = C_B V_{BE}$.
\item En déduire $C_A = \dfrac{C_B \cdot V_{BE}}{V_A}$.
\end{enumerate}
\end{methode}

\begin{exemplebox}[Lecture graphique, tableau d'avancement et calcul de $C_A$]
On dose $V_A = \SI{10,0}{mL}$ d'une solution d'acide chlorhydrique de concentration inconnue $C_A$ par de la soude à $C_B = \SI{0,050}{mol.L^{-1}}$. La courbe $\mathrm{pH} = f(V_B)$, exploitée par la méthode des tangentes parallèles, donne un point d'équivalence de coordonnées $V_{BE} = \SI{12,5}{mL}$ et $\mathrm{pH}_E = 7{,}0$.

\textbf{Tableau d'avancement (quantités de matière en mol) :}
\[
\begin{array}{|c|c|c|c|}
\hline
 & \ce{H3O+} & \ce{OH-} & \ce{H2O} \\
\hline
\text{État initial} & C_A \times 10{,}0 \times 10^{-3} & 0 & \text{solvant} \\
\hline
\text{À l'équivalence} (x_{\text{max}}) & 0 & 0 & \text{excès} \\
\hline
\end{array}
\]
À l'équivalence : $x_{\text{max}} = n_{\text{initial}}(\ce{H3O+}) = n_{\text{versé}}(\ce{OH-}) = C_B \times V_{BE}$.
\[ C_A \times 10{,}0 \times 10^{-3} = 0{,}050 \times 12{,}5 \times 10^{-3} \]

\textbf{Calcul de la concentration :}
\[ C_A = \frac{C_B \cdot V_{BE}}{V_A} = \frac{0{,}050 \times 12{,}5}{10{,}0} = \SI{0,0625}{mol.L^{-1}} \]

\textbf{Vérification de cohérence.} Les quantités à l'équivalence : $n_B = C_B V_{BE} = 0{,}050 \times 12{,}5 \times 10^{-3} = \SI{6,25e-4}{mol}$ et $n_A = C_A V_A = 0{,}0625 \times 10{,}0 \times 10^{-3} = \SI{6,25e-4}{mol}$. Les deux quantités sont égales : l'équivalence est bien vérifiée. De plus, $\mathrm{pH}_E = 7$ confirme qu'il s'agit bien d'un dosage acide fort -- base forte.
\end{exemplebox}

\begin{exoresolu}[Dosage d'un détartrant à base d'acide chlorhydrique]
\textbf{Énoncé.} Un détartrant ménager contient de l'acide chlorhydrique. On prélève $V_A = \SI{20,0}{mL}$ d'une solution diluée 10 fois de ce détartrant et on la dose par de la soude à $C_B = \SI{0,100}{mol.L^{-1}}$. Le suivi pH-métrique donne les résultats suivants (extraits) :

\begin{center}
\begin{tabularx}{0.85\linewidth}{|l|X|X|X|X|X|X|}
\hline
\rowcolor{popblueL} $V_B$ (mL) & 0 & 5,0 & 10,0 & 14,0 & 15,0 & 16,0 \\
\hline pH & 1,0 & 1,2 & 1,5 & 2,2 & 7,0 & 11,2 \\
\hline
\end{tabularx}
\end{center}

\begin{enumerate}
\item Écrire l'équation-bilan de la réaction de dosage et préciser ses caractéristiques.
\item Déterminer $V_{BE}$ et justifier.
\item Calculer la concentration $C_A$ de la solution diluée, puis la concentration $C_0$ du détartrant commercial.
\end{enumerate}

\tcblower
\textbf{Solution détaillée.}

\textbf{1.} L'acide chlorhydrique (\ce{H3O+} + \ce{Cl-}) est un acide fort ; la soude (\ce{Na+} + \ce{OH-}) est une base forte. La réaction de dosage est :
\[ \ce{H3O+ + OH- -> 2 H2O} \]
Elle est \textbf{totale} ($K = 10^{14} \gg 10^4$), \textbf{rapide} et \textbf{exothermique} : elle convient parfaitement à un dosage.

\textbf{2.} Le saut de pH se produit entre $V_B = \SI{14,0}{mL}$ (pH = 2,2) et $V_B = \SI{16,0}{mL}$ (pH = 11,2). Le point où $\mathrm{pH} = 7$ correspond à $V_{BE} = \SI{15,0}{mL}$. C'est le milieu du saut : le point d'équivalence est $E\,(\SI{15,0}{mL}\,;\,7{,}0)$.

\textbf{3.} À l'équivalence : $C_A \cdot V_A = C_B \cdot V_{BE}$, d'où
\[ C_A = \frac{C_B \cdot V_{BE}}{V_A} = \frac{0{,}100 \times 15{,}0}{20{,}0} = \SI{0,075}{mol.L^{-1}} \]
Le détartrant a été dilué 10 fois, donc :
\[ C_0 = 10 \times C_A = \SI{0,75}{mol.L^{-1}} \]
\end{exoresolu}

\courssub{Dosage réciproque : base forte par acide fort}

Que se passe-t-il si l'on inverse les rôles, c'est-à-dire si l'on place la soude dans le bécher et l'acide chlorhydrique dans la burette ? La réaction de dosage est \textbf{la même} : \ce{H3O+ + OH- -> 2 H2O}, toujours totale, rapide et exothermique. Seul le sens de variation du pH change.

Au départ, le bécher contient la base forte : le pH initial est \emph{élevé} (par exemple pH = 13 pour une soude à \SI{0,10}{mol.L^{-1}}). Au fur et à mesure de l'ajout d'acide, les ions \ce{OH-} sont consommés : le pH \textbf{décroît}, d'abord lentement (base en excès), puis brutalement au voisinage de l'équivalence (saut de pH \emph{descendant}), avant de se stabiliser sur un palier acide (acide en excès). À l'équivalence, le bécher ne contient toujours que du chlorure de sodium en solution : $\mathrm{pH}_E = 7$, et la relation d'équivalence s'écrit de manière symétrique :

\[ C_B \cdot V_B = C_A \cdot V_{AE} \]

\begin{popfigure}
\begin{center}
\begin{tikzpicture}
\begin{axis}[
  width=0.95\linewidth, height=7cm,
  xlabel={$V_A$ (mL)}, ylabel={pH},
  xmin=0, xmax=40, ymin=0, ymax=14,
  xtick={0,5,10,15,20,25,30,35,40},
  ytick={0,2,4,6,7,8,10,12,14},
  grid=both, grid style={popline!40},
  axis line style={popdark},
  label style={popdark},
]
\addplot[popcurve,domain=0:39.6,samples=400]
  {7 - 6*tanh((x-20)/1.6) + 6*tanh((0-20)/1.6) - 0.02*x};
\addplot[popgold, dashed, domain=0:20] {7};
\addplot[popgold, dashed] coordinates {(20,0) (20,7)};
\addplot[only marks, mark=*, mark size=2.5pt, poporange] coordinates {(20,7)};
\node[poporange, anchor=west, font=\bfseries] at (axis cs:21,7.8) {$E\,(V_{AE} = \SI{20,0}{mL}\,;\,7)$};
\node[popmuted, anchor=east, font=\small] at (axis cs:18,12.8) {base en excès};
\node[popmuted, anchor=west, font=\small] at (axis cs:23,1.4) {acide en excès};
\end{axis}
\end{tikzpicture}
\end{center}
\legende{Dosage réciproque : \SI{20,0}{mL} de soude à \SI{0,10}{mol.L^{-1}} dosés par \ce{HCl} à \SI{0,10}{mol.L^{-1}}. La courbe est décroissante, mais l'équivalence est identique : $V_{AE} = \SI{20,0}{mL}$, $\mathrm{pH}_E = 7$.}
\end{popfigure}

\begin{aretenir}[Dosage réciproque : ce qui change, ce qui ne change pas]
\begin{itemize}
\item \textbf{Change} : le sens de variation (courbe décroissante), le pH initial (élevé), le réactif en excès dans chaque zone.
\item \textbf{Ne change pas} : la réaction de dosage \ce{H3O+ + OH- -> 2 H2O}, l'existence d'un saut de pH, le point d'équivalence comme point d'inflexion, et $\mathrm{pH}_E = 7$ à \SI{25}{\celsius}.
\end{itemize}
\end{aretenir}

\begin{attention}[Erreur fréquente : « l'équivalence, c'est toujours pH = 7 »]
C'est \textbf{faux} en général ! La valeur $\mathrm{pH}_E = 7$ n'est vraie \emph{que} pour le dosage d'un acide fort par une base forte (ou réciproquement), car le sel formé (\ce{NaCl}) est sans action sur l'eau. Nous verrons dans les sections suivantes que le dosage d'un acide \emph{faible} par une base forte donne $\mathrm{pH}_E > 7$, et celui d'une base \emph{faible} par un acide fort donne $\mathrm{pH}_E < 7$. La bonne définition, valable dans tous les cas, est : \textbf{l'équivalence est le point d'inflexion de la courbe, au milieu du saut de pH}, repéré par la méthode des tangentes parallèles ou par le maximum de la dérivée. Chercher systématiquement le point où la courbe coupe la droite $\mathrm{pH} = 7$ est un piège classique de l'examen.
\end{attention}

\begin{exercice}[À toi de jouer]
On dose $V_A = \SI{25,0}{mL}$ d'une solution d'acide chlorhydrique par de la soude à $C_B = \SI{0,080}{mol.L^{-1}}$. La méthode des tangentes parallèles, appliquée à la courbe $\mathrm{pH} = f(V_B)$, donne $V_{BE} = \SI{16,0}{mL}$ et $\mathrm{pH}_E = 7{,}0$.
\begin{enumerate}
\item Calculer la concentration $C_A$ de la solution d'acide.
\item Quel est le pH initial de la solution dosée ?
\item Si l'on avait dosé la soude par l'acide (dosage réciproque) avec les mêmes concentrations, quelle serait l'allure de la courbe et la valeur du pH à l'équivalence ?
\end{enumerate}
\end{exercice}

\begin{corrige}[Exercice n°1]
\textbf{1.} À l'équivalence : $C_A V_A = C_B V_{BE}$, donc
\[ C_A = \frac{0{,}080 \times 16{,}0}{25{,}0} = \SI{0,0512}{mol.L^{-1}} \approx \SI{5,1e-2}{mol.L^{-1}}. \]

\textbf{2.} L'acide chlorhydrique est un acide fort totalement dissocié : $\mathrm{pH} = -\log C_A = -\log(0{,}0512) \approx 1{,}3$.

\textbf{3.} La courbe serait \textbf{décroissante} : départ à pH élevé (soude en excès), saut de pH descendant, puis palier acide. La réaction étant toujours celle d'un acide fort sur une base forte, le pH à l'équivalence resterait $\mathrm{pH}_E = 7$ à \SI{25}{\celsius}.
\end{corrige}

En résumé, le dosage pH-métrique d'un acide fort par une base forte fournit une courbe en « S » dont le point d'inflexion — le point d'équivalence $E$, de pH égal à 7 — se détermine par les tangentes parallèles ou par le maximum de la dérivée, et dont l'abscisse $V_{BE}$ permet de remonter à la concentration inconnue grâce à $C_A V_A = C_B V_{BE}$. Mais cette belle symétrie se brise dès que l'un des réactifs est faible : c'est ce que nous découvrirons dans la section suivante, en dosant l'acide éthanoïque du vinaigre par la soude.

\coursec{Dosage d'un acide faible par une base forte}

Dans la section précédente, nous avons dosé un acide fort (l'acide chlorhydrique) par une base forte (la soude) : la courbe $pH = f(V_B)$ partait d'un pH très acide, présentait un grand saut vertical et franchissait l'équivalence à $pH_E = 7$. Mais que se passe-t-il lorsque l'acide dosé est \emph{faible}, comme l'acide éthanoïque du vinaigre ? C'est précisément la situation que rencontre le technicien de laboratoire qui veut vérifier le degré d'acidité d'un vinaigre de table à Douala ou à Bafoussam. Nous allons voir que la courbe change d'allure, que l'équivalence ne se situe plus à pH 7, et surtout qu'un nouveau point remarquable apparaît : le \cle{point de demi-équivalence}, véritable trésor puisqu'il permet de mesurer expérimentalement le $pK_a$ d'un couple acide/base.

\courssub{La réaction de dosage : acide éthanoïque et soude}

\textbf{Équation-bilan et caractéristiques}\par

Plaçons-nous dans la situation de référence : on dose un volume $V_A$ d'une solution d'\cle{acide éthanoïque} \ce{CH3COOH} (acide faible, $pK_a = 4{,}8$) de concentration $C_A$, par une solution de \cle{soude} \ce{NaOH} (base forte) de concentration $C_B$ versée depuis la burette.

L'acide éthanoïque étant faible, il est \emph{très peu dissocié} dans l'eau : c'est bien la forme \ce{CH3COOH} qui prédomine dans le bécher au départ. Les ions hydroxyde \ce{OH-} apportés par la soude attaquent directement les molécules d'acide :

\[ \ce{CH3COOH + OH- -> CH3COO- + H2O} \]

\begin{propriete}[Caractéristiques de la réaction de dosage]
La réaction entre l'acide éthanoïque et les ions hydroxyde est :
\begin{itemize}
\item \textbf{totale} : sa constante d'équilibre vaut
\[ K = \frac{K_a}{K_e} = 10^{14 - pK_a} = 10^{14 - 4{,}8} = 10^{9{,}2} \gg 10^4 \]
\item \textbf{rapide} : c'est une réaction ionique, quasi instantanée ;
\item \textbf{exothermique} : elle dégage de la chaleur.
\end{itemize}
Ces trois qualités (totale, rapide, unique) en font une excellente \cle{réaction de dosage}.
\end{propriete}

\begin{attention}[Ne pas confondre les deux réactions]
L'erreur classique consiste à écrire l'équation de \emph{dissociation} de l'acide dans l'eau, \ce{CH3COOH + H2O <=> CH3COO- + H3O+} (réaction \emph{limitée}, d'où la double flèche), à la place de l'équation de \emph{dosage} \ce{CH3COOH + OH- -> CH3COO- + H2O} (réaction \emph{totale}, flèche simple). Au baccalauréat, la réaction de dosage s'écrit toujours avec \ce{OH-} comme réactif et une flèche simple.
\end{attention}

\courssub{Allure de la courbe $pH = f(V_B)$}

Suivons pas à pas ce qui se passe dans le bécher lorsqu'on verse progressivement la soude, en raisonnant sur les espèces présentes.

\begin{itemize}
\item \textbf{Au départ ($V_B = 0$)} : la solution contient un acide faible. Le pH initial est donc \emph{plus élevé} que celui d'un acide fort de même concentration : pour $C_A = \SI{0,1}{mol.L^{-1}}$, on lit $pH_i \approx 2{,}9$ (au lieu de 1,0 pour HCl).
\item \textbf{Avant l'équivalence} : les \ce{OH-} versés transforment \ce{CH3COOH} en \ce{CH3COO-}. Le mélange contient alors le couple \ce{CH3COOH}/\ce{CH3COO-} en quantités comparables : c'est une \cle{solution tampon}, le pH monte \emph{lentement}.
\item \textbf{À l'équivalence ($V_B = V_{BE}$)} : tout l'acide a été transformé en ion éthanoate \ce{CH3COO-}, qui est une \emph{base faible}. La solution est donc \textbf{basique} : $pH_E > 7$.
\item \textbf{Après l'équivalence} : les \ce{OH-} s'accumulent ; le pH rejoint progressivement celui de la soude versée.
\end{itemize}

\begin{popfigure}
\begin{center}
\begin{tikzpicture}
\begin{axis}[
  width=0.92\linewidth, height=7.2cm,
  xlabel={$V_B$ (mL)}, ylabel={pH},
  xmin=0, xmax=40, ymin=0, ymax=14,
  xtick={0,8,16,24,32,40}, ytick={0,2,4,6,8,10,12,14},
  grid=both, grid style={popline!40},
  legend style={at={(0.03,0.97)}, anchor=north west, font=\small}
]
% Courbe acide faible - base forte (approximation lissée)
\addplot[popcurve, domain=0:15.2, samples=150] {2.87 + 0.9*ln(1+x/(16-x))/ln(10) + 1.05*ln(1+x)/ln(10)};
\addplot[popcurve, domain=15.2:16.8, samples=100] {8.7 + 3.2*tanh(8*(x-16))};
\addplot[popcurve, domain=16.8:40, samples=100] {11.7 + 1.0*ln(1+(x-16.8)/3)/ln(10) + 0.35*ln(1+(x-16.8))/ln(10)};
% Point E
\addplot[only marks, mark=*, mark size=2.5pt, poporange] coordinates {(16,8.7)};
\node[poporange, font=\small\bfseries, anchor=south west] at (axis cs:16.4,8.7) {$E\,(16\,;\,8{,}7)$};
% Point H
\addplot[only marks, mark=*, mark size=2.5pt, poppurple] coordinates {(8,4.8)};
\node[poppurple, font=\small\bfseries, anchor=south east] at (axis cs:7.8,4.9) {$H\,(8\,;\,4{,}8)$};
% Pointilles
\addplot[dashed, popmuted, domain=0:16] coordinates {(16,0) (16,8.7) (0,8.7)};
\addplot[dashed, popmuted] coordinates {(8,0) (8,4.8) (0,4.8)};
\node[popmuted, font=\footnotesize, anchor=north] at (axis cs:16,0) {$V_{BE}$};
\node[popmuted, font=\footnotesize, anchor=north] at (axis cs:8,0) {$\dfrac{V_{BE}}{2}$};
\node[popmuted, font=\footnotesize, anchor=west] at (axis cs:0.3,8.7) {$pH_E$};
\node[popmuted, font=\footnotesize, anchor=west] at (axis cs:0.3,4.8) {$pK_a$};
\end{axis}
\end{tikzpicture}
\end{center}
\legende{Courbe de dosage de \SI{20,0}{mL} d'acide éthanoïque à \SI{0,10}{mol.L^{-1}} par la soude à \SI{0,10}{mol.L^{-1}}. Le point $E$ est le point d'équivalence ($pH_E \approx 8{,}7 > 7$) ; le point $H$ est le point de demi-équivalence ($pH = pK_a = 4{,}8$).}
\end{popfigure}

Comparée à la courbe acide fort–base forte, cette courbe présente trois différences essentielles à retenir :

\begin{aretenir}[Signatures d'un dosage acide faible–base forte]
\begin{itemize}
\item le \textbf{pH initial est plus élevé} (acide faible peu dissocié) ;
\item le \textbf{saut de pH est moins important} et commence vers pH 7 ;
\item le \textbf{pH à l'équivalence est supérieur à 7} ($pH_E \approx 8{,}7$ ici).
\end{itemize}
\end{aretenir}

\courssub{Pourquoi le pH à l'équivalence est-il supérieur à 7 ?}

C'est LA question que l'examinateur adore poser. La réponse tient en une phrase : \emph{regarde ce qui reste dans le bécher à l'équivalence}.

À l'équivalence, tout l'acide \ce{CH3COOH} a été consommé et transformé en ion éthanoate \ce{CH3COO-}. La solution contient donc de l'\cle{éthanoate de sodium} \ce{CH3COO- + Na+} : les ions \ce{Na+} sont spectateurs (indifférents vis-à-vis de l'eau), mais l'ion \ce{CH3COO-} est la \emph{base conjuguée} de l'acide éthanoïque. Il réagit (faiblement) avec l'eau :

\[ \ce{CH3COO- + H2O <=> CH3COOH + OH-} \]

Cette réaction limitée produit des ions hydroxyde \ce{OH-} : la solution est donc \textbf{basique} à l'équivalence, d'où $pH_E > 7$ (environ 8,7 pour nos concentrations).

\begin{attention}[Erreur fréquente n°1]
Écrire « le pH à l'équivalence vaut toujours 7 » est \textbf{faux}. La valeur $pH_E = 7$ n'est valable que pour le dosage d'un \emph{acide fort par une base forte} (les ions formés, \ce{Na+} et \ce{Cl-}, sont indifférents). Le pH à l'équivalence dépend de la \textbf{nature de l'espèce formée} : base faible formée $\Rightarrow pH_E > 7$ ; acide faible formé $\Rightarrow pH_E < 7$ ; espèces indifférentes $\Rightarrow pH_E = 7$.
\end{attention}

\courssub{Le point de demi-équivalence : une mine d'or}

\textbf{Définition et démonstration}\par

Intuition d'abord : à l'équivalence, on a versé exactement de quoi transformer \emph{tout} l'acide. À \emph{mi-chemin} de l'équivalence, on n'a donc transformé que la \emph{moitié} de l'acide : il reste autant de \ce{CH3COOH} qu'il s'est formé de \ce{CH3COO-}. Or, la relation de Henderson-Hasselbalch nous dit que :

\[ pH = pK_a + \log\frac{[\ce{CH3COO-}]}{[\ce{CH3COOH}]} \]

Si $[\ce{CH3COOH}] = [\ce{CH3COO-}]$, alors le logarithme vaut $\log 1 = 0$ et il reste simplement $pH = pK_a$.

\begin{definition}[Point de demi-équivalence]
Le \cle{point de demi-équivalence} $H$ est le point de la courbe de dosage où la \textbf{moitié de l'acide a réagi}. Il correspond au volume versé $V_B = \dfrac{V_{BE}}{2}$. En ce point :
\[ [\ce{CH3COOH}] = [\ce{CH3COO-}] \qquad \text{donc} \qquad \boxed{pH = pK_a} \]
Ses coordonnées sont $\left(\dfrac{V_{BE}}{2}\,;\,pK_a\right)$.
\end{definition}

\textbf{Exploitation : mesurer le $pK_a$ d'un couple}\par

Ce résultat est formidable : une simple courbe de dosage pH-métrique permet de \emph{mesurer expérimentalement} le $pK_a$ d'un couple acide/base, sans aucun calcul de concentration.

\begin{methode}[Déterminer le $pK_a$ d'un couple sur une courbe de dosage]
\begin{enumerate}
\item Repérer le point d'équivalence $E$ sur la courbe (méthode des tangentes parallèles ou de la dérivée) et lire $V_{BE}$.
\item Calculer $V_{BE}/2$ et reporter cette valeur sur l'axe des volumes.
\item Monter verticalement jusqu'à la courbe, puis lire le pH correspondant sur l'axe des ordonnées.
\item Ce pH est le $pK_a$ du couple : $pK_a = pH\left(V_{BE}/2\right)$.
\end{enumerate}
\end{methode}

\begin{exemplebox}[Le vinaigre nous livre son $pK_a$]
On dose $V_A = \SI{20,0}{mL}$ de vinaigre dilué par de la soude. La méthode des tangentes donne $V_{BE} = \SI{16,0}{mL}$.
\begin{itemize}
\item Demi-équivalence : $V_B = \dfrac{16{,}0}{2} = \SI{8,0}{mL}$.
\item Lecture graphique à $V_B = \SI{8,0}{mL}$ : $pH = 4{,}8$.
\item Conclusion : $pK_a(\ce{CH3COOH}/\ce{CH3COO-}) = 4{,}8$ — on retrouve exactement la valeur tabulée de l'acide éthanoïque !
\end{itemize}
\end{exemplebox}

\textbf{La zone tampon autour de la demi-équivalence}\par

Observe la courbe autour du point $H$ : elle y est presque \emph{plate}. Le pH varie très peu lorsqu'on ajoute de la soude dans cette région. C'est normal : le mélange contient alors des quantités comparables d'acide \ce{CH3COOH} et de base conjuguée \ce{CH3COO-} — c'est une \cle{solution tampon}. Une solution tampon \emph{s'oppose} aux variations de pH lors de l'ajout modéré d'acide ou de base, ou lors d'une dilution. Le pouvoir tampon est \emph{maximal} précisément à la demi-équivalence, où $pH = pK_a$.

\begin{popfigure}
\begin{center}
\begin{tikzpicture}[font=\small]
% Axe horizontal = avancement du dosage
\draw[-{Stealth}, thick, popdark] (0,0) -- (11.5,0) node[below left=2pt and -30pt] {$V_B$ croissant};
% Zones
\fill[popblueL, opacity=0.5] (0,0.15) rectangle (3.2,2.6);
\fill[poppurpleL, opacity=0.5] (3.2,0.15) rectangle (6.8,2.6);
\fill[popgreenL, opacity=0.5] (6.8,0.15) rectangle (8.6,2.6);
\fill[poporangeL, opacity=0.5] (8.6,0.15) rectangle (11.2,2.6);
% Marques
\draw[thick, popdark] (6.8,0) -- (6.8,-0.15) node[below] {$\dfrac{V_{BE}}{2}$};
\draw[thick, popdark] (8.6,0) -- (8.6,-0.15) node[below] {$V_{BE}$};
% Textes des zones
\node[align=center, popblue!70!black] at (1.6,1.4) {\textbf{\ce{CH3COOH}}\\ prédomine\\ $pH < pK_a$};
\node[align=center, poppurple!70!black] at (5.0,1.4) {\textbf{zone tampon}\\ $[\ce{CH3COOH}] \approx [\ce{CH3COO-}]$\\ le pH varie peu};
\node[align=center, popgreen!60!black] at (7.7,1.4) {\textbf{\ce{CH3COO-}}\\ prédomine\\ $pH > pK_a$};
\node[align=center, poporange!80!black] at (9.9,1.4) {\textbf{\ce{OH-}}\\ en excès\\ $pH \gg 7$};
% Flèche pKa
\draw[-{Stealth}, poppurple, thick] (5.0,2.75) -- (6.8,2.75);
\node[poppurple, above, font=\footnotesize] at (6.0,2.75) {$pH = pK_a$ en $H$};
\end{tikzpicture}
\end{center}
\legende{Évolution des espèces prédominantes le long du dosage de \ce{CH3COOH} par \ce{NaOH} : l'acide domine au début, le couple coexiste dans la zone tampon autour de la demi-équivalence, la base conjuguée \ce{CH3COO-} domine ensuite, puis les ions \ce{OH-} s'accumulent après l'équivalence.}
\end{popfigure}

\courssub{Comparaison avec le dosage acide fort–base forte}

Superposons les deux courbes à \emph{même concentration} (\SI{0,10}{mol.L^{-1}}) pour bien fixer les différences.

\begin{popfigure}
\begin{center}
\begin{tikzpicture}
\begin{axis}[
  width=0.92\linewidth, height=7.5cm,
  xlabel={$V_B$ (mL)}, ylabel={pH},
  xmin=0, xmax=40, ymin=0, ymax=14,
  xtick={0,8,16,24,32,40}, ytick={0,2,4,6,8,10,12,14},
  grid=both, grid style={popline!40},
  legend style={at={(0.03,0.97)}, anchor=north west, font=\small, draw=popline}
]
% Acide fort - base forte
\addplot[popcurve, popgreen, domain=0:15.5, samples=120] {1 + 0.55*ln(1+x)/ln(10) + 0.9*ln(1+x/(16-x))/ln(10)};
\addplot[popcurve, popgreen, domain=15.5:16.5, samples=80] {7 + 4.0*tanh(9*(x-16))};
\addplot[popcurve, popgreen, domain=16.5:40, samples=100] {11.0 + 1.2*ln(1+(x-16.5)/2.5)/ln(10) + 0.4*ln(1+(x-16.5))/ln(10)};
\addlegendentry{\ce{HCl} (acide fort) par \ce{NaOH}}
% Acide faible - base forte
\addplot[popcurve, popblue, domain=0:15.2, samples=150] {2.87 + 0.9*ln(1+x/(16-x))/ln(10) + 1.05*ln(1+x)/ln(10)};
\addplot[popcurve, popblue, domain=15.2:16.8, samples=100] {8.7 + 3.2*tanh(8*(x-16))};
\addplot[popcurve, popblue, domain=16.8:40, samples=100] {11.7 + 1.0*ln(1+(x-16.8)/3)/ln(10) + 0.35*ln(1+(x-16.8))/ln(10)};
\addlegendentry{\ce{CH3COOH} (acide faible) par \ce{NaOH}}
% Points remarquables
\addplot[only marks, mark=*, mark size=2.2pt, popgreen!60!black] coordinates {(16,7)};
\node[popgreen!60!black, font=\footnotesize\bfseries, anchor=west] at (axis cs:16.5,6.6) {$E_1\,(pH_E = 7)$};
\addplot[only marks, mark=*, mark size=2.2pt, popblue!70!black] coordinates {(16,8.7)};
\node[popblue!70!black, font=\footnotesize\bfseries, anchor=west] at (axis cs:16.5,9.2) {$E_2\,(pH_E \approx 8{,}7)$};
\addplot[only marks, mark=*, mark size=2.2pt, poppurple] coordinates {(8,4.8)};
\node[poppurple, font=\footnotesize\bfseries, anchor=south east] at (axis cs:7.8,5.0) {$H\,(pH = pK_a)$};
\end{axis}
\end{tikzpicture}
\end{center}
\legende{Comparaison des courbes de dosage à même concentration (\SI{0,10}{mol.L^{-1}}) : l'acide faible part d'un pH plus élevé, son saut de pH est plus court et son équivalence est basique. En revanche, le volume à l'équivalence $V_{BE} = \SI{16}{mL}$ est \emph{identique} dans les deux cas.}
\end{popfigure}

\begin{aretenir}[Ce qui change, ce qui ne change pas]
\begin{center}
\begin{tabularx}{\linewidth}{|X|X|X|}
\hline
\rowcolor{popblueL} \textbf{Grandeur} & \textbf{Acide fort / base forte} & \textbf{Acide faible / base forte} \\
\hline
pH initial ($C = \SI{0,1}{mol.L^{-1}}$) & 1,0 & 2,9 \\
\hline
Amplitude du saut de pH & grande ($\approx 4 \to 10$) & plus faible ($\approx 7 \to 11$) \\
\hline
$pH_E$ à l'équivalence & $= 7$ & $> 7$ (ici $\approx 8{,}7$) \\
\hline
Demi-équivalence & sans intérêt particulier & $pH = pK_a$ (point $H$) \\
\hline
Relation à l'équivalence & \multicolumn{2}{c|}{$C_A \cdot V_A = C_B \cdot V_{BE}$ \quad (identique !)} \\
\hline
\end{tabularx}
\end{center}
\end{aretenir}

\begin{attention}[La relation à l'équivalence ne change jamais]
Beaucoup d'élèves croient que la faiblesse de l'acide modifie la relation d'équivalence. C'est faux : la réaction de dosage étant \emph{totale}, mole pour mole, on a toujours $n(\ce{OH-})_{\text{versé}} = n(\ce{CH3COOH})_{\text{initial}}$, c'est-à-dire :
\[ C_A \cdot V_A = C_B \cdot V_{BE} \]
Seule la \emph{forme de la courbe} (pH initial, saut, $pH_E$) dépend de la force de l'acide — pas la stœchiométrie.
\end{attention}

\courssub{Exercice résolu : exploitation complète d'une courbe}

\begin{exoresolu}[Dosage de l'acide éthanoïque par la soude]
On dose $V_A = \SI{20,0}{mL}$ d'une solution d'acide éthanoïque de concentration $C_A$ inconnue par de la soude à $C_B = \SI{0,100}{mol.L^{-1}}$. La courbe pH-métrique donne : équivalence à $V_{BE} = \SI{16,0}{mL}$ avec $pH_E = 8{,}7$ ; à $V_B = \SI{8,0}{mL}$, on lit $pH = 4{,}8$.
\begin{enumerate}
\item Écrire l'équation-bilan de la réaction de dosage.
\item Déterminer $C_A$.
\item Justifier la valeur $pH_E > 7$.
\item Déterminer le $pK_a$ du couple et le $K_a$ correspondant.
\end{enumerate}
\tcblower
\textbf{1. Équation-bilan.} L'acide faible réagit totalement avec les ions hydroxyde :
\[ \ce{CH3COOH + OH- -> CH3COO- + H2O} \qquad (K = 10^{9{,}2} \gg 10^4, \text{ réaction totale et rapide}) \]

\textbf{2. Concentration $C_A$.} À l'équivalence : $C_A \cdot V_A = C_B \cdot V_{BE}$, donc
\[ C_A = \frac{C_B \cdot V_{BE}}{V_A} = \frac{0{,}100 \times 16{,}0}{20{,}0} = \SI{0,080}{mol.L^{-1}} \]

\textbf{3. Justification de $pH_E > 7$.} À l'équivalence, la solution ne contient plus d'acide \ce{CH3COOH} : elle contient l'ion éthanoate \ce{CH3COO-}, base faible, qui capte des protons de l'eau selon \ce{CH3COO- + H2O <=> CH3COOH + OH-}. La production d'ions \ce{OH-} rend la solution basique : $pH_E = 8{,}7 > 7$.

\textbf{4. Détermination du $pK_a$.} Le volume $V_B = \SI{8,0}{mL}$ est exactement $\dfrac{V_{BE}}{2} = \dfrac{16{,}0}{2}$ : c'est le point de demi-équivalence $H$, où $[\ce{CH3COOH}] = [\ce{CH3COO-}]$. Donc :
\[ pK_a = pH(H) = 4{,}8 \qquad \text{et} \qquad K_a = 10^{-pK_a} = 10^{-4{,}8} \approx 1{,}6 \times 10^{-5} \]
On retrouve les valeurs tabulées du couple \ce{CH3COOH}/\ce{CH3COO-}.
\end{exoresolu}

\begin{savaistu}[Le degré d'acidité du vinaigre]
C'est grâce à ce dosage que l'on détermine le \cle{degré d'acidité} du vinaigre vendu sur les marchés de Yaoundé ou de Garoua : \textbf{1 degré (1°) correspond à \SI{1}{g} d'acide éthanoïque pour \SI{100}{g} de vinaigre}. Un vinaigre à 6° contient donc environ \SI{60}{g} d'acide éthanoïque par litre, soit une concentration voisine de \SI{1}{mol.L^{-1}} — qu'il faut diluer dix fois avant de le doser au laboratoire ! La réglementation impose d'ailleurs un minimum d'acidité pour qu'un produit puisse porter le nom de « vinaigre ». Le même principe de dosage s'applique au contrôle de l'acidité du vin de palme, qui s'« aigrit » naturellement lorsque l'acide éthanoïque s'y accumule.
\end{savaistu}

\begin{exercice}[À toi de jouer]
On dose $V_A = \SI{10,0}{mL}$ d'une solution d'acide méthanoïque \ce{HCOOH} par de la soude à $C_B = \SI{0,050}{mol.L^{-1}}$. L'équivalence est atteinte pour $V_{BE} = \SI{12,0}{mL}$.
\begin{enumerate}
\item Écrire l'équation-bilan du dosage.
\item Calculer la concentration $C_A$ de l'acide.
\item Le pH à l'équivalence est-il égal, supérieur ou inférieur à 7 ? Justifier.
\item À quel volume versé lit-on directement le $pK_a$ du couple \ce{HCOOH}/\ce{HCOO-} sur la courbe ?
\end{enumerate}
\end{exercice}

\begin{corrige}[Exercice]
\textbf{1.} \ce{HCOOH + OH- -> HCOO- + H2O} (réaction totale, rapide, exothermique).

\textbf{2.} À l'équivalence : $C_A = \dfrac{C_B \cdot V_{BE}}{V_A} = \dfrac{0{,}050 \times 12{,}0}{10{,}0} = \SI{0,060}{mol.L^{-1}}$.

\textbf{3.} $pH_E > 7$ : à l'équivalence, la solution contient l'ion méthanoate \ce{HCOO-}, base faible qui réagit avec l'eau selon \ce{HCOO- + H2O <=> HCOOH + OH-}, d'où une solution basique.

\textbf{4.} Au point de demi-équivalence : $V_B = \dfrac{V_{BE}}{2} = \SI{6,0}{mL}$. Le pH lu à ce volume est le $pK_a$ du couple (environ 3,8 pour l'acide méthanoïque).
\end{corrige}

En résumé, le dosage d'un acide faible par une base forte se reconnaît à trois signatures : un départ moins acide, un saut de pH raccourci et une équivalence \emph{basique}. Son point de demi-équivalence offre en prime une mesure directe du $pK_a$. Il nous reste à voir comment rendre cette équivalence \emph{visible à l'œil nu}, sans pH-mètre : ce sera l'objet du dosage colorimétrique et du choix judicieux de l'indicateur coloré.

\coursec{Dosage d'un acide fort par une base faible}

Dans la section précédente, nous avons dosé un acide \emph{faible} (l'acide éthanoïque du vinaigre) par une base \emph{forte} (la soude) : le saut de pH était moins ample et le pH à l'équivalence était supérieur à 7, parce qu'il se formait une base faible, l'ion éthanoate. Que se passe-t-il maintenant si c'est la \emph{base} qui est faible ? C'est la situation que tu rencontres lorsqu'on dose un détartrant (acide chlorhydrique) par une solution d'ammoniac, base faible bien connue des produits ménagers. Intuitivement, on pourrait craindre qu'une base « faible » soit incapable de neutraliser complètement un acide fort. Nous allons voir que cette crainte est infondée : la réaction de dosage reste \cle{totale}, mais la \emph{courbe} change de visage, et surtout le pH à l'équivalence devient \emph{inférieur} à 7.

\courssub{La réaction de dosage : acide chlorhydrique et ammoniac}

\textbf{Situation de référence.} On place dans un bécher un volume $V_A$ d'acide chlorhydrique (\ce{H3O+} + \ce{Cl-}) de concentration $C_A$ connue approximativement, et on verse progressivement une solution d'ammoniac \ce{NH3} de concentration $C_B$ connue avec précision (solution titrante). L'ammoniac est une base faible : seule une faible partie des molécules \ce{NH3} réagit avec l'eau selon $\ce{NH3 + H2O <=> NH4+ + OH-}$, avec $K_a(\ce{NH4+}/\ce{NH3}) \approx \SI{6,3e-10}{}$, soit $pK_a \approx 9,2$.

\begin{definition}[Réaction entre un acide fort et une base faible]
La réaction entre les ions oxonium \ce{H3O+} d'un acide fort et une base faible comme l'ammoniac est \cle{totale}, \cle{rapide} et \cle{exothermique}. Son équation-bilan s'écrit :
\[ \ce{H3O+ + NH3 -> NH4+ + H2O} \]
Sa constante d'équilibre vaut $K = \dfrac{1}{K_a(\ce{NH4+}/\ce{NH3})} \approx \dfrac{1}{6,3\times10^{-10}} \approx 1,6\times10^{9} \gg 10^{4}$ : elle est donc quantitative et peut servir de \cle{réaction de dosage}, même si la base est faible.
\end{definition}

Retiens bien cette idée essentielle : la « faiblesse » de l'ammoniac concerne sa réaction \emph{avec l'eau}, qui est limitée. Mais face aux ions \ce{H3O+}, acide bien plus fort que l'eau, l'ammoniac réagit \emph{intégralement} : chaque mole de \ce{NH3} versée consomme exactement une mole de \ce{H3O+}, tant qu'il en reste. Le produit formé est l'ion ammonium \ce{NH4+}, qui est un \emph{acide faible} : c'est lui qui donnera son caractère acide à la solution à l'équivalence.

\begin{attention}[Piège fréquent]
Beaucoup d'élèves écrivent l'équation $\ce{H3O+ + OH- -> 2 H2O}$ pour \emph{tous} les dosages. C'est faux ici : la solution d'ammoniac contient très peu d'ions \ce{OH-} ; l'espèce basique majoritaire est la molécule \ce{NH3}. L'équation-bilan du dosage doit donc faire intervenir \ce{NH3}, pas \ce{OH-}. De même, ne crois jamais qu'une base faible « ne peut pas doser » un acide fort : la réaction de dosage est totale, seule l'allure de la courbe change.
\end{attention}

\courssub{Étude de la courbe pH = f(VB)}

\textbf{Protocole.} Le protocole est identique à celui des sections précédentes : burette contenant la solution d'ammoniac, bécher contenant l'acide chlorhydrique, agitation magnétique, électrode du pH-mètre plongeant dans le mélange. On relève le pH après chaque ajout de volume $V_B$ de base.

\textbf{Lecture de la courbe, zone par zone.}

\begin{itemize}
\item \textbf{Début du dosage} ($V_B = 0$) : la solution contient uniquement l'acide fort, le pH initial est très faible (pH $= 1$ pour $C_A = \SI{0,10}{mol.L^{-1}}$).
\item \textbf{Avant l'équivalence} ($V_B < V_{BE}$) : l'ammoniac versé est intégralement consommé ; il reste des ions \ce{H3O+} en excès, qui imposent un pH faible. Le pH croît lentement : c'est l'acide fort restant qui gouverne.
\item \textbf{À l'équivalence} ($V_B = V_{BE}$) : on observe un saut de pH, mais \emph{moins ample} que lors du dosage par une base forte. Le point d'équivalence $E$ se situe à un pH \emph{inférieur à 7} (typiquement $\text{pH}_E \approx 5,3$ pour des concentrations voisines de \SI{0,10}{mol.L^{-1}}). La détermination graphique de $V_{BE}$ se fait par la \textbf{méthode des tangentes parallèles} aux branches avant et après l'équivalence (ou par la méthode de la dérivée, notion).
\item \textbf{Après l'équivalence} ($V_B > V_{BE}$) : l'ammoniac versé en excès s'accumule en présence des ions \ce{NH4+} formés. On obtient un mélange du couple $\ce{NH4+}/\ce{NH3}$ : c'est une \cle{solution tampon}. Le pH croît très lentement et se stabilise au voisinage du $pK_a$ du couple, vers 9,2. Remarque que le pH \emph{ne monte jamais} vers 12--13 comme avec la soude : l'ammoniac étant une base faible, même en grand excès, elle ne peut pas rendre la solution très basique.
\end{itemize}

\begin{popfigure}
\begin{center}
\begin{tikzpicture}
\begin{axis}[
  width=0.92\linewidth, height=7.2cm,
  xlabel={$V_B$ (mL)}, ylabel={pH},
  xmin=0, xmax=50, ymin=0, ymax=14,
  xtick={0,10,20,30,40,50}, ytick={0,2,4,6,7,8,10,12,14},
  grid=major, grid style={popline!50},
  axis line style={popdark},
  label style={popdark},
]
\addplot[popcurve,domain=0:50,samples=200]{1+8.2/(1+exp(-(x-20)*0.9))};
\addplot[popline, dashed] coordinates {(20,0) (20,5.3) (0,5.3)};
\addplot[popline, dashed] coordinates {(40,0) (40,9.2) (0,9.2)};
\addplot[popline, dotted] coordinates {(0,7) (50,7)};
\node[circle,fill=poporange,inner sep=2.2pt,label={[popdark]above right:{$E\,(\text{pH}_E \approx 5,3)$}}] at (axis cs:20,5.3) {};
\node[circle,fill=poppurple,inner sep=2.2pt,label={[popdark]above left:{$I\,(\text{pH} = pK_a \approx 9,2)$}}] at (axis cs:40,9.2) {};
\node[popmuted, anchor=west] at (axis cs:41,7) {pH $=7$};
\node[popmuted, anchor=south west] at (axis cs:1,1.15) {acide fort en excès};
\node[popmuted, anchor=south east] at (axis cs:49,9.5) {tampon \ce{NH4+}/\ce{NH3}};
\end{axis}
\end{tikzpicture}
\end{center}
\legende{Dosage pH-métrique de \SI{20,0}{mL} d'acide chlorhydrique à \SI{0,10}{mol.L^{-1}} par l'ammoniac à \SI{0,10}{mol.L^{-1}}. Le saut de pH est moins ample qu'avec une base forte et le point d'équivalence $E$ se situe à $\text{pH}_E < 7$. Le point $I$ (à $V_B = 2V_{BE}$) correspond à $[\ce{NH4+}] = [\ce{NH3}]$ dans la zone tampon.}
\end{popfigure}

\begin{propriete}[Caractéristiques du dosage acide fort -- base faible]
Pour le dosage d'un acide fort par une base faible :
\begin{itemize}
\item le pH initial est très faible (imposé par l'acide fort) ;
\item le saut de pH à l'équivalence est \cle{moins marqué} que pour le dosage par une base forte ;
\item le pH à l'équivalence est \cle{inférieur à 7} ;
\item après l'équivalence, le pH plafonne au voisinage du $pK_a$ du couple de la base faible : le point $I$ de la courbe où pH $= pK_a$ (ici 9,2) correspond au mélange équimolaire $[\ce{NH4+}] = [\ce{NH3}]$ dans la zone tampon ; il est atteint pour $V_B = 2 V_{BE}$.
\end{itemize}
\end{propriete}

\courssub{Pourquoi le pH à l'équivalence est-il inférieur à 7 ?}

C'est le point le plus délicat — et le plus souvent demandé au Baccalauréat. Raisonnons sur les \emph{espèces présentes} dans le bécher à l'équivalence.

À l'équivalence, tout l'acide \ce{H3O+} et tout l'ammoniac \ce{NH3} versés ont été consommés (réaction totale). Que reste-t-il ? Des ions \ce{Cl-}, spectateurs, sans action sur l'eau, et des ions \ce{NH4+} formés par la réaction. Or l'ion ammonium est un \emph{acide faible} : il réagit partiellement avec l'eau selon :
\[ \ce{NH4+ + H2O <=> NH3 + H3O+} \]
Cette réaction libère des ions oxonium \ce{H3O+} : la solution à l'équivalence est donc \emph{acide}, d'où $\text{pH}_E < 7$.

Autrement dit, à l'équivalence, tout se passe comme si on avait préparé une solution de \cle{chlorure d'ammonium} (\ce{NH4+} + \ce{Cl-}) : une solution de sel d'acide fort et de base faible, qui est acide.

\begin{methode}[Prévoir le signe de $\text{pH}_E - 7$ sans calcul]
\begin{enumerate}
\item Écris l'équation-bilan de la réaction de dosage.
\item Identifie l'espèce acido-basique \emph{formée} à l'équivalence (l'ion ou la molécule conjuguée).
\item Conclus :
  \begin{itemize}
  \item si l'espèce formée est un \textbf{acide faible} (ex. \ce{NH4+}), alors $\text{pH}_E < 7$ ;
  \item si l'espèce formée est une \textbf{base faible} (ex. \ce{CH3COO-}), alors $\text{pH}_E > 7$ ;
  \item si l'espèce formée est \textbf{indifférente} (ex. \ce{Na+}, \ce{Cl-}, dosage acide fort--base forte), alors $\text{pH}_E = 7$.
  \end{itemize}
\end{enumerate}
Généralisation capitale : le pH à l'équivalence dépend de la \emph{nature de l'espèce conjuguée formée}. Seul le dosage acide fort--base forte donne $\text{pH}_E = 7$.
\end{methode}

\begin{popfigure}
\begin{center}
\begin{tikzpicture}[scale=1]
\draw[-{Stealth}, thick, popdark] (0,0) -- (13,0) node[below left=2pt and -30pt, popdark]{pH};
\foreach \x/\p in {1/2, 3/4, 5/6, 7/8, 9/10, 11/12}
  \draw[popdark] (\x,0.08) -- (\x,-0.08) node[below, popdark]{\small \p};
\fill[popblue!25] (0,0.35) rectangle (9.2,1.15);
\fill[poporange!30] (9.2,0.35) rectangle (13,1.15);
\draw[very thick, poppurple] (9.2,0.3) -- (9.2,1.2) node[above, poppurple]{$pK_a = 9,2$};
\node[popdark] at (4.6,0.75) {\ce{NH4+} prédomine};
\node[popdark] at (11.1,0.75) {\ce{NH3} prédomine};
\node[popmuted, anchor=west] at (0,1.6) {\small Zone de la courbe après l'équivalence : mélange tampon \ce{NH4+}/\ce{NH3}};
\draw[-{Stealth}, popmuted] (5.2,1.5) -- (8.5,1.25);
\end{tikzpicture}
\end{center}
\legende{Diagramme de prédominance du couple \ce{NH4+}/\ce{NH3}. À l'équivalence du dosage (pH $\approx 5,3$), l'espèce présente est l'acide \ce{NH4+} ; dans la zone tampon après l'équivalence, les deux espèces coexistent, avec pH $= pK_a$ lorsque $[\ce{NH4+}] = [\ce{NH3}]$ (point $I$).}
\end{popfigure}

\courssub{Exploitation quantitative de l'équivalence}

La définition de l'équivalence ne change pas : c'est le moment où les réactifs ont été mélangés dans les proportions stoechiométriques de l'équation-bilan. Comme les coefficients de \ce{H3O+} et \ce{NH3} sont tous deux égaux à 1 :

\begin{aretenir}[Relation à l'équivalence]
À l'équivalence du dosage d'un acide fort par une base faible (coefficients 1--1) :
\[ C_A \cdot V_A = C_B \cdot V_{BE} \]
Cette relation est \emph{identique} à celle du dosage acide fort--base forte : la stoechiométrie ne dépend pas de la force des réactifs, car la réaction de dosage est totale dans les deux cas.
\end{aretenir}

\begin{exemplebox}[Exemple chiffré]
On dose $V_A = \SI{25,0}{mL}$ d'acide chlorhydrique de concentration $C_A = \SI{0,080}{mol.L^{-1}}$ par une solution d'ammoniac de concentration $C_B = \SI{0,10}{mol.L^{-1}}$. Le volume équivalent est :
\[ V_{BE} = \frac{C_A \cdot V_A}{C_B} = \frac{0,080 \times 25,0}{0,10} = \SI{20,0}{mL}. \]
Vérification par les quantités de matière : $n(\ce{H3O+}) = 0,080 \times 25,0\times10^{-3} = \SI{2,0e-3}{mol}$ et $n(\ce{NH3})_{versée} = 0,10 \times 20,0\times10^{-3} = \SI{2,0e-3}{mol}$. Les quantités sont égales : c'est bien l'équivalence. À ce moment, le bécher contient \SI{2,0e-3}{mol} d'ions \ce{NH4+} dans un volume total de \SI{45,0}{mL}, soit $[\ce{NH4+}] \approx \SI{4,4e-2}{mol.L^{-1}}$ : une solution de chlorure d'ammonium, acide, ce qui confirme $\text{pH}_E < 7$.
\end{exemplebox}

\courssub{Bilan comparatif des trois types de dosage}

Rassemblons maintenant tout ce que nous avons appris dans les sections 2, 3 et 4. Ce tableau est un classique des sujets de Baccalauréat : sache le reproduire et le commenter.

\begin{popfigure}
\begin{center}
\begin{tikzpicture}
\node[anchor=north west, inner sep=0pt] at (0,0) {
\begin{tabular}{|p{3.1cm}|p{3.1cm}|p{2.2cm}|p{2.6cm}|p{2.6cm}|}
\hline
\rowcolor{popblueL} \textbf{Type de dosage} & \textbf{Exemple} & \textbf{pH à l'équivalence} & \textbf{Amplitude du saut} & \textbf{Indicateur adapté} \\
\hline
Acide fort -- base forte & \ce{H3O+} + \ce{OH-} (HCl / soude) & $\text{pH}_E = 7$ & très grand (env. 3 à 11) & BBT (ou hélianthine, phénolphtaléine) \\
\hline
\rowcolor{popcream} Acide faible -- base forte & \ce{CH3COOH} + \ce{OH-} (vinaigre / soude) & $\text{pH}_E > 7$ & moyen (env. 5 à 11) & phénolphtaléine \\
\hline
Acide fort -- base faible & \ce{H3O+} + \ce{NH3} (HCl / ammoniac) & $\text{pH}_E < 7$ & moyen (env. 3 à 9) & hélianthine \\
\hline
\end{tabular}
};
\end{tikzpicture}
\end{center}
\legende{Comparaison des trois dosages étudiés au programme. L'indicateur coloré est choisi de sorte que sa zone de virage contienne le pH à l'équivalence (nous y reviendrons en détail dans la section 5).}
\end{popfigure}

\begin{exoresolu}[Vérification d'une étiquette de solution d'acide]
Une solution commerciale d'acide chlorhydrique porte la mention « concentration : \SI{0,50}{mol.L^{-1}} ». Pour la vérifier, on dilue d'abord cette solution 10 fois, puis on dose $V_A = \SI{20,0}{mL}$ de la solution diluée par une solution d'ammoniac de concentration $C_B = \SI{0,100}{mol.L^{-1}}$. Le virage de l'hélianthine et le saut de pH sont observés pour $V_{BE} = \SI{10,2}{mL}$. L'étiquette est-elle exacte ?
\tcblower
\textbf{Équation-bilan :} $\ce{H3O+ + NH3 -> NH4+ + H2O}$, réaction totale, coefficients 1--1.

\textbf{À l'équivalence :} $C_{A,\text{dil}} \cdot V_A = C_B \cdot V_{BE}$, d'où
\[ C_{A,\text{dil}} = \frac{C_B \cdot V_{BE}}{V_A} = \frac{0,100 \times 10,2}{20,0} = \SI{5,10e-2}{mol.L^{-1}}. \]

\textbf{Remontée à la solution commerciale} (dilution par 10) :
\[ C_{A,\text{comm}} = 10 \times C_{A,\text{dil}} = \SI{0,510}{mol.L^{-1}}. \]

\textbf{Conclusion :} l'écart relatif vaut $\dfrac{0,510 - 0,50}{0,50} = 2\,\%$, inférieur à la tolérance habituelle de 5\,\% : l'étiquette est \emph{conforme}. On note au passage que le pH à l'équivalence est inférieur à 7 (formation de \ce{NH4+}), ce qui justifie le choix de l'hélianthine (zone de virage 3,1--4,4) plutôt que la phénolphtaléine.
\end{exoresolu}

\courssub{Complément Bac : dosage d'une base faible par un acide fort}

Par \emph{analogie} — et parce que les sujets de Baccalauréat adorent cette variante — étudions la situation symétrique : on place l'ammoniac dans le bécher et on verse l'acide chlorhydrique depuis la burette. La réaction de dosage est \emph{la même} :
\[ \ce{NH3 + H3O+ -> NH4+ + H2O} \qquad (K \approx 1,6\times10^{9}, \text{ totale}) \]
Mais la courbe pH $= f(V_A)$ est maintenant \cle{décroissante} :
\begin{itemize}
\item au départ, le pH est celui de la base faible : élevé mais modéré (pH $\approx 11$ pour une solution à \SI{0,1}{mol.L^{-1}}, et non 13 comme la soude) ;
\item avant l'équivalence, le mélange \ce{NH3} restant / \ce{NH4+} formé constitue un tampon : le pH décroît lentement ; à la \cle{demi-équivalence} ($V_A = V_{AE}/2$), on a $[\ce{NH3}] = [\ce{NH4+}]$, donc pH $= pK_a = 9,2$ — c'est le point de demi-équivalence $H$, qui permet de déterminer expérimentalement le $pK_a$ du couple ;
\item à l'équivalence, la solution ne contient que \ce{NH4+} (et \ce{Cl-}) : $\text{pH}_E < 7$, comme précédemment ;
\item après l'équivalence, l'acide fort en excès impose un pH très faible.
\end{itemize}

\begin{attention}[Deux points à ne pas confondre]
Le point de demi-équivalence $H$ (pH $= pK_a$) et le point d'équivalence $E$ (saut de pH) sont \emph{distincts} : $H$ se situe à $V_{AE}/2$ sur la partie tampon de la courbe (avant l'équivalence), $E$ au milieu du saut. Confondre leurs abscisses est une erreur classique qui fausse la détermination du $pK_a$. Repère toujours d'abord $E$ (méthode des tangentes), puis divise son abscisse par deux pour trouver $H$.
\end{attention}

\begin{savaistu}[Le chlorure d'ammonium, un sel aux mille usages]
Le sel formé à l'équivalence de ce dosage, le chlorure d'ammonium \ce{NH4Cl}, n'est pas qu'une curiosité de laboratoire. Il est utilisé comme \emph{engrais azoté} dans certaines cultures (l'ion \ce{NH4+} est une source d'azote assimilable par les plantes), comme électrolyte dans les \emph{piles sèches} Leclanché — les piles « bâton » classiques que l'on trouve encore dans les boutiques de quartier à Yaoundé ou Douala — et comme flux de soudure en métallurgie. Son caractère acide en solution, que tu viens de démontrer, est d'ailleurs mis à profit pour nettoyer les surfaces métalliques avant soudure : il dissout les oxydes !
\end{savaistu}

\begin{aretenir}[L'essentiel de la section]
\begin{itemize}
\item La réaction $\ce{H3O+ + NH3 -> NH4+ + H2O}$ est totale ($K = 1/K_a \gg 10^{4}$) : une base faible dose parfaitement un acide fort.
\item La courbe pH $= f(V_B)$ présente un saut moins ample qu'avec une base forte, et $\text{pH}_E < 7$ car il se forme l'acide faible \ce{NH4+}.
\item La relation à l'équivalence reste $C_A V_A = C_B V_{BE}$ : la stoechiométrie ne dépend pas de la force des réactifs.
\item Le pH à l'équivalence dépend de l'espèce conjuguée formée : acide faible $\Rightarrow$ $\text{pH}_E < 7$ ; base faible $\Rightarrow$ $\text{pH}_E > 7$ ; espèce indifférente $\Rightarrow$ $\text{pH}_E = 7$.
\item Pour le dosage inverse (base faible par acide fort), la courbe est décroissante et le point de demi-équivalence $H$ (à $V_{AE}/2$) donne pH $= pK_a$.
\end{itemize}
\end{aretenir}

\coursec{Dosage colorimétrique et choix de l'indicateur}

Dans la section précédente, nous avons suivi le dosage de l'ammoniac par l'acide chlorhydrique à l'aide d'un pH-mètre : point par point, nous avons tracé la courbe $\mathrm{pH} = f(V_A)$ et repéré l'équivalence par la méthode des tangentes. Cette méthode \cle{pH-métrique} est précise, mais elle exige un appareil étalonné, du temps et de la patience. Or, au laboratoire comme dans une savonnerie artisanale ou une brasserie, on a souvent besoin d'un résultat \emph{rapide} et \emph{suffisamment exact}. C'est là qu'intervient le \cle{dosage colorimétrique} : on remplace l'œil électronique du pH-mètre par l'œil du chimiste, aidé d'une substance qui change de couleur au bon moment : l'\cle{indicateur coloré acido-basique}. Toute cette section répond à deux questions : comment ça marche, et surtout \emph{comment choisir le bon indicateur} — car c'est l'une des questions les plus fréquentes au Baccalauréat.

\courssub{Principe du dosage colorimétrique}

\textbf{Idée intuitive.} Imagine que tu verses progressivement de la soude dans un vinaigre. Tant que l'équivalence n'est pas atteinte, il reste de l'acide ; dès qu'elle est dépassée, la soude est en excès. Le pH, lui, « saute » brutalement au voisinage de l'équivalence. Si l'on pouvait \emph{voir} ce saut, on connaîtrait le volume équivalent sans aucun appareil. C'est exactement le rôle de l'indicateur coloré : c'est un « révélateur de pH » qui change de couleur dans une zone étroite de pH.

\begin{definition}[Indicateur coloré acido-basique]
Un \cle{indicateur coloré acido-basique} est un couple acide/base, noté $\ce{HIn}/\ce{In-}$, dont la forme acide $\ce{HIn}$ et la forme basique $\ce{In-}$ \textbf{n'ont pas la même couleur}. L'équilibre mis en jeu s'écrit :
\[ \ce{HIn(aq) + H2O(l) <=> In-(aq) + H3O+(aq)} \qquad \text{avec } K_a = \frac{[\ce{In-}][\ce{H3O+}]}{[\ce{HIn}]} \]
En milieu très acide, l'équilibre est déplacé vers $\ce{HIn}$ : on voit la \textbf{teinte acide}. En milieu très basique, il est déplacé vers $\ce{In-}$ : on voit la \textbf{teinte basique}.
\end{definition}

\textbf{Pourquoi le changement de couleur n'est-il pas instantané ?} L'œil humain ne perçoit la teinte acide pure que lorsque $\ce{HIn}$ domine largement (environ $[\ce{HIn}] > 10\,[\ce{In-}]$, soit $\mathrm{pH} < \mathrm{p}K_a - 1$), et la teinte basique pure que lorsque $\ce{In-}$ domine ($\mathrm{pH} > \mathrm{p}K_a + 1$). Entre les deux, on observe une couleur intermédiaire, dite \textbf{teinte sensible}.

\begin{definition}[Zone de virage]
La \cle{zone de virage} d'un indicateur coloré est l'intervalle de pH dans lequel il change progressivement de couleur. Elle encadre le $\mathrm{p}K_a$ du couple $\ce{HIn}/\ce{In-}$ :
\[ \text{zone de virage} \approx \mathrm{p}K_a(\ce{HIn}/\ce{In-}) \pm 1 \]
Au milieu de la zone ($\mathrm{pH} = \mathrm{p}K_a$), les deux formes sont présentes en quantités égales : c'est la teinte sensible.
\end{definition}

\begin{attention}[L'indicateur est lui-même un acide !]
Un indicateur coloré est un \emph{acide faible} : il consomme donc une petite quantité du réactif titrant. C'est pourquoi on n'introduit que \textbf{2 à 3 gouttes} d'indicateur dans le bécher. En mettre davantage fausserait le volume équivalent — et ternirait la couleur, rendant le virage moins net. Question classique du Bac : « Pourquoi ne verse-t-on que quelques gouttes d'indicateur ? »
\end{attention}

\courssub{Les trois indicateurs du programme}

Le programme de Terminale en retient trois, qu'il faut connaître avec leurs zones de virage et leurs couleurs :

\begin{center}
\begin{tabularx}{\linewidth}{|l|c|X|X|c|}
\hline
\rowcolor{popblueL} \textbf{Indicateur} & $\mathrm{p}K_a$ & \textbf{Teinte acide} & \textbf{Teinte basique} & \textbf{Zone de virage} \\
\hline
Hélianthine (méthylorange) & 3,7 & rouge & jaune-orange & 3,1 -- 4,4 \\
\hline
Bleu de bromothymol (BBT) & 7,0 & jaune & bleu & 6,0 -- 7,6 \\
\hline
Phénolphtaléine & 9,4 & incolore & rose (fuchsia) & 8,2 -- 10,0 \\
\hline
\end{tabularx}
\end{center}

\begin{popfigure}
\begin{center}
\begin{tikzpicture}[x=0.78cm]
% axe pH
\draw[popdark, thick, -{Stealth}] (0,0) -- (14.6,0) node[below left=2pt and -14pt]{\small pH};
\foreach \x in {0,1,...,14} {
  \draw[popdark] (\x,0) -- (\x,-0.12) node[below]{\scriptsize \x};
}
% Hélianthine
\fill[red!70] (0,0.45) rectangle (3.1,0.95);
\fill[orange!80!yellow] (3.1,0.45) rectangle (4.4,0.95);
\fill[yellow!80!orange] (4.4,0.45) rectangle (14,0.95);
\draw[popdark] (0,0.45) rectangle (14,0.95);
\draw[popdark, dashed] (3.1,0.45) -- (3.1,0.95);
\draw[popdark, dashed] (4.4,0.45) -- (4.4,0.95);
\node[right, popdark] at (14.15,0.7) {\small Hélianthine};
\node[white] at (1.5,0.7) {\scriptsize rouge};
\node[popdark] at (9,0.7) {\scriptsize jaune-orange};
\node[popdark] at (3.75,1.12) {\scriptsize 3,1 -- 4,4};
% BBT
\fill[yellow!85!orange] (0,1.45) rectangle (6,1.95);
\fill[popgreen!60] (6,1.45) rectangle (7.6,1.95);
\fill[popblue!80] (7.6,1.45) rectangle (14,1.95);
\draw[popdark] (0,1.45) rectangle (14,1.95);
\draw[popdark, dashed] (6,1.45) -- (6,1.95);
\draw[popdark, dashed] (7.6,1.45) -- (7.6,1.95);
\node[right, popdark] at (14.15,1.7) {\small BBT};
\node[popdark] at (3,1.7) {\scriptsize jaune};
\node[white] at (11,1.7) {\scriptsize bleu};
\node[popdark] at (6.8,2.12) {\scriptsize 6,0 -- 7,6};
% Phénolphtaléine
\fill[popcream] (0,2.45) rectangle (8.2,2.95);
\fill[poppink!40] (8.2,2.45) rectangle (10,2.95);
\fill[poppink!85!popink] (10,2.45) rectangle (14,2.95);
\draw[popdark] (0,2.45) rectangle (14,2.95);
\draw[popdark, dashed] (8.2,2.45) -- (8.2,2.95);
\draw[popdark, dashed] (10,2.45) -- (10,2.95);
\node[right, popdark] at (14.15,2.7) {\small Phénolphtaléine};
\node[popmuted] at (4,2.7) {\scriptsize incolore};
\node[white] at (12,2.7) {\scriptsize rose};
\node[popdark] at (9.1,3.12) {\scriptsize 8,2 -- 10,0};
\end{tikzpicture}
\end{center}
\legende{Zones de virage des trois indicateurs du programme sur l'échelle de pH. La zone colorée intermédiaire (entre les pointillés) est la zone de virage, centrée sur le $\mathrm{p}K_a$ de l'indicateur.}
\end{popfigure}

\begin{popfigure}
\begin{center}
\begin{tikzpicture}[scale=0.9]
% Bécher 1 : hélianthine
\draw[popdark, thick] (-5.2,0) -- (-5.2,1.5) (-3.8,0) -- (-3.8,1.5) (-5.2,0) arc (180:360:0.7);
\fill[red!60] (-5.13,0.08) rectangle (-3.87,0.85);
\node[popdark] at (-4.5,-0.55) {\scriptsize hélianthine};
\node[popdark] at (-4.5,-0.95) {\scriptsize pH $< 3{,}1$ : rouge};
% Bécher 2 : BBT
\draw[popdark, thick] (-1.2,0) -- (-1.2,1.5) (0.2,0) -- (0.2,1.5) (-1.2,0) arc (180:360:0.7);
\fill[popgreen!55] (-1.13,0.08) rectangle (0.13,0.85);
\node[popdark] at (-0.5,-0.55) {\scriptsize BBT};
\node[popdark] at (-0.5,-0.95) {\scriptsize pH $\approx 7$ : vert};
% Bécher 3 : phénolphtaléine
\draw[popdark, thick] (2.8,0) -- (2.8,1.5) (4.2,0) -- (4.2,1.5) (2.8,0) arc (180:360:0.7);
\fill[poppink!70] (2.87,0.08) rectangle (4.13,0.85);
\node[popdark] at (3.5,-0.55) {\scriptsize phénolphtaléine};
\node[popdark] at (3.5,-0.95) {\scriptsize pH $> 10$ : rose};
\end{tikzpicture}
\end{center}
\legende{Trois béchers témoins : chaque indicateur exhibe une couleur caractéristique selon le pH du milieu. Au cours d'un dosage, on observe le passage net d'une teinte à l'autre.}
\end{popfigure}

\begin{savaistu}[Le chou rouge, indicateur du marché]
Pas besoin de laboratoire pour faire de la chimie des indicateurs ! Le jus de \textbf{chou rouge}, que l'on trouve dans les marchés de Yaoundé ou de Douala, contient des \emph{anthocyanes} : il est \textbf{rouge} en milieu acide (ajoute du jus de citron ou du vinaigre), \textbf{violet} vers la neutralité et \textbf{vert-bleu} en milieu basique (ajoute un peu de lessive de cendres ou de savon dissous). Le jus de bissap, lui aussi, vire au rouge vif en présence d'acide : c'est pour cela qu'on l'acidifie naturellement avec du citron. La chimie est dans la cuisine !
\end{savaistu}

\courssub{La règle d'or du choix de l'indicateur}

\begin{aretenir}[Règle de choix]
Un indicateur coloré convient pour un dosage si et seulement si sa \textbf{zone de virage contient le pH à l'équivalence} $\mathrm{pH}_E$ du dosage — ou, de façon équivalente, si sa zone de virage est \textbf{incluse dans le saut de pH} de la courbe. Le virage (changement de couleur) signale alors l'équivalence : $V_{BE}$ lu = $V_E$ réel.
\end{aretenir}

\begin{methode}[Choisir judicieusement un indicateur coloré]
\begin{enumerate}
\item \textbf{Identifier} la nature des réactifs : acide fort ou faible ? base forte ou faible ?
\item \textbf{Déterminer le pH à l'équivalence} $\mathrm{pH}_E$ : par lecture graphique sur la courbe $\mathrm{pH} = f(V)$, ou par le raisonnement sur la nature de la solution à l'équivalence (sel d'acide fort et de base forte : $\mathrm{pH}_E = 7$ ; base faible restante : $\mathrm{pH}_E > 7$ ; acide faible restant : $\mathrm{pH}_E < 7$).
\item \textbf{Comparer} $\mathrm{pH}_E$ aux zones de virage des indicateurs disponibles.
\item \textbf{Retenir} l'indicateur dont la zone de virage contient $\mathrm{pH}_E$ ; en cas d'hésitation entre plusieurs, choisir celui dont la zone est la mieux centrée sur $\mathrm{pH}_E$ et entièrement comprise dans le saut de pH.
\end{enumerate}
\end{methode}

Appliquons cette règle aux trois dosages étudiés dans les sections précédentes :

\begin{center}
\begin{tabularx}{\linewidth}{|X|c|c|X|}
\hline
\rowcolor{popblueL} \textbf{Dosage} & $\mathrm{pH}_E$ & \textbf{Indicateur conseillé} & \textbf{Remarque} \\
\hline
Acide fort -- base forte & 7,0 & BBT & saut de pH très grand (3 à 11) : les trois indicateurs conviennent, le BBT est idéal \\
\hline
Acide faible -- base forte & $> 7$ (souvent 8 -- 9) & phénolphtaléine & le saut démarre vers pH 5--6 : l'hélianthine et le BBT virent beaucoup trop tôt \\
\hline
Acide fort -- base faible & $< 7$ (souvent 5 -- 6) & \textbf{BBT} & le saut couvre 3--8 ; BBT (6,0--7,6) est dans la partie haute du saut, hélianthine (3,1--4,4) dans la partie basse. BBT est plus centré sur $\mathrm{pH}_E$ \\
\hline
\end{tabularx}
\end{center}

\begin{popfigure}
\begin{center}
\begin{tikzpicture}
\begin{axis}[
  width=0.92\linewidth, height=7.2cm,
  xlabel={$V_B$ (mL)}, ylabel={pH},
  xmin=0, xmax=30, ymin=0, ymax=14,
  xtick={0,5,10,15,20,25,30}, ytick={0,2,4,6,8,10,12,14},
  grid=both, grid style={popline!40},
  axis line style={popdark}, label style={popdark}, tick label style={popdark},
]
% bandes de virage
\fill[poppink!35] (axis cs:0,8.2) rectangle (axis cs:30,10);
\fill[popgreen!25] (axis cs:0,6) rectangle (axis cs:30,7.6);
\fill[orange!30] (axis cs:0,3.1) rectangle (axis cs:30,4.4);
\node[popink, anchor=west] at (axis cs:21,9.1) {\scriptsize phénolphtaléine};
\node[popgreen!60!black, anchor=west] at (axis cs:21,6.8) {\scriptsize BBT};
\node[poporange!80!black, anchor=west] at (axis cs:21,3.75) {\scriptsize hélianthine};
% courbe acide faible - base forte (Ve = 20)
\addplot[popcurve, domain=0.1:19.4, samples=150] {4.8 + log10(x/(20-x))};
\addplot[popcurve, domain=20.6:30, samples=80] {14 + log10((x-20)*0.1/(x+20))}; % approximation post-equivalence
\addplot[popcurve, domain=19.4:20.6, samples=40] {8.7 + 3.3*(x-20)};
% point d'équivalence
\addplot[only marks, mark=*, mark size=2pt, popink] coordinates {(20,8.7)};
\draw[popink, dashed] (axis cs:20,0) -- (axis cs:20,8.7) -- (axis cs:0,8.7);
\node[popink, anchor=south west] at (axis cs:20.3,8.75) {\scriptsize $E\,(\mathrm{pH}_E \approx 8{,}7)$};
\end{axis}
\end{tikzpicture}
\end{center}
\legende{Dosage d'un acide faible ($\mathrm{p}K_a = 4{,}8$) par la soude : seule la zone de virage de la phénolphtaléine contient $\mathrm{pH}_E \approx 8{,}7$ et coupe la courbe dans son saut quasi vertical. Le BBT et l'hélianthine croisent la courbe bien avant l'équivalence, là où elle est presque horizontale : l'erreur sur $V_{BE}$ serait énorme.}
\end{popfigure}

La figure ci-dessus montre un point capital souvent mal compris : ce n'est pas seulement la valeur de $\mathrm{pH}_E$ qui compte, c'est le fait que le virage se produise \textbf{dans la partie quasi verticale} du saut de pH. Là, une seule goutte de réactif fait varier le pH de plusieurs unités : l'indicateur vire d'un coup, et le volume lu est précis à la goutte près.

\courssub{Conséquences d'un mauvais choix d'indicateur}

\begin{attention}[Erreur fréquente n°1 : la phénolphtaléine avec une base faible]
Au Bac, beaucoup de candidats écrivent machinalement « on ajoute de la phénolphtaléine » quel que soit le dosage. Catastrophe ! Pour le dosage de l'\textbf{ammoniac} (base faible) par l'\textbf{acide chlorhydrique}, $\mathrm{pH}_E \approx 5$--$6$. Or la phénolphtaléine vire entre 8,2 et 10 : en versant l'acide, elle devient incolore \emph{bien avant} l'équivalence. Le volume $V_{AE}$ mesuré est trop faible, donc la concentration calculée de la base est \textbf{sous-estimée}. L'indicateur rigoureux est le \textbf{BBT} (zone 6,0--7,6 incluse dans le saut de pH) ; l'hélianthine (zone 3,1--4,4) vire après l'équivalence et surestime la concentration.
\end{attention}

\begin{attention}[Erreur fréquente n°2 : raisonner « à l'envers »]
Quand on dose une \emph{base} placée dans le bécher par un \emph{acide} versé depuis la burette, le pH \emph{diminue} : on part de la teinte basique pour aller vers la teinte acide. Le virage attendu est alors, par exemple, « passage du jaune-orange au rouge » pour l'hélianthine — et non l'inverse. Toujours annoncer le changement de couleur \textbf{dans le sens du dosage réalisé}.
\end{attention}

De manière générale, un indicateur mal choisi vire \emph{avant} ou \emph{après} l'équivalence : le volume équivalent lu $V_{BE}$ est erroné, et comme $C_A = \dfrac{C_B \cdot V_{BE}}{V_A}$, toute la chaîne de calcul s'effondre. En revanche, dans le cas acide fort--base forte, le saut de pH est si grand (de 3 à 11 environ) que les trois indicateurs virent quasiment à la même goutte : le choix y est moins critique, le BBT restant le plus rigoureux.

\courssub{Exemple chiffré : le degré d'acidité d'un vinaigre}

\begin{exoresolu}[Choix de l'indicateur pour doser un vinaigre]
\textbf{Énoncé.} On dose $V_A = \SI{10,0}{mL}$ d'une solution de vinaigre dilué au dixième (acide éthanoïque, $\mathrm{p}K_a(\ce{CH3COOH}/\ce{CH3COO-}) = 4{,}8$) par de la soude à $C_B = \SI{0,10}{mol.L^{-1}}$. Le suivi pH-métrique donne un point d'équivalence à $V_{BE} = \SI{12,0}{mL}$ avec $\mathrm{pH}_E \approx 8{,}7$. On dispose de l'hélianthine, du BBT et de la phénolphtaléine.
\begin{enumerate}
\item Quel indicateur choisir ? Justifier et décrire le virage observé.
\item Calculer la concentration $C_A$ du vinaigre dilué, puis celle du vinaigre commercial.
\item Que se serait-il passé avec le BBT ?
\end{enumerate}
\tcblower
\textbf{Solution détaillée.}
\begin{enumerate}
\item La réaction de dosage est :
\[ \ce{CH3COOH(aq) + OH-(aq) -> CH3COO-(aq) + H2O(l)} \]
réaction quasi totale ($K = 10^{14-4{,}8} = 10^{9{,}2} \gg 10^4$) et rapide : c'est bien une réaction de dosage. À l'équivalence, la solution contient la base faible \ce{CH3COO-}, d'où $\mathrm{pH}_E \approx 8{,}7 > 7$. Seule la \textbf{phénolphtaléine} a une zone de virage (8,2 -- 10,0) qui contient $\mathrm{pH}_E$. Le virage observé : la solution, \emph{incolore}, devient \emph{rose pâle persistant} à la première goutte de soude en excès.
\item À l'équivalence : $C_A \cdot V_A = C_B \cdot V_{BE}$, donc
\[ C_A = \frac{C_B \cdot V_{BE}}{V_A} = \frac{0{,}10 \times 12{,}0}{10{,}0} = \SI{0,12}{mol.L^{-1}}. \]
Le vinaigre commercial, dix fois plus concentré : $C = 10 \times 0{,}12 = \SI{1,2}{mol.L^{-1}}$, soit environ $1{,}2 \times 60 = \SI{72}{g.L^{-1}}$ d'acide éthanoïque, c'est-à-dire un vinaigre à $7{,}2$ degrés — conforme à l'étiquette d'un vinaigre de table classique.
\item Avec le BBT (zone 6,0 -- 7,6), le virage du jaune au bleu se produirait \emph{avant} l'équivalence, sur la partie ascendante de la courbe : on lirait $V_{BE} < \SI{12,0}{mL}$, et la concentration trouvée serait \textbf{sous-estimée}. Le BBT est donc à rejeter ici.
\end{enumerate}
\end{exoresolu}

\courssub{Colorimétrie ou pH-métrie : lequel choisir ?}

\begin{center}
\begin{tabularx}{\linewidth}{|X|X|}
\hline
\rowcolor{popblueL} \textbf{Dosage pH-métrique} & \textbf{Dosage colorimétrique} \\
\hline
Précis : $V_E$ déterminé à la goutte près par la méthode des tangentes & Rapide : une seule manipulation, résultat immédiat \\
\hline
Fournit la courbe complète (accès au $\mathrm{p}K_a$ par la demi-équivalence) & Ne donne que le point d'équivalence \\
\hline
Nécessite un pH-mètre étalonné, coûteux et fragile & Nécessite seulement quelques gouttes d'indicateur \\
\hline
Indispensable si la solution est colorée ou trouble & Inutilisable si la solution masque la couleur de l'indicateur \\
\hline
\end{tabularx}
\end{center}

\begin{aretenir}[L'essentiel de la section]
\begin{itemize}
\item Un indicateur coloré est un couple acide/base $\ce{HIn}/\ce{In-}$ dont les deux formes ont des couleurs différentes ; sa zone de virage vaut environ $\mathrm{p}K_a \pm 1$.
\item Règle d'or : la zone de virage doit \textbf{contenir le pH à l'équivalence} (et couper le saut de pH).
\item Acide fort--base forte : BBT (les trois passent). Acide faible--base forte : phénolphtaléine. Acide fort--base faible : \textbf{BBT} (privilégié) ou hélianthine.
\item On ne verse que 2 à 3 gouttes d'indicateur, car c'est un acide faible qui consomme du réactif.
\item Un mauvais choix décale le virage, fausse $V_{BE}$ et donc la concentration calculée.
\end{itemize}
\end{aretenir}

\begin{exercice}[À toi de jouer]
On dose $\SI{20,0}{mL}$ d'une solution d'ammoniac par de l'acide chlorhydrique à $\SI{0,050}{mol.L^{-1}}$. Le pH à l'équivalence vaut 5,5 et le volume équivalent est $\SI{15,0}{mL}$.
\begin{enumerate}
\item Écrire l'équation-bilan de la réaction de dosage.
\item Choisir, en justifiant, l'indicateur coloré approprié et décrire le virage observé.
\item Calculer la concentration de la solution d'ammoniac.
\end{enumerate}
\end{exercice}

\begin{corrige}[Exercice : dosage de l'ammoniac]
\begin{enumerate}
\item $\ce{NH3(aq) + H3O+(aq) -> NH4+(aq) + H2O(l)}$, réaction quasi totale ($K = 1/\mathrm{K}_a(\ce{NH4+}/\ce{NH3}) = 10^{9{,}2} \gg 10^4$) et rapide.
\item $\mathrm{pH}_E = 5{,}5$. Le saut de pH s'étend approximativement de 3,7 à 7,3. Le \textbf{BBT} (zone 6,0 -- 7,6) a sa zone de virage incluse dans la partie haute de ce saut, au plus près de $\mathrm{pH}_E$ : c'est l'indicateur optimal. L'hélianthine (zone 3,1 -- 4,4) se situe dans la partie basse du saut, elle vire après l'équivalence et surestimerait la concentration. La phénolphtaléine vire bien avant l'équivalence. Avec le BBT, on observe le passage du \textbf{bleu} (milieu basique initial) au \textbf{jaune} (milieu acide en excès) ; au point d'équivalence, la teinte est \textbf{verte} (mélange bleu/jaune).
\item $C_{\ce{NH3}} = \dfrac{C_{\ce{HCl}} \cdot V_{AE}}{V_{\ce{NH3}}} = \dfrac{0{,}050 \times 15{,}0}{20{,}0} = \SI{3,75e-2}{mol.L^{-1}}$.
\end{enumerate}
\end{corrige}

\coursec{Applications : vérification d'étiquettes et degré d'acidité}

Dans les sections précédentes, tu as appris à choisir judicieusement un indicateur coloré pour repérer l'équivalence d'un dosage et à tracer les courbes pH-métriques caractéristiques. Il est temps de mettre tout cela en pratique : les dosages acido-basiques ne sont pas des exercices scolaires artificiels, ce sont des outils quotidiens des laboratoires de contrôle qualité. Quand un industriel annonce sur une bouteille de vinaigre « acidité : 6° », ou qu'un fabricant de déboucheur liquide affiche « hydroxyde de sodium à 4,0 mol·L⁻¹ », comment vérifier que l'étiquette dit vrai ? C'est exactement la démarche que nous allons construire ici, pas à pas, avec deux problématiques concrètes : la vérification de l'étiquette d'une solution commerciale, puis la détermination du degré d'acidité d'un vinaigre.

\courssub{Vérifier l'étiquette d'une solution commerciale}

\textbf{La problématique.} Au laboratoire du lycée, on trouve souvent des flacons commerciaux : « acide chlorhydrique concentré, environ 37 \% en masse », « détergent contenant de l'hydroxyde de sodium », etc. Ces indications sont parfois approximatives, et le produit peut avoir vieilli (la soude, par exemple, se carbonate lentement à l'air en présence de CO₂ : \ce{2 NaOH + CO2 -> Na2CO3 + H2O}). Le technicien doit donc \cle{vérifier} la concentration réelle par un dosage.

\textbf{Pourquoi faut-il diluer d'abord ?} Intuition : essaie de doser directement une solution d'acide chlorhydrique à \SI{12}{mol.L^{-1}} avec une burette de \SI{25}{mL} de soude \SI{0,10}{mol.L^{-1}}... il faudrait prélever moins d'un millilitre de solution commerciale, volume impossible à mesurer avec précision à la pipette jaugée (volume minimal usuel : \SI{5,0}{mL} ou \SI{10,0}{mL}), et la réaction serait dangereusement exothermique. La solution : \cle{diluer} la solution commerciale, c'est-à-dire prélever un volume $V_0$ connu (à la pipette jaugée) et l'introduire dans une fiole jaugée de volume $V$, complétée avec de l'eau distillée.

\begin{definition}[Facteur de dilution]
Le \cle{facteur de dilution} est le rapport $F = \dfrac{C_0}{C} = \dfrac{V}{V_0}$, où $C_0$ est la concentration molaire de la solution commerciale (solution mère) et $C$ celle de la solution diluée (solution fille). On dit « dilution au $1/F$ » ou « dilution $F$ fois ».
\end{definition}

\begin{exemplebox}[Préparation d'une dilution au 1/20]
On veut diluer 20 fois un acide chlorhydrique commercial. On prélève $V_0 = \SI{5,0}{mL}$ à la pipette jaugée, on les verse dans une fiole jaugée de $V = \SI{100,0}{mL}$ contenant déjà un peu d'eau distillée, on complète au trait de jauge et on homogénéise. Vérification : $F = \dfrac{100,0}{5,0} = 20$. Si la solution commerciale vaut $C_0 = \SI{1,0}{mol.L^{-1}}$, la solution diluée vaut $C = \SI{0,050}{mol.L^{-1}}$, parfaitement dosable par une base de concentration voisine.
\end{exemplebox}

\begin{attention}[Toujours verser l'acide dans l'eau]
Pour diluer un acide concentré, on verse \textbf{l'acide dans l'eau}, jamais l'inverse : la dissolution est très exothermique et les premières gouttes d'eau sur l'acide concentré provoqueraient des projections brûlantes. Port de la blouse, des lunettes et des gants obligatoire.
\end{attention}

\textbf{Exploitation du dosage.} On dose ensuite un volume $V_A$ de la solution \emph{diluée} par une base titrée de concentration $C_B$ (suivi pH-métrique ou indicateur coloré). À l'équivalence :
\[ C_A \cdot V_A = C_B \cdot V_{BE} \quad\Longrightarrow\quad C_A = \frac{C_B \cdot V_{BE}}{V_A} \]
puis on \emph{remonte} à la solution commerciale : $C_0 = F \times C_A$. On compare enfin $C_0$ à la valeur annoncée sur l'étiquette.

\begin{methode}[Vérifier une étiquette : la démarche en 5 étapes]
\begin{enumerate}
\item \textbf{Diluer} la solution commerciale d'un facteur $F$ connu (pipette jaugée + fiole jaugée).
\item \textbf{Doser} un volume $V_A$ de solution diluée par le réactif titrant (suivi pH-métrique ou indicateur coloré).
\item \textbf{Repérer} l'équivalence : $V_{BE}$ (méthode des tangentes parallèles ou virage de l'indicateur).
\item \textbf{Calculer} $C_A = \dfrac{C_B V_{BE}}{V_A}$ puis $C_0 = F \times C_A$.
\item \textbf{Comparer} à l'étiquette et conclure de façon critique (écart relatif, incertitudes).
\end{enumerate}
\end{methode}

\courssub{Le degré d'acidité d'un vinaigre}

\textbf{La problématique.} Le vinaigre est une solution aqueuse d'\cle{acide éthanoïque} \ce{CH3COOH}, obtenue par oxydation de l'éthanol des boissons fermentées (vin, vin de palme, bili-bili, bière) par des bactéries du genre \emph{Acetobacter}. Au Cameroun comme en France, la réglementation impose une acidité minimale (souvent 5° ou 6°) pour qu'un produit porte le nom « vinaigre ». Mais que signifie ce « degré » ?

\begin{definition}[Degré d'acidité d'un vinaigre (degré acétimétrique)]
Le \cle{degré d'acidité} d'un vinaigre est la \textbf{masse, en grammes, d'acide éthanoïque \ce{CH3COOH} contenue dans \SI{100}{g} de vinaigre}. Un vinaigre à 6° contient donc \SI{6}{g} d'acide éthanoïque pour \SI{100}{g} de vinaigre.
\end{definition}

Comme la masse volumique du vinaigre est très proche de celle de l'eau ($\rho \approx \SI{1,00}{g.mL^{-1}}$), on assimile \SI{100}{g} de vinaigre à \SI{100}{mL} : le degré correspond alors numériquement à la concentration massique $C_m$ exprimée en \si{g} pour \SI{100}{mL}. Si $C_m$ est en \si{g.L^{-1}}, on a $d = \dfrac{C_m}{10}$.

\textbf{Protocole expérimental.} Le vinaigre commercial (environ \SI{1}{mol.L^{-1}}) est trop concentré et souvent coloré : on le dilue au 1/10 et on dose la solution diluée.

\begin{experience}[Dosage du vinaigre dilué]
\textbf{Dispositif et protocole :}
\begin{enumerate}
\item Pipetter $V_0 = \SI{10,0}{mL}$ de vinaigre commercial, les introduire dans une fiole jaugée de \SI{100,0}{mL}, compléter à l'eau distillée : solution $S$ diluée 10 fois ($F = 10$).
\item Prélever $V_A = \SI{10,0}{mL}$ de $S$ dans un erlenmeyer ; ajouter 2 gouttes de phénolphtaléine (ou plonger la sonde d'un pH-mètre).
\item Remplir la burette avec la soude titrée à $C_B = \SI{0,050}{mol.L^{-1}}$.
\item Verser la soude par fractions décroissantes en agitant ; noter le volume $V_{BE}$ au virage persistant de la phénolphtaléine (rose pâle), ou tracer $pH = f(V_B)$.
\end{enumerate}
\textbf{Observations :} la solution vire au rose pâle vers $V_B \approx \SI{13}{mL}$ ; la courbe pH-métrique montre un saut de pH centré sur l'équivalence, avec $pH_E > 7$ (acide faible dosé par une base forte).
\textbf{Interprétation :} la réaction de dosage est totale :
\[ \ce{CH3COOH + OH- -> CH3COO- + H2O} \]
\end{experience}

\begin{popfigure}
\begin{center}
\begin{tikzpicture}[scale=0.9, every node/.style={font=\small}]
% Etape 1 : bouteille vinaigre
\draw[thick, popdark] (0,0.4) -- (0,1.8) -- (0.25,2.1) -- (0.25,2.6) -- (0.75,2.6) -- (0.75,2.1) -- (1,1.8) -- (1,0.4) -- cycle;
\fill[popgold!50] (0.06,0.46) -- (0.06,1.3) -- (0.94,1.3) -- (0.94,0.46) -- cycle;
\node[align=center] at (0.5,-0.3) {Vinaigre\\ commercial};
% pipette
\draw[->, thick, popblue] (1.2,1.6) -- (2.2,1.6) node[midway, above]{pipette \SI{10,0}{mL}};
% Etape 2 : fiole jaugee
\begin{scope}[xshift=2.6cm]
\draw[thick, popdark] (0.35,2.6) -- (0.35,1.7) (0.75,2.6) -- (0.75,1.7);
\draw[thick, popdark] (0.35,1.7) .. controls (0.35,1.2) and (0,1.1) .. (0,0.5) .. controls (0,0) and (1.1,0) .. (1.1,0.5) .. controls (1.1,1.1) and (0.75,1.2) .. (0.75,1.7);
\draw[popdark] (0.33,2.2) -- (0.77,2.2) node[midway, right] {trait de jauge};
\fill[popgold!30] (0.08,0.5) .. controls (0.08,0.12) and (1.02,0.12) .. (1.02,0.5) .. controls (1.02,0.9) and (0.08,0.9) .. cycle;
\node[align=center] at (0.55,-0.3) {Fiole jaugée \SI{100}{mL}\\ dilution $1/10$};
\end{scope}
\draw[->, thick, popblue] (3.95,1.6) -- (4.95,1.6) node[midway, above]{prélèvement \SI{10,0}{mL}};
% Etape 3 : erlenmeyer + burette
\begin{scope}[xshift=5.4cm]
\draw[thick, popdark] (0.9,3.3) -- (0.9,1.9) (1.15,3.3) -- (1.15,1.9);
\draw[thick, popdark] (0.9,1.9) -- (0.9,1.75) .. controls (0.9,1.6) .. (1.02,1.55);
\fill[popblue!40] (0.92,2.6) rectangle (1.13,3.25);
\draw[thick, popdark] (0.35,1.15) -- (0.85,0.05) -- (1.7,0.05) -- (2.2,1.15);
\draw[thick, popdark] (0.35,1.15) -- (0.35,1.35) (2.2,1.15) -- (2.2,1.35);
\fill[poppink!40] (0.42,1.1) -- (0.88,0.12) -- (1.67,0.12) -- (2.13,1.1) -- cycle;
\node[align=center] at (1.27,-0.5) {Dosage par \ce{NaOH} \SI{0,050}{mol.L^{-1}}\\ phénolphtaléine ou pH-mètre};
\node[font=\footnotesize, popdark] at (1.55,2.9) {\ce{NaOH}};
\end{scope}
\end{tikzpicture}
\legende{Chaîne de dilution et de dosage du vinaigre : dilution au 1/10 en fiole jaugée, puis dosage d'une prise d'essai de \SI{10,0}{mL} par la soude titrée.}
\end{center}
\end{popfigure}

\courssub{Exploitation : de l'équivalence au degré d'acidité}

L'exploitation se fait en cascade, en remontant la chaîne de dilution. À l'équivalence, les réactifs ont été mélangés dans les proportions stœchiométriques :
\[ n(\ce{CH3COOH})_{\text{dosé}} = n(\ce{OH-})_{\text{versé}} \quad\Longrightarrow\quad C_A \cdot V_A = C_B \cdot V_{BE} \]

\begin{enumerate}
\item Concentration molaire de la solution \textbf{diluée} : $C_A = \dfrac{C_B \cdot V_{BE}}{V_A}$.
\item Concentration molaire du \textbf{vinaigre commercial} : $C_0 = F \times C_A = 10 \times C_A$.
\item \textbf{Concentration massique} : $C_m = C_0 \times M(\ce{CH3COOH})$ avec $M = \SI{60,0}{g.mol^{-1}}$.
\item \textbf{Degré d'acidité} : masse d'acide dans \SI{100}{g} de vinaigre $\approx$ masse dans \SI{100}{mL}, soit $d = \dfrac{C_m}{10}$ (avec $C_m$ en \si{g.L^{-1}}).
\end{enumerate}

\begin{popfigure}
\begin{center}
\begin{tikzpicture}[node distance=0.55cm, every node/.style={font=\small},
boite/.style={draw=popblue, thick, fill=popblueL, rounded corners=3pt, align=center, minimum height=0.9cm, inner sep=4pt},
fleche/.style={->, thick, poporange}]
\node[boite] (vbe) {$V_{BE}$ lu à la burette\\ (ou sur la courbe)};
\node[boite, right=of vbe] (ca) {$C_A = \dfrac{C_B V_{BE}}{V_A}$\\ (solution diluée)};
\node[boite, right=of ca] (c0) {$C_0 = 10 \times C_A$\\ (vinaigre commercial)};
\node[boite, below=0.7cm of c0] (cm) {$C_m = C_0 \times M$\\ $M = \SI{60,0}{g.mol^{-1}}$};
\node[boite, left=of cm, fill=popgoldL, draw=popgold] (deg) {Degré $d = \dfrac{C_m}{10}$\\ comparaison à l'étiquette};
\draw[fleche] (vbe) -- (ca);
\draw[fleche] (ca) -- (c0);
\draw[fleche] (c0) -- (cm);
\draw[fleche] (cm) -- (deg);
\end{tikzpicture}
\legende{Organigramme de la démarche d'exploitation : du volume équivalent au degré d'acidité, en remontant la dilution.}
\end{center}
\end{popfigure}

\begin{exoresolu}[Degré d'acidité d'un vinaigre du marché]
\textbf{Énoncé.} Un vinaigre commercial porte la mention « 6° ». On le dilue 10 fois. On dose $V_A = \SI{10,0}{mL}$ de solution diluée par de la soude à $C_B = \SI{0,050}{mol.L^{-1}}$. Le virage de la phénolphtaléine (confirmé par le suivi pH-métrique) a lieu à $V_{BE} = \SI{13,4}{mL}$. On donne $M(\ce{CH3COOH}) = \SI{60,0}{g.mol^{-1}}$. Déterminer le degré d'acidité et vérifier l'étiquette.
\tcblower
\textbf{Solution détaillée.}

Équation-bilan de la réaction de dosage (totale) :
\[ \ce{CH3COOH + OH- -> CH3COO- + H2O} \]

\textbf{Étape 1} — Concentration de la solution diluée :
\[ C_A = \frac{C_B \cdot V_{BE}}{V_A} = \frac{0,050 \times 13,4}{10,0} = \SI{0,0670}{mol.L^{-1}} \]

\textbf{Étape 2} — Concentration du vinaigre commercial (on multiplie par le facteur de dilution $F=10$) :
\[ C_0 = 10 \times C_A = \SI{0,670}{mol.L^{-1}} \]

\textbf{Étape 3} — Concentration massique :
\[ C_m = C_0 \times M = 0,670 \times 60,0 = \SI{40,2}{g.L^{-1}} \]

\textbf{Étape 4} — Degré d'acidité :
Le degré $d$ est la masse d'acide (en g) pour \SI{100}{g} de vinaigre. Avec $\rho \approx \SI{1,00}{g.mL^{-1}}$, \SI{100}{g} $\approx$ \SI{100}{mL} = \SI{0,100}{L}.
Masse d'acide dans \SI{100}{mL} $= C_m \times 0,100 = 40,2 \times 0,100 = \SI{4,02}{g}$.
Donc $d \approx 4,0°$.

\textbf{Analyse critique :} Le résultat trouvé ($d \approx 4,0°$) est significativement inférieur à la mention « 6° » de l'étiquette. Cet écart (environ 33 \%) dépasse largement les incertitudes expérimentales usuelles (5 à 10 \%). Soit l'échantillon est non conforme (fraude, altération), soit les données numériques de l'énoncé ($V_{BE} = \SI{13,4}{mL}$) correspondent à un vinaigre à 4°. Pour un vinaigre conforme à 6° ($C_0 \approx \SI{1,00}{mol.L^{-1}}$), on aurait attendu $V_{BE} \approx \SI{20,0}{mL}$ dans ces conditions.

\textbf{Résultat retenu :} Le vinaigre dosé a un degré d'acidité d'environ \textbf{4,0°}. L'étiquette « 6° » n'est \textbf{pas vérifiée} pour cet échantillon.
\end{exoresolu}

\begin{attention}[Erreurs fréquentes au Bac]
\begin{itemize}
\item \textbf{Oublier le facteur de dilution :} Conclure que le vinaigre vaut $C_A = \SI{0,067}{mol.L^{-1}}$ et déclarer l'étiquette fausse, alors qu'on a simplement oublié de multiplier par $F = 10$.
\item \textbf{Confondre les unités :} Mélanger concentration molaire (\si{mol.L^{-1}}), concentration massique (\si{g.L^{-1}}) et degré (g pour \SI{100}{g}). Réflexe : à chaque étape, écris l'unité et demande-toi « De quelle solution parle-t-on : diluée ou commerciale ? ».
\item \textbf{Mauvaise conversion degré / concentration massique :} Retenir que $d = C_m / 10$ (si $C_m$ en \si{g.L^{-1}}) car \SI{100}{mL} = \SI{0,1}{L}.
\end{itemize}
\end{attention}

\begin{exemplebox}[Exemple chiffré de référence (vinaigre conforme 6°)]
Vinaigre dilué 10 fois ; prise d'essai $V_A = \SI{10,0}{mL}$ ; soude à $C_B = \SI{0,050}{mol.L^{-1}}$ ; $V_{BE} = \SI{20,0}{mL}$.
\begin{align*}
C_A &= \frac{0,050 \times 20,0}{10,0} = \SI{0,100}{mol.L^{-1}} \\
C_0 &= 10 \times 0,100 = \SI{1,00}{mol.L^{-1}} \\
C_m &= 1,00 \times 60,0 = \SI{60,0}{g.L^{-1}} \\
d &= \frac{60,0}{10} = 6,0°
\end{align*}
Compte tenu des incertitudes expérimentales (lecture burette $\pm \SI{0,1}{mL}$, virage indicateur), un écart relatif de l'ordre de 10 \% par rapport à la mention de l'étiquette reste acceptable : un vinaigre annoncé à 6° dont le dosage donne entre 5,4° et 6,6° est jugé \textbf{conforme}.
\end{exemplebox}

\begin{popfigure}
\begin{center}
\begin{tikzpicture}
\begin{axis}[
width=0.9\linewidth, height=6.5cm,
xlabel={$V_B$ (mL)}, ylabel={pH},
xmin=0, xmax=25, ymin=2, ymax=12,
grid=major, grid style={popline!40},
tick label style={font=\small}, label style={font=\small},
clip=false]
% Courbe sigmoïde approximative pour acide faible / base forte
% pH initial ~2.9, pH équivalence ~8.7, pKa=4.76
\addplot[popcurve, domain=0:24.5, samples=400, thick] 
{2.9 + (8.7-2.9)/(1+exp(-(x-13.4)*1.2)) + 0.02*(x-13.4)};
% Point E
\addplot[mark=*, mark size=2.5pt, poporange] coordinates {(13.4,8.7)};
\node[font=\small, poporange, anchor=west] at (axis cs:14,8.2) {$E$ ($V_{BE} = \SI{13,4}{mL}$, $pH_E \approx 8,7$)};
\draw[dashed, poporange, thin] (axis cs:13.4,2) -- (axis cs:13.4,8.7);
\draw[dashed, poporange, thin] (axis cs:0,8.7) -- (axis cs:13.4,8.7);
% Point H demi-equivalence
\addplot[mark=*, mark size=2.5pt, popgreen] coordinates {(6.7,4.76)};
\node[font=\small, popgreen!60!black, anchor=south east] at (axis cs:6.5,4.9) {$H$ : $pH = pK_a = 4,76$};
\draw[dashed, popgreen, thin] (axis cs:6.7,2) -- (axis cs:6.7,4.76);
\draw[dashed, popgreen, thin] (axis cs:0,4.76) -- (axis cs:6.7,4.76);
% Zone virage phénolphtaléine
\fill[poppink!20] (axis cs:0,8.2) rectangle (axis cs:25,10);
\node[font=\footnotesize, poppink!60!black, rotate=90, anchor=south] at (axis cs:25.2,9) {Zone virage phénolphtaléine};
\end{axis}
\end{tikzpicture}
\legende{Courbe de dosage du vinaigre dilué (\SI{10,0}{mL}, $C_A \approx \SI{0,067}{mol.L^{-1}}$) par la soude \SI{0,050}{mol.L^{-1}} : point d'équivalence $E$ ($pH_E > 7$) et point de demi-équivalence $H$ ($pH = pK_a$). La zone de virage de la phénolphtaléine (8,2--10) contient bien $pH_E$.}
\end{center}
\end{popfigure}

\courssub{Incertitudes et sources d'erreur}

Un résultat de dosage n'a de valeur que si l'on discute sa précision. Les principales sources d'erreur :

\begin{itemize}
\item \textbf{Lecture de la burette} : une graduation au dixième de millilitre donne une incertitude d'environ $\pm \SI{0,05}{mL}$ à $\pm \SI{0,1}{mL}$ sur $V_{BE}$, soit près de 1 \% d'erreur relative pour $V_{BE} \approx \SI{13}{mL}$.
\item \textbf{Repérage du virage} : en dosage colorimétrique, l'œil retient souvent une goutte de trop ; le virage doit être \emph{franc et persistant} (au moins 30 s). Un dépassement de \SI{0,2}{mL} fausse le résultat de 1,5 \%.
\item \textbf{Précision de la dilution} : pipette et fiole jaugées de classe A garantissent $\pm \SI{0,05}{mL}$ (pour \SI{10}{mL} et \SI{100}{mL}) ; un trait de jauge mal ajusté se répercute directement sur $C_0$ via le facteur $F$.
\item \textbf{Choix de l'indicateur} : la phénolphtaléine (zone de virage 8,2 -- 10) convient parfaitement car $pH_E \approx 8,7$ ; l'hélianthine (3,1 -- 4,4) donnerait un virage bien avant l'équivalence et un degré d'acidité largement sous-estimé.
\item \textbf{Carbonatation de la soude} : si la solution de \ce{NaOH} n'est pas fraîche, elle a réagi avec le \ce{CO2} de l'air, sa concentration réelle est inférieure à la valeur nominale, ce qui surestime $C_0$.
\end{itemize}

Au total, une incertitude relative de 5 à 10 \% est réaliste en TP de lycée : c'est pourquoi on compare le résultat à l'étiquette avec un esprit critique, et non au centième près.

\begin{savaistu}[Le contrôle qualité au Cameroun]
Ces mêmes dosages acido-basiques sont pratiqués quotidiennement au \textbf{LANACOME} (Laboratoire National de Contrôle de Qualité des Médicaments et d'Experts) à Yaoundé, ainsi que dans les laboratoires de la SONARA (raffinage), des brasseries (brassage, contrôle acidité bière/bissap) ou des savonneries artisanales (huile de palme, lessive) : vérification de la concentration des solutions injectables, acidité des vinaigres et jus de fruits, contrôle de la soude dans les lessives. Le degré d'acidité que tu viens de calculer est exactement la grandeur qu'un inspecteur des normes (ANOR) vérifie avant d'autoriser la mise sur le marché d'un vinaigre !
\end{savaistu}

\courssub{Bilan de la leçon : la démarche complète d'un dosage}

\begin{aretenir}[La démarche d'un dosage acido-basique en 6 temps]
\begin{enumerate}
\item \textbf{Protocole} : diluer si nécessaire (facteur $F$), prélever un volume exact $V_A$, choisir le suivi (pH-métrique ou indicateur adapté au $pH_E$).
\item \textbf{Réaction} : écrire l'équation-bilan de la réaction de dosage (totale, rapide) : \ce{H3O+ + OH- -> 2H2O} ou \ce{CH3COOH + OH- -> CH3COO- + H2O}.
\item \textbf{Courbe ou virage} : tracer $pH = f(V_B)$ ou observer le changement de teinte persistant.
\item \textbf{Équivalence} : repérer $V_{BE}$ (méthode des tangentes parallèles, dérivée ou virage) ; à l'équivalence, $C_A V_A = C_B V_{BE}$.
\item \textbf{Calcul} : remonter à la solution commerciale ($\times F$), convertir si besoin en concentration massique ou en degré.
\item \textbf{Conclusion critique} : comparer à l'étiquette en tenant compte des incertitudes ; un écart inférieur à 10 \% est généralement acceptable.
\end{enumerate}
\end{aretenir}

\begin{exercice}[À toi de jouer]
Une solution commerciale de déboucheur contient de l'hydroxyde de sodium. On la dilue 50 fois, puis on dose $V_B = \SI{20,0}{mL}$ de solution diluée par de l'acide chlorhydrique à $C_A = \SI{0,10}{mol.L^{-1}}$. L'équivalence est atteinte pour $V_{AE} = \SI{16,0}{mL}$.
\begin{enumerate}
\item Écrire l'équation-bilan de la réaction de dosage.
\item Calculer la concentration de la solution diluée, puis celle du déboucheur commercial.
\item Le fabricant annonce \SI{4,0}{mol.L^{-1}}. L'étiquette est-elle conforme à 10 \% près ?
\end{enumerate}
\end{exercice}

\begin{corrige}[Exercice]
1. \ce{H3O+ + OH- -> 2H2O} (réaction totale, rapide, exothermique).
2. À l'équivalence : $n(\ce{H3O+}) = n(\ce{OH-}) \Rightarrow C_A V_{AE} = C_B V_B$.
$C_B = \dfrac{C_A V_{AE}}{V_B} = \dfrac{0,10 \times 16,0}{20,0} = \SI{0,080}{mol.L^{-1}}$.
Solution commerciale : $C_0 = 50 \times 0,080 = \SI{4,0}{mol.L^{-1}}$.
3. Valeur trouvée = valeur annoncée = \SI{4,0}{mol.L^{-1}}. Écart relatif nul : l'étiquette est parfaitement conforme.
\end{corrige}

\newpage

\coursec{Travaux pratiques}

\courssub{Objectifs du TP}

À la fin de ce TP, tu seras capable de :

\begin{itemize}
\item réaliser une \cle{dilution} précise d'une solution commerciale (vinaigre) à l'aide d'une pipette jaugée et d'une fiole jaugée ;
\item mettre en \oe uvre un \cle{dosage pH-métrique} d'un acide faible (l'acide éthanoïque \ce{CH3COOH}) par une base forte (la soude \ce{NaOH}) ;
\item tracer la courbe $\mathrm{pH} = f(V_B)$ et déterminer le volume équivalent $V_{BE}$ par la \cle{méthode des tangentes parallèles} ;
\item vérifier expérimentalement qu'à la \cle{demi-équivalence}, $\mathrm{pH} = \mathrm{p}K_a$ du couple \ce{CH3COOH}/\ce{CH3COO-} ;
\item réaliser un \cle{dosage colorimétrique} avec la phénolphtaléine et comparer les deux déterminations de $V_{BE}$ ;
\item calculer le \cle{degré d'acidité} du vinaigre et conclure sur la conformité de l'étiquette.
\end{itemize}

\begin{savaistu}[Pourquoi doser un vinaigre ?]
Le \cle{degré d'acidité} d'un vinaigre est la masse d'acide éthanoïque, en grammes, contenue dans \SI{100}{mL} de vinaigre. Un vinaigre « à \SI{6}{\percent} » (ou « 6 degrés ») contient donc \SI{6}{g} d'acide éthanoïque pour \SI{100}{mL}. Au marché de Mokolo ou de Mvog-Mbi, on trouve aussi bien des vinaigres industriels étiquetés que des vinaigres artisanaux issus de la fermentation du vin de palme ou du bili-bili : leur acidité réelle est rarement contrôlée. Ce TP te donne les moyens de vérifier une étiquette, comme le ferait un laboratoire de contrôle qualité.
\end{savaistu}

\courssub{Matériel et produits}

\begin{center}
\begin{tabularx}{\linewidth}{|l|X|}
\hline
\rowcolor{popblue!40} \textbf{Matériel} & \textbf{Produits} \\
\hline
Burette graduée de \SI{25}{mL} & Vinaigre commercial étiqueté « 6\textdegree » \\
\rowcolor{popblueL} Pipette jaugée de \SI{10,0}{mL} (+ poire aspirante) & Solution de soude titrée à $C_B = \SI{0,050}{mol.L^{-1}}$ \\
Fiole jaugée de \SI{100,0}{mL} & Phénolphtaléine (solution alcoolique) \\
\rowcolor{popblueL} Bécher de \SI{100}{mL} (2) & Eau distillée (pissette) \\
pH-mètre étalonné + sonde (ou papier pH) & \\
\rowcolor{popblueL} Agitateur magnétique + barreau (ou agitateur en verre) & \\
Support, pince, entonnoir & \\
\hline
\end{tabularx}
\end{center}

\begin{attention}[Solutions de repli avec du matériel local]
\begin{itemize}
\item \textbf{Pas de pH-mètre ?} Réalise uniquement le dosage colorimétrique à la phénolphtaléine (3 essais concordants). Le papier pH permet un suivi grossier du saut de pH.
\item \textbf{Pas de soude titrée ?} Elle peut être préparée à partir de pastilles de soude du commerce (produit d'entretien), mais il faudra d'abord l'étalonner avec une solution d'acide chlorhydrique titrée, car la soude se carbonate à l'air.
\item \textbf{Pas d'agitateur magnétique ?} Agite doucement le bécher à la main ou avec un agitateur en verre après chaque ajout de soude.
\item \textbf{Pas de vinaigre étiqueté ?} Utilise du vinaigre artisanal de vin de palme : le TP devient alors une \emph{mesure} du degré d'acidité, et non plus une vérification.
\end{itemize}
\end{attention}

\courssub{Sécurité}

\begin{attention}[Consignes de sécurité]
\begin{itemize}
\item La \textbf{soude} est \cle{corrosive} (pictogramme : deux gouttes attaquant une main et une surface métallique). Porte obligatoirement une \textbf{blouse}, des \textbf{lunettes de protection} et des \textbf{gants}. En cas de contact avec la peau, rince abondamment à l'eau froide pendant plusieurs minutes et préviens l'enseignant.
\item La \textbf{phénolphtaléine} est en solution dans l'éthanol : produit \textbf{inflammable} (pictogramme : flamme). Aucune flamme nue sur la paillasse.
\item Ne jamais pipetter à la bouche : utilise toujours la \textbf{poire aspirante} (propipette).
\item La dilution du vinaigre ne présente pas de danger, mais évite de respirer directement les vapeurs au-dessus de la fiole.
\item Les solutions usées (vinaigre neutralisé) peuvent être jetées à l'évier avec beaucoup d'eau courante.
\end{itemize}
\end{attention}

\courssub{Schéma du dispositif}

\begin{popfigure}
\begin{center}
\begin{tikzpicture}[scale=0.95, every node/.style={font=\small}]
% Support
\draw[popdark, line width=1.2pt] (-3.2,0) -- (-3.2,6.4);
\draw[popdark, line width=1.6pt] (-4.0,0) -- (-2.4,0);
% Pince
\draw[popdark, line width=1pt] (-3.2,5.2) -- (-1.1,5.2);
\draw[popdark, line width=1pt] (-1.1,5.35) -- (-1.1,5.05);
% Burette
\draw[popblue, line width=1pt] (-0.85,6.3) -- (-0.85,1.9);
\draw[popblue, line width=1pt] (-0.35,6.3) -- (-0.35,1.9);
\draw[popblue!50, fill=popblue!20] (-0.83,5.9) rectangle (-0.37,2.0);
\draw[popblue, line width=1pt] (-0.85,1.9) -- (-0.7,1.55) -- (-0.5,1.55) -- (-0.35,1.9);
% Robinet
\draw[popdark, fill=popdark] (-0.95,1.45) rectangle (-0.25,1.65);
\draw[popdark, line width=1pt] (-0.6,1.45) -- (-0.6,1.15);
% Graduations
\foreach \y in {2.4,2.9,3.4,3.9,4.4,4.9,5.4} \draw[popblue!70] (-0.85,\y) -- (-0.65,\y);
\node[popblue, anchor=west] at (-0.2,5.6) {soude $C_B = \SI{0,050}{mol.L^{-1}}$};
% Goutte
\draw[popblue, fill=popblue!60] (-0.6,0.85) circle (0.05);
% Bécher
\draw[popdark, line width=1.2pt] (-1.6,-0.05) -- (-1.6,-1.9) -- (0.4,-1.9) -- (0.4,-0.05);
\draw[poporange!60, fill=poporange!15] (-1.55,-1.0) rectangle (0.35,-1.85);
\node[poporange!80!black, anchor=west] at (0.55,-1.5) {vinaigre dilué $V_A = \SI{10,0}{mL}$};
% Barreau magnétique
\draw[popdark, fill=popmuted!40] (-0.9,-1.8) rectangle (-0.3,-1.68);
% Sonde pH
\draw[popgreen!70!black, line width=1.4pt] (1.6,1.6) -- (0.1,-0.4);
\draw[popgreen!70!black, fill=popgreen!30] (0.1,-0.4) -- (0.0,-1.2) -- (-0.25,-1.2) -- (-0.14,-0.35) -- cycle;
\node[popgreen!60!black, anchor=west] at (1.55,1.75) {sonde du pH-mètre};
% pH-mètre
\draw[popdark, fill=popcream, line width=1pt] (2.6,0.6) rectangle (4.6,2.0);
\draw[popgreen!70!black, fill=popgreen!20] (2.85,1.25) rectangle (4.35,1.75);
\node[popgreen!60!black] at (3.6,1.5) {pH = \num{3,4}};
\node[popmuted, font=\scriptsize] at (3.6,0.9) {pH-mètre};
% Agitateur
\draw[popmuted, fill=popmuted!20] (-2.0,-2.0) rectangle (0.8,-2.5);
\node[popmuted, font=\scriptsize] at (-0.6,-2.25) {agitateur magnétique};
\end{tikzpicture}
\end{center}
\legende{Dispositif de dosage pH-métrique du vinaigre dilué par la soude.}
\end{popfigure}

\courssub{Protocole}

\textbf{Étape 1 — Dilution du vinaigre (facteur 10).}\par
\begin{enumerate}
\item À l'aide de la pipette jaugée munie de la poire aspirante, prélever $V_0 = \SI{10,0}{mL}$ de vinaigre commercial.
\item Les introduire dans la fiole jaugée de \SI{100,0}{mL} contenant déjà un peu d'eau distillée.
\item Compléter avec de l'eau distillée jusqu'au trait de jauge, en ajustant le bas du ménisque. Boucher et homogénéiser par retournements (au moins 10 fois). On obtient la solution $S$, dix fois moins concentrée que le vinaigre.
\end{enumerate}

\textbf{Étape 2 — Dosage pH-métrique.}\par
\begin{enumerate}
\setcounter{enumi}{3}
\item Rincer et remplir la burette avec la soude à $C_B = \SI{0,050}{mol.L^{-1}}$, chasser la bulle d'air, ajuster le zéro.
\item Prélever $V_A = \SI{10,0}{mL}$ de la solution $S$ dans le bécher ; ajouter environ \SI{20}{mL} d'eau distillée pour immerger correctement la sonde (l'ajout d'eau ne modifie pas $V_{BE}$).
\item Étalonner le pH-mètre (solutions tampons pH 4 et pH 7), rincer la sonde à l'eau distillée, la plonger dans le bécher. Lancer l'agitation douce.
\item Relever le pH après chaque ajout de soude : tous les \SI{1}{mL} de $V_B = 0$ à \SI{15}{mL}, puis tous les \SI{0,2}{mL} dans la zone du saut de pH (environ 19 à \SI{22}{mL}), puis de nouveau tous les \SI{1}{mL} jusqu'à \SI{25}{mL}.
\end{enumerate}

\textbf{Étape 3 — Dosage colorimétrique.}\par
\begin{enumerate}
\setcounter{enumi}{7}
\item Recommencer le dosage sur une nouvelle prise de \SI{10,0}{mL} de $S$, en ajoutant \textbf{3 gouttes de phénolphtaléine}.
\item Verser la soude jusqu'au virage persistant (rose pâle stable au moins 30 s). Noter le volume $V'_{BE}$. Refaire deux essais concordants.
\end{enumerate}

\courssub{Observations et résultats attendus}

Exemple de tableau de mesures obtenu avec un vinaigre étiqueté « 6\textdegree » :

\begin{center}
\small
\begin{tabular}{|c|c|c|c|c|c|c|c|c|c|c|}
\hline
\rowcolor{poporange!50} $V_B$ (mL) & 0 & 2 & 4 & 6 & 8 & 10 & 12 & 14 & 16 & 18 \\
\hline
pH & 2,9 & 3,8 & 4,2 & 4,4 & 4,6 & 4,8 & 5,0 & 5,2 & 5,4 & 5,7 \\
\hline
\rowcolor{poporange!50} $V_B$ (mL) & 19 & 19,6 & 19,8 & 20,0 & 20,2 & 20,4 & 21 & 22 & 24 & 25 \\
\hline
pH & 6,0 & 6,6 & 7,3 & 8,7 & 10,4 & 10,8 & 11,1 & 11,3 & 11,5 & 11,6 \\
\hline
\end{tabular}
\end{center}

Observations attendues :

\begin{itemize}
\item le pH initial est voisin de 2,9 (acide faible : pH $> -\log C_A$) ;
\item la courbe $\mathrm{pH} = f(V_B)$ est croissante, avec un \emph{saut de pH} modéré (environ de 6,5 à 10,5) autour de $V_{BE} \approx \SI{20,0}{mL}$ ;
\item à la demi-équivalence ($V_B = V_{BE}/2 \approx \SI{10,0}{mL}$), on lit $\mathrm{pH} \approx 4{,}8 = \mathrm{p}K_a$ du couple \ce{CH3COOH}/\ce{CH3COO-} ;
\item le pH à l'équivalence est \textbf{supérieur à 7} (environ 8,7) : la solution contient alors la base faible \ce{CH3COO-} ;
\item au dosage colorimétrique, la solution vire au rose pâle pour $V'_{BE} \approx \SI{20,0}{mL}$, valeur très proche de celle du dosage pH-métrique.
\end{itemize}

\begin{propriete}[Caractéristiques du dosage acide faible / base forte]
Pour le dosage d'un acide faible (AH) par une base forte (HO$^-$) :
\begin{itemize}
\item L'équation-bilan est $\ce{AH + OH- -> A- + H2O}$ (totale, rapide, unique).
\item La courbe $\mathrm{pH} = f(V_B)$ présente un saut de pH moins marqué que pour fort/fort.
\item \textbf{À la demi-équivalence} ($V_B = V_{BE}/2$) : $\mathrm{pH} = \mathrm{p}K_a$ du couple $\ce{AH}/\ce{A-}$.
\item \textbf{À l'équivalence} : le pH est \textbf{supérieur à 7} (solution de la base conjuguée $\ce{A-}$, hydrolyse basique).
\end{itemize}
\end{propriete}

\begin{methode}[Détermination graphique de $V_{BE}$ par la méthode des tangentes parallèles]
\begin{enumerate}
\item Tracer la courbe $\mathrm{pH} = f(V_B)$ en point par point (abscisse : $V_B$, ordonnée : pH).
\item Tracer la tangente à la branche montante \textbf{avant} le saut (zone quasi-linéaire).
\item Tracer la tangente à la branche montante \textbf{après} le saut (zone quasi-linéaire).
\item Tracer la \textbf{bissectrice} de l'angle formé par ces deux tangentes (ou tracer une droite parallèle aux deux tangentes et passant par le point d'inflexion visible).
\item L'abscisse du point d'intersection de cette bissectrice (ou de cette parallèle) avec la courbe est le volume équivalent $V_{BE}$.
\end{enumerate}
\end{methode}

\courssub{Exploitation}

\begin{exoresolu}[Exploitation complète du TP]
\textbf{Questions.}
\begin{enumerate}
\item Écrire l'équation-bilan de la réaction de dosage. Justifier qu'elle peut servir de réaction de dosage.
\item Définir l'équivalence et établir la relation entre $C_A$, $V_A$, $C_B$ et $V_{BE}$.
\item Calculer la concentration $C_A$ de la solution diluée $S$, puis la concentration $C_0$ du vinaigre commercial.
\item En déduire le degré d'acidité du vinaigre. L'étiquette « 6\textdegree » est-elle conforme ? (On donne $M(\ce{CH3COOH}) = \SI{60}{g.mol^{-1}}$.)
\item Pourquoi la phénolphtaléine (zone de virage 8,2 -- 10,0) est-elle un bon choix ici ? L'hélianthine (3,1 -- 4,4) conviendrait-elle ?
\end{enumerate}
\tcblower
\textbf{Solution détaillée.}
\begin{enumerate}
\item La réaction entre l'acide éthanoïque et les ions hydroxyde s'écrit :
\[ \ce{CH3COOH + OH- -> CH3COO- + H2O} \]
C'est une réaction acide/base entre l'acide \ce{CH3COOH} (couple \ce{CH3COOH}/\ce{CH3COO-}, $\mathrm{p}K_a = 4,8$) et la base \ce{OH-} (couple \ce{H2O}/\ce{OH-}, $\mathrm{p}K_a = 14$ via $K_e = 10^{-14}$).
Sa constante d'équilibre vaut $K = \frac{K_a(\ce{CH3COOH})}{K_e} = \frac{10^{-4,8}}{10^{-14}} = 10^{9,2} \gg 10^4$ : la réaction est \textbf{totale}. Elle est de plus \textbf{rapide} et \textbf{unique} (pas de réaction parasite) : elle convient parfaitement pour un dosage.
\item À l'\cle{équivalence}, les réactifs ont été mélangés dans les proportions stœchiométriques : la quantité de matière d'acide éthanoïque initialement présente est égale à la quantité de matière d'ions hydroxyde versée.
\[ n(\ce{CH3COOH})_{\text{initial}} = n(\ce{OH-})_{\text{versé}} \quad \Rightarrow \quad C_A \cdot V_A = C_B \cdot V_{BE} \]
\item Concentration de la solution diluée $S$ :
\[ C_A = \frac{C_B \cdot V_{BE}}{V_A} = \frac{\SI{0,050}{mol.L^{-1}} \times \SI{20,0}{mL}}{\SI{10,0}{mL}} = \SI{0,10}{mol.L^{-1}} \]
Le vinaigre a été dilué 10 fois ($V_0 = \SI{10,0}{mL}$ dans \SI{100,0}{mL}) : $C_0 = 10 \times C_A = \SI{1,0}{mol.L^{-1}}$.
\item Concentration massique du vinaigre commercial :
\[ \gamma = C_0 \times M = \SI{1,0}{mol.L^{-1}} \times \SI{60}{g.mol^{-1}} = \SI{60}{g.L^{-1}} \]
Le degré d'acidité est la masse d'acide dans \SI{100}{mL} (soit \SI{0,100}{L}) :
\[ \text{Degré} = \gamma \times \SI{0,100}{L} = \SI{60}{g.L^{-1}} \times \SI{0,100}{L} = \SI{6,0}{g} \]
Le vinaigre est donc à \textbf{6,0\textdegree}. \textbf{Conclusion :} l'écart avec l'étiquette est nul (ou inférieur à l'incertitude expérimentale, de l'ordre de 5 \%) : l'étiquette est \textbf{conforme}.
\item Le pH à l'équivalence vaut environ 8,7 : il appartient à la zone de virage de la phénolphtaléine (8,2 -- 10,0), d'où un virage net au bon volume. L'hélianthine virerait vers pH 3--4, \emph{bien avant} l'équivalence (zone tampon) : le volume repéré serait trop faible et le résultat faux. L'hélianthine est donc à proscrire pour ce dosage.
\end{enumerate}
\end{exoresolu}

\begin{attention}[Erreurs fréquentes à éviter]
\begin{itemize}
\item Oublier le \textbf{facteur de dilution} : la concentration calculée à partir du dosage est celle de la solution $S$, pas celle du vinaigre !
\item Confondre \textbf{degré d'acidité} (g pour \SI{100}{mL}) et concentration massique (g/L) : il y a un facteur 10 entre les deux.
\item Verser la soude trop vite près du saut de pH : on « rate » l'équivalence. Réduis les ajouts à \SI{0,2}{mL} dès que le pH varie sensiblement.
\item Ajouter trop de phénolphtaléine : 2 à 3 gouttes suffisent ; un excès trouble le virage et consomme de la soude.
\item Croire que l'eau ajoutée pour immerger la sonde fausse le dosage : elle dilue, mais ne change pas la \emph{quantité de matière} d'acide, donc pas $V_{BE}$.
\item Écrire $\ce{H3O+ + OH- -> 2H2O}$ pour ce dosage : c'est l'équation du dosage fort/fort. Ici, l'acide est faible, l'équation-bilan fait intervenir la forme moléculaire \ce{CH3COOH}.
\end{itemize}
\end{attention}

\courssub{Conclusion}

Ce TP t'a permis de vérifier expérimentalement l'étiquette d'un produit de consommation courante, exactement comme le fait un technicien de laboratoire de contrôle qualité. Le dosage pH-métrique a fourni une détermination précise du volume équivalent ($V_{BE} \approx \SI{20,0}{mL}$) par la méthode des tangentes, et a confirmé deux propriétés essentielles du dosage d'un acide faible par une base forte : $\mathrm{pH} = \mathrm{p}K_a$ à la demi-équivalence et $\mathrm{pH} > 7$ à l'équivalence. Le dosage colorimétrique à la phénolphtaléine, plus rapide et moins coûteux, donne un résultat concordant : c'est la méthode de choix pour un contrôle de routine, à condition de bien choisir l'indicateur (zone de virage contenant le pH à l'équivalence). Le degré d'acidité trouvé (6,0\textdegree) valide l'étiquette du vinaigre étudié. La même démarche s'applique au dosage des détergents, des jus de fruits (bissap, foléré) ou des eaux minérales : tu disposes désormais d'une compétence expérimentale complète, exigible au Baccalauréat et utile dans la vie courante.

\newpage

\coursec{Méthodes et exercices résolus}

\coursec{Leçon 09 : Réactions acide/base et dosages}

\begin{savaistu}[Situation déclenchante : le vinaigre de palme et le contrôle qualité]
Au marché de Mokolo à Yaoundé, comme à Nkoulouloun à Douala, les vendeuses de vinaigre de palme (issu de la fermentation du vin de palme) ou de bili-bili vantent l'acidité de leur produit : « \emph{C'est du costaud, ça dégraisse bien !} ». Mais comment vérifier scientifiquement cette acidité ? L'étiquette d'un déboucheur (acide chlorhydrique) annonce 23\,\% massique ; est-ce exact ? En chimie analytique, le \textbf{dosage acido-basique} permet de répondre à ces questions. Cette leçon te donne les outils rigoureux (pH-métrie, colorimétrie, choix d'indicateur) pour titrer une espèce chimique, vérifier une étiquette commerciale et déterminer le \textbf{degré d'acidité} d'un vinaigre, conformément au programme officiel de Terminale C, D, E.
\end{savaistu}

\courssub{Réaction acide fort -- base forte : étude et équivalence}

\begin{definition}[Réaction de neutralisation acide fort -- base forte]
Lorsqu'un acide fort (totalement dissocié en \ce{H3O+}) rencontre une base forte (totalement dissociée en \ce{OH-}), la réaction effective ne met en jeu que les ions \ce{H3O+} et \ce{OH-}. Les contre-ions (\ce{Cl-}, \ce{Na+}, \ce{NO3-}, \ce{K+}...) sont des \textbf{ions spectateurs}.
\[ \ce{H3O+ + OH- -> 2H2O} \qquad (K = 1/K_e = 10^{14} \text{ à } \SI{25}{\celsius})
\]
\end{definition}

\begin{propriete}[Caractéristiques de la réaction]
Cette réaction est :
\begin{itemize}
\item \textbf{Totale} : la constante d'équilibre $K = 10^{14}$ est immense ; les réactifs sont consommés jusqu'à épuisement de l'un d'eux (réactif limitant).
\item \textbf{Rapide} : cinétique quasi instantanée à l'échelle macroscopique (vitesse limitée par la diffusion/mélange).
\item \textbf{Exothermique} : $\Delta_r H^\circ \approx \SI{-57,1}{kJ.mol^{-1}}$ (chaleur libérée, température qui monte).
\end{itemize}
\end{propriete}

\begin{aretenir}[Équivalence acido-basique (monoacide / monobase)]
À l'équivalence, les réactifs ont été introduits dans les proportions stoechiométriques exactes. Il n'y a plus ni acide titré, ni base titrante en excès.
\[ n(\ce{H3O+})_{\text{initial}} = n(\ce{OH-})_{\text{versé}} \quad\Longrightarrow\quad \boxed{C_A V_A = C_B V_{BE}} \]
$C_A, V_A$ : concentration et volume de la solution titrée (acide).\\
$C_B, V_{BE}$ : concentration et volume équivalent de la solution titrante (base).
\end{aretenir}

\begin{methode}[Équation de dosage et équivalence : les réflexes]
\begin{enumerate}
\item \textbf{Identifier} la nature des réactifs : acide fort (\ce{H3O+}), acide faible (\ce{HA}), base forte (\ce{OH-}), base faible (\ce{B}).
\item \textbf{Écrire l'équation-bilan} avec les bonnes espèces (ions spectateurs exclus) :
\begin{itemize}
\item acide fort + base forte : \ce{H3O+ + OH- -> 2H2O} ;
\item acide faible + base forte : \ce{HA + OH- -> A- + H2O} ;
\item acide fort + base faible : \ce{H3O+ + B -> BH+ + H2O}.
\end{itemize}
\item \textbf{Préciser les caractéristiques} : réaction totale, rapide, exothermique.
\item \textbf{Définir l'équivalence} : $C_A V_A = C_B V_{BE}$ (stœchiométrie 1/1 pour monoacide/monobase).
\item \textbf{Exploiter} : connaissant trois grandeurs, calculer la quatrième (souvent $C_A$).
\end{enumerate}
\end{methode}

\begin{exoresolu}[Dosage d'une solution d'acide chlorhydrique]
On dose $V_A = \SI{20,0}{mL}$ d'une solution d'acide chlorhydrique de concentration $C_A$ inconnue par une solution d'hydroxyde de sodium de concentration $C_B = \SI{0,100}{mol.L^{-1}}$. L'équivalence est atteinte pour $V_{BE} = \SI{16,4}{mL}$.
\begin{enumerate}
\item Écrire l'équation-bilan de la réaction de dosage et donner ses caractéristiques.
\item Définir l'équivalence et calculer $C_A$.
\end{enumerate}
\tcblower
\textbf{1. Équation-bilan.} L'acide chlorhydrique est un acide fort, totalement ionisé en \ce{H3O+} et \ce{Cl-} ; la soude est une base forte, totalement dissociée en \ce{Na+} et \ce{OH-}. Les ions \ce{Na+} et \ce{Cl-} sont spectateurs. La réaction de dosage s'écrit :
\[ \ce{H3O+ + OH- -> 2H2O} \]
Cette réaction est \textbf{totale} ($K = 10^{14}$), \textbf{rapide} et \textbf{exothermique}.

\textbf{2. Équivalence et calcul.} À l'équivalence, les réactifs ont été introduits dans les proportions stoechiométriques : $n(\ce{H3O+})_{\text{initial}} = n(\ce{OH-})_{\text{versé}}$, soit :
\[ C_A V_A = C_B V_{BE} \quad\Longrightarrow\quad C_A = \frac{C_B \times V_{BE}}{V_A} = \frac{0{,}100 \times 16{,}4}{20{,}0}. \]
\[ \boxed{C_A = \SI{8,20e-2}{mol.L^{-1}}} \]
\textbf{Conclusion :} la solution d'acide chlorhydrique a une concentration de \SI{0,0820}{mol.L^{-1}}. (Vérification : $V_{BE} < V_A$, donc $C_A < C_B$, cohérent.)
\end{exoresolu}

\courssub{Dosage pH-métrique d'un acide fort par une base forte}

\begin{experience}[Dosage pH-métrique : dispositif et protocole]
\textbf{Matériel :} Bécher (solution titrée), burette graduée (solution titrante), pH-mètre étalonné (deux tampons : pH 4,01 et 7,00 ou 10,01), barreau aimanté + agitateur magnétique, thermomètre.
\textbf{Protocole :}
\begin{enumerate}
\item Verser $V_A$ de la solution à doser dans le bécher, immerger l'électrode, agiter.
\item Noter le pH initial ($V_B = 0$).
\item Verser la solution titrante par fractions (ex. \SI{1}{mL} puis \SI{0,5}{mL} puis \SI{0,1}{mL} près du saut), noter $(V_B\,;\,\mathrm{pH})$ après stabilisation à chaque ajout.
\item Continuer bien après l'équivalence (excès de base).
\end{enumerate}
\textbf{Observations :} La courbe $\mathrm{pH} = f(V_B)$ présente un \textbf{saut de pH} quasi vertical autour de $V_{BE}$. À l'équivalence, $\mathrm{pH}_E = 7,0$ (solution neutre, seuls ions spectateurs \ce{Na+}, \ce{Cl-}).
\end{experience}

\begin{popfigure}
\begin{tikzpicture}
\begin{axis}[
    width=\linewidth, height=0.55\linewidth,
    xlabel={Volume de soude versé $V_B$ (mL)}, ylabel={pH},
    xmin=0, xmax=25, ymin=0, ymax=14,
    grid=major, grid style={popmuted}, tick label style={font=\small},
    axis line style={popdark}, tick style={popdark},
]
% Courbe avant équivalence : pH = -log10( (n_A - n_B) / V_total )
% n_A = 1 mmol, n_B = 0.1 * x mmol, V_total = (10 + x) mL
% [H+] = (1 - 0.1x) / (10 + x) mol/L
\addplot[popblue, thick, smooth, domain=0:9.9, samples=200] { -log10( (1 - 0.1*x) / (10 + x) ) };
% Courbe après équivalence : pH = 14 + log10( (n_B - n_A) / V_total )
% [OH-] = (0.1x - 1) / (10 + x) mol/L
\addplot[popblue, thick, smooth, domain=10.1:25, samples=200] { 14 + log10( (0.1*x - 1) / (10 + x) ) };
% Saut de pH à l'équivalence (trait vertical)
\draw[popblue, thick] (axis cs:10,3.5) -- (axis cs:10,10.5);
% Point d'équivalence
\node[circle, fill=popblue, inner sep=2pt, draw=popblue, pin={[pin distance=1.2cm, pin edge={popdark}]90:{$E(V_{BE}=10,0;\mathrm{pH}_E=7,0)$}}] at (axis cs:10,7) {};
% Tangentes aux branches (acide et base)
\draw[poporange, dashed, thick] (axis cs:0,1) -- (axis cs:9,1) node[right, font=\small, poporange] {Tangente 1 (acide)};
\draw[poporange, dashed, thick] (axis cs:11,13) -- (axis cs:20,13) node[right, font=\small, poporange] {Tangente 2 (base)};
% Médiatrice (équidistante des deux tangentes horizontales)
\draw[poppurple, thick] (axis cs:0,7) -- (axis cs:25,7) node[right, font=\small, poppurple] {Médiatrice (équidistante)};
\end{axis}
\end{tikzpicture}
\legende{Courbe de dosage pH-métrique acide fort / base forte ($C_A=C_B=\SI{0,1}{mol.L^{-1}}, V_A=\SI{10}{mL}$). Le saut de pH est centré sur $\mathrm{pH}=7$. La méthode des tangentes parallèles (orange) et leur médiatrice (violet) permettent de déterminer $V_{BE}$ graphiquement.}
\end{popfigure}

\begin{methode}[Tracer la courbe de dosage et repérer l'équivalence (méthode des tangentes)]
\begin{enumerate}
\item \textbf{Reporter les points} $(V_B\,;\,\mathrm{pH})$ sur un graphe soigné : volume en abscisse, pH en ordonnée, échelles adaptées (souvent \SI{1}{cm} pour \SI{2}{mL} et \SI{1}{cm} pour 1 unité pH).
\item \textbf{Relier les points} par une courbe lisse (jamais des segments de droite) : la courbe est croissante (dosage d'un acide par une base).
\item \textbf{Identifier le saut de pH} : zone quasi verticale de la courbe.
\item \textbf{Méthode des tangentes parallèles} :
\begin{itemize}
\item tracer deux tangentes à la courbe, parallèles entre elles, de part et d'autre du saut (une sur la branche acide, une sur la branche basique) ;
\item tracer la droite équidistante de ces deux tangentes (médiatrice) ;
\item son intersection avec la courbe donne le \textbf{point d'équivalence} $E(V_{BE}\,;\,\mathrm{pH}_E)$.
\end{itemize}
\item \textbf{Méthode de la dérivée} (notion) : le point d'équivalence correspond au maximum de la courbe $\dfrac{\mathrm{d}\mathrm{pH}}{\mathrm{d}V_B} = f(V_B)$.
\item \textbf{Lire} $V_{BE}$ et $\mathrm{pH}_E$, puis exploiter $C_A V_A = C_B V_{BE}$.
\end{enumerate}
\end{methode}

\begin{exoresolu}[Exploitation graphique d'un dosage acide fort -- base forte]
Le dosage pH-métrique de $V_A = \SI{10,0}{mL}$ d'acide chlorhydrique par la soude à $C_B = \SI{0,050}{mol.L^{-1}}$ donne la courbe $\mathrm{pH}=f(V_B)$. La méthode des tangentes parallèles fournit $V_{BE} = \SI{12,0}{mL}$ et $\mathrm{pH}_E = 7{,}0$.
\begin{enumerate}
\item Justifier la valeur du pH à l'équivalence.
\item Calculer $C_A$.
\item Quelle est la concentration de la solution obtenue à l'équivalence en ions \ce{Na+} ?
\end{enumerate}
\tcblower
\textbf{1. pH à l'équivalence.} À l'équivalence, la solution ne contient que les ions spectateurs \ce{Na+} et \ce{Cl-} (sans action sur l'eau) et de l'eau : c'est une solution de chlorure de sodium, \textbf{neutre}. D'où $\mathrm{pH}_E = 7{,}0$ à \SI{25}{\celsius}.

\textbf{2. Concentration $C_A$.} À l'équivalence :
\[ C_A = \frac{C_B V_{BE}}{V_A} = \frac{0{,}050 \times 12{,}0}{10{,}0} = \SI{6,0e-2}{mol.L^{-1}}. \]

\textbf{3. Concentration en \ce{Na+} à l'équivalence.} Les ions sodium apportés : $n(\ce{Na+}) = C_B V_{BE} = 0{,}050 \times 12{,}0\times10^{-3} = \SI{6,0e-4}{mol}$. Le volume total est $V_A + V_{BE} = \SI{22,0}{mL}$ :
\[ [\ce{Na+}] = \frac{6{,}0\times10^{-4}}{22{,}0\times10^{-3}} = \SI{2,7e-2}{mol.L^{-1}}. \]
\textbf{Conclusion :} $C_A = \SI{0,060}{mol.L^{-1}}$ ; à l'équivalence la solution est neutre et contient \ce{Na+} et \ce{Cl-} à \SI{0,027}{mol.L^{-1}} chacun.
\end{exoresolu}

\courssub{Dosage d'un acide faible par une base forte}

\begin{definition}[Point de demi-équivalence]
Au cours du dosage d'un acide faible \ce{HA} par une base forte, le \textbf{point de demi-équivalence} correspond au volume $V_B = \dfrac{V_{BE}}{2}$. À ce stade, la moitié de l'acide initial a réagi : $n(\ce{HA}) = n(\ce{A-})$. La solution est un \textbf{mélange tampon} de concentration maximale et $\mathrm{pH} = \mathrm{p}K_a(\ce{HA}/\ce{A-})$.
\end{definition}

\begin{popfigure}
\begin{tikzpicture}
\begin{axis}[
    width=\linewidth, height=0.55\linewidth,
    xlabel={Volume de soude versé $V_B$ (\SI{}{mL})}, ylabel={pH},
    xmin=0, xmax=45, ymin=0, ymax=14,
    grid=major, grid style={popmuted}, tick label style={font=\small},
    axis line style={popdark}, tick style={popdark},
]
% Simulated curve Weak Acid (pKa=4.8, 0.1M) / Strong Base (0.1M), Va=20mL -> Vbe=20mL
% Before eq: Henderson-Hasselbalch
\addplot[popblue, thick, domain=0:19.9, samples=150] { 4.8 + log10( (0.1*x) / (0.1*20 - 0.1*x) ) };
% Vertical jump at eq
\draw[popblue, thick] (axis cs:20,7) -- (axis cs:20,10.5);
% After eq: excess OH-
\addplot[popblue, thick, domain=20.1:45, samples=100] { 14 + log10( (0.1*x - 0.1*20) / (20 + x) * 1000 ) };
% Half-equivalence point
\node[circle, fill=popblue, inner sep=2pt, draw=popdark, pin={[pin distance=1.5cm, pin edge={popdark}]135:{$E/2\,(V_B=\SI{10,0}{\milli\liter};\mathrm{pH}=4,8=\mathrm{p}K_a)$}}] at (axis cs:10,4.8) {};
% Equivalence point
\node[circle, fill=popblue, inner sep=2pt, draw=popdark, pin={[pin distance=1.5cm, pin edge={popdark}]45:{$E\,(V_{BE}=\SI{20,0}{\milli\liter};\mathrm{pH}_E\approx 8,7)$}}] at (axis cs:20,8.7) {};
% Buffer zone bracket (simplified without decorations library)
\draw[popgreen, thick] (axis cs:2,2.5) -- (axis cs:18,2.5) node[midway, below=5pt, popgreen] {Zone tampon};
\draw[popgreen, thick] (axis cs:2,2.5) -- (axis cs:2,2.7);
\draw[popgreen, thick] (axis cs:18,2.5) -- (axis cs:18,2.7);
\end{axis}
\end{tikzpicture}
\legende{Courbe de dosage pH-métrique acide faible (éthanoïque, $\mathrm{p}K_a=4,8$) / base forte. On observe la zone tampon (plateau), le point de demi-équivalence $E/2$ ($\mathrm{pH}=\mathrm{p}K_a$), un saut de pH moins vertical et un pH à l'équivalence $>7$ (basique, ions \ce{CH3COO-}).}
\end{popfigure}

\begin{methode}[Décrire et justifier l'allure de la courbe $\mathrm{pH} = f(V_B)$ pour acide faible / base forte]
\begin{enumerate}
\item \textbf{Avant le dosage} ($V_B = 0$) : pH acide, $\mathrm{pH} \approx \tfrac{1}{2}(\mathrm{p}K_a - \log C_A)$.
\item \textbf{Avant l'équivalence} ($0 < V_B < V_{BE}$) : le titrant (\ce{OH-}) est limitant. Formation d'un mélange \ce{HA}/\ce{A-} : \textbf{zone tampon}. Le pH augmente lentement.
\item \textbf{À la demi-équivalence} ($V_B = V_{BE}/2$) : $[\ce{HA}] = [\ce{A-}]$, d'où $\boxed{\mathrm{pH} = \mathrm{p}K_a}$.
\item \textbf{À l'équivalence} ($V_B = V_{BE}$) : tout \ce{HA} transformé en \ce{A-}. L'ion \ce{A-} est une base faible (\ce{A- + H2O <=> HA + OH-}). La solution est \textbf{basique} : $\mathrm{pH}_E > 7$.
\item \textbf{Après l'équivalence} ($V_B > V_{BE}$) : excès de \ce{OH-} (base forte). pH tend vers celui de la solution titrante ($\mathrm{pH} = 14 + \log C_B$).
\end{enumerate}
\end{methode}

\begin{exoresolu}[Prévision des points remarquables d'une courbe]
On dose $V_A = \SI{20,0}{mL}$ d'une solution d'acide éthanoïque de concentration $C_A = \SI{0,100}{mol.L^{-1}}$ ($\mathrm{p}K_a(\ce{CH3COOH}/\ce{CH3COO-}) = 4{,}8$) par la soude à $C_B = \SI{0,100}{mol.L^{-1}}$.
\begin{enumerate}
\item Calculer le volume équivalent $V_{BE}$.
\item Donner, en le justifiant, le pH à la demi-équivalence.
\item Le pH à l'équivalence est-il inférieur, égal ou supérieur à 7 ? Justifier.
\end{enumerate}
\tcblower
\textbf{1. Volume équivalent.} La réaction de dosage est :
\[ \ce{CH3COOH + OH- -> CH3COO- + H2O} \]
À l'équivalence : $C_A V_A = C_B V_{BE}$, d'où :
\[ V_{BE} = \frac{C_A V_A}{C_B} = \frac{0{,}100 \times 20{,}0}{0{,}100} = \SI{20,0}{mL}. \]

\textbf{2. Demi-équivalence.} À $V_B = \dfrac{V_{BE}}{2} = \SI{10,0}{mL}$, la moitié de l'acide a été transformée : $[\ce{CH3COOH}] = [\ce{CH3COO-}]$. La relation d'Henderson-Hasselbalch donne :
\[ \mathrm{pH} = \mathrm{p}K_a + \log\frac{[\ce{CH3COO-}]}{[\ce{CH3COOH}]} = \mathrm{p}K_a + \log 1 = \boxed{4{,}8}. \]

\textbf{3. pH à l'équivalence.} À l'équivalence, la solution contient des ions éthanoate \ce{CH3COO-}, qui constituent une \textbf{base faible} (ils réagissent avec l'eau : \ce{CH3COO- + H2O <=> CH3COOH + OH-}). La solution est donc \textbf{basique} : $\mathrm{pH}_E > 7$ (en pratique, voisin de 8,7 pour ces concentrations).

\textbf{Conclusion :} $V_{BE} = \SI{20,0}{mL}$ ; $\mathrm{pH}_{E/2} = 4{,}8$ ; $\mathrm{pH}_E > 7$.
\end{exoresolu}

\courssub{Dosage d'un acide fort par une base faible}

\begin{popfigure}
\begin{tikzpicture}
\begin{axis}[
    width=\linewidth, height=0.55\linewidth,
    xlabel={Volume d'acide versé $V_A$ (mL)}, ylabel={pH},
    xmin=0, xmax=45, ymin=0, ymax=14,
    grid=major, grid style={popmuted}, tick label style={font=\small},
    axis line style={popdark}, tick style={popdark},
]
% Simulated curve Weak Base (NH3, pKb=4.74 -> pKa NH4+=9.26, 0.1M) / Strong Acid (0.1M), Vb=20mL -> Vae=20mL
% Before eq: Buffer NH4+/NH3. pH = pKa + log([B]/[BH+]) = 9.26 + log((Cb*Vb - Ca*Va)/(Ca*Va))
\addplot[popblue, thick, smooth, domain=0.1:19.9, samples=150] { 9.26 + log10( (20 - x) / x ) };
% Vertical jump at eq
\draw[popblue, thick] (axis cs:20,5.5) -- (axis cs:20,2.5);
% After eq: excess H3O+
\addplot[popblue, thick, smooth, domain=20.1:45, samples=100] { -log10( 0.1*(x - 20) / (20 + x) ) };
% Half-equivalence
\node[circle, fill=popblue, inner sep=2pt, draw=popdark, pin={45:{$E/2\,(V_A=10,0;\mathrm{pH}=9,26=\mathrm{p}K_a)$}}, pin distance=1.5cm, pin edge={popdark}] at (axis cs:10,9.26) {};
% Equivalence
\node[circle, fill=popblue, inner sep=2pt, draw=popdark, pin={-90:{$E\,(V_{AE}=20,0;\mathrm{pH}_E\approx 5,3)$}}, pin distance=1.5cm, pin edge={popdark}] at (axis cs:20,5.3) {};
% Buffer zone label (replacing brace decoration unavailable in allowed libraries)
\node[popgreen, font=\bfseries] at (axis cs:10,11.5) {Zone tampon \ce{NH3}/\ce{NH4+}};
\end{axis}
\end{tikzpicture}
\legende{Courbe de dosage pH-métrique base faible (ammoniac, $\mathrm{p}K_a(\ce{NH4+}/\ce{NH3})=9,26$) / acide fort. La courbe est décroissante. Zone tampon \ce{NH3}/\ce{NH4+}, demi-équivalence $\mathrm{pH}=\mathrm{p}K_a$, saut de pH, équivalence acide ($\mathrm{pH}_E < 7$, ions \ce{NH4+}).}
\end{popfigure}

\begin{methode}[Particularités du dosage acide fort / base faible (ex. \ce{HCl} / \ce{NH3})]
\begin{enumerate}
\item \textbf{Équation-bilan :} \ce{H3O+ + NH3 -> NH4+ + H2O} (totale, rapide).
\item \textbf{Avant l'équivalence :} excès de base faible \ce{NH3} + \ce{NH4+} formé $\rightarrow$ \textbf{zone tampon} (pH basique, décroît lentement).
\item \textbf{Demi-équivalence :} $[\ce{NH3}] = [\ce{NH4+}] \Rightarrow \mathrm{pH} = \mathrm{p}K_a(\ce{NH4+}/\ce{NH3}) \approx 9,26$.
\item \textbf{À l'équivalence :} seule \ce{NH4+} (acide faible) en solution $\rightarrow$ \textbf{pH acide} ($\mathrm{pH}_E < 7$, typiquement 5--6).
\item \textbf{Après l'équivalence :} excès d'acide fort \ce{H3O+} $\rightarrow$ pH chute vers valeurs très acides.
\end{enumerate}
\end{methode}

\courssub{Dosage colorimétrique et choix de l'indicateur}

\begin{propriete}[Critère de choix d'un indicateur coloré]
Un indicateur coloré \ce{InH}/\ce{In-} convient pour un dosage si sa \textbf{zone de virage} (intervalle de pH où le changement de couleur est visible) est \textbf{incluse dans le saut de pH} de la courbe de titration et \textbf{contient le pH à l'équivalence} $\mathrm{pH}_E$.
\end{propriete}

\begin{popfigure}
\begin{tikzpicture}
% pH scale vertical
\draw[popdark, thick] (0,0) -- (0,14);
\foreach \y in {0,1,...,14} \draw[popdark] (-0.1,\y) -- (0.1,\y) node[right, font=\tiny] {\y};
% Indicator ranges
\draw[poporange, line width=4pt] (0.5,3.1) -- (0.5,4.4) node[right, midway, font=\small, poporange] {Hélianthine (3,1--4,4)};
\draw[popblue, line width=4pt] (1.5,6.0) -- (1.5,7.6) node[right, midway, font=\small, popblue] {BBT (6,0--7,6)};
\draw[poppink, line width=4pt] (2.5,8.2) -- (2.5,10.0) node[right, midway, font=\small, poppink] {Phénolphtaléine (8,2--10,0)};
% Equivalence zones
\node[popgreen, align=center, font=\small] at (4.5,7) {Acide fort / Base forte\\$\mathrm{pH}_E=7$ $\rightarrow$ BBT};
\node[popgreen, align=center, font=\small] at (4.5,9) {Acide faible / Base forte\\$\mathrm{pH}_E>7$ $\rightarrow$ Phénolphtaléine};
\node[popgreen, align=center, font=\small] at (4.5,4) {Acide fort / Base faible\\$\mathrm{pH}_E<7$ $\rightarrow$ Hélianthine};
\end{tikzpicture}
\legende{Échelle de pH et zones de virage des trois indicateurs du programme. Le choix se fait par inclusion de la zone de virage dans le saut de pH (zone hachurée grise schématique) autour de $\mathrm{pH}_E$.}
\end{popfigure}

\begin{methode}[Choix de l'indicateur coloré]
\begin{enumerate}
\item \textbf{Déterminer le pH à l'équivalence} (lecture graphique ou raisonnement : 7 pour fort/fort ; $>7$ pour faible/fort ; $<7$ pour fort/faible).
\item \textbf{Choisir l'indicateur dont la zone de virage contient $\mathrm{pH}_E$} :
\begin{center}
\begin{tabularx}{\linewidth}{|l|c|c|}
\hline
\rowcolor{popblueL} \textbf{Indicateur} & \textbf{Zone de virage} & \textbf{Dosage adapté} \\
\hline
Hélianthine & 3,1 -- 4,4 & acide fort -- base faible (\ce{HCl}/\ce{NH3}) \\
\hline
Bleu de bromothymol (BBT) & 6,0 -- 7,6 & acide fort -- base forte (\ce{HCl}/\ce{NaOH}) \\
\hline
Phénolphtaléine & 8,2 -- 10,0 & acide faible -- base forte (\ce{CH3COOH}/\ce{NaOH}) \\
\hline
\end{tabularx}
\end{center}
\item \textbf{Vérifier} que le virage se produit dans le saut de pH : le changement de teinte doit être net pour quelques gouttes de titrant versé.
\item \textbf{Observer} : repérer l'équivalence au changement persistant de couleur (ex. apparition du rose pâle persistant de la phénolphtaléine).
\end{enumerate}
\end{methode}

\begin{exoresolu}[Choix d'indicateur pour trois dosages]
On dispose de trois indicateurs : hélianthine (3,1--4,4), BBT (6,0--7,6), phénolphtaléine (8,2--10,0). Choisir, en justifiant, l'indicateur adapté à chacun des dosages suivants :
\begin{enumerate}
\item dosage de l'acide chlorhydrique par la soude ($\mathrm{pH}_E = 7{,}0$) ;
\item dosage de l'acide éthanoïque du vinaigre par la soude ($\mathrm{pH}_E \approx 8{,}7$) ;
\item dosage de l'ammoniac par l'acide chlorhydrique ($\mathrm{pH}_E \approx 5{,}5$).
\end{enumerate}
\tcblower
\textbf{Principe :} un indicateur convient si sa zone de virage est incluse dans le saut de pH et contient le pH à l'équivalence.

\textbf{1. \ce{HCl} / soude :} $\mathrm{pH}_E = 7{,}0 \in [6{,}0\,;\,7{,}6]$ : on choisit le \textbf{BBT} (virage du jaune au bleu, vert à l'équivalence). La phénolphtaléine resterait utilisable car le saut de pH est très large (environ 3 à 11), mais le BBT est le choix optimal.

\textbf{2. Acide éthanoïque / soude :} $\mathrm{pH}_E \approx 8{,}7 \in [8{,}2\,;\,10{,}0]$ : on choisit la \textbf{phénolphtaléine} (virage de l'incolore au rose). L'hélianthine virerait bien avant l'équivalence : erreur grossière.

\textbf{3. Ammoniac / \ce{HCl} :} $\mathrm{pH}_E \approx 5{,}5$ : seule la zone de virage de l'\textbf{hélianthine} (3,1--4,4) est proche et incluse dans le saut de pH (qui s'étend environ de 6,5 à 3) : on choisit l'hélianthine (virage du jaune orangé au rouge).

\textbf{Conclusion :} BBT, phénolphtaléine, hélianthine respectivement.
\end{exoresolu}

\courssub{Applications : vérification d'étiquettes et degré d'acidité}

\begin{methode}[Calculs de dosage : tableau d'avancement et dilutions]
\begin{enumerate}
\item \textbf{Écrire l'équation} de la réaction de dosage (réactif titré + titrant).
\item \textbf{Dresser un tableau d'avancement} si la situation est complexe (mélange non équivalent, excès) : à l'équivalence, l'avancement maximal consomme les deux réactifs simultanément.
\item \textbf{Appliquer la relation d'équivalence} en tenant compte de la stoechiométrie : pour \ce{H2SO4} (diacide) dosé par la soude, $2\,C_A V_A = C_B V_{BE}$.
\item \textbf{Attention aux dilutions préalables} : si l'échantillon a été dilué $n$ fois avant le dosage, remonter à la concentration de la solution commerciale : $C_0 = n \times C_{\text{dosée}}$.
\item \textbf{Convertir} si demandé : concentration massique $C_m = C \times M$ ; pourcentage massique ; degré d'acidité d'un vinaigre (grammes d'acide éthanoïque pour \SI{100}{g} de vinaigre, assimilé à \SI{100}{mL} car $d \approx 1$).
\end{enumerate}
\end{methode}

\begin{experience}[TP : Détermination du degré d'acidité d'un vinaigre (bili-bili ou vinaigre de palme)]
\textbf{Objectif :} Vérifier la conformité d'un vinaigre local (degré $\ge 5^\circ$).
\textbf{Matériel :} Fiole jaugée \SI{100}{mL} (dilution 10), pipette jaugée \SI{10}{mL}, erlenmeyer \SI{100}{mL}, burette \SI{50}{mL} (soude \SI{0,100}{mol.L^{-1}}), phénolphtaléine, bécher, entonnoir.
\textbf{Protocole :}
\begin{enumerate}
\item \textbf{Dilution 10} : Pipeter \SI{10,0}{mL} de vinaigre commercial dans la fiole \SI{100}{mL}, compléter à l'eau distillée, homogénéiser.
\item \textbf{Dosage} : Pipeter \SI{10,0}{mL} de solution diluée dans l'erlenmeyer, ajouter 2 gouttes de phénolphtaléine.
\item \textbf{Titrage} : Remplir la burette de soude, chasser les bulles, noter $V_{B0}$. Titrer jusqu'au virage rose \textbf{persistant} (30 s). Noter $V_{B1}$. $V_{BE} = V_{B1} - V_{B0}$.
\item \textbf{Concordance} : Refaire 2 fois. Écarts $< \SI{0,2}{mL}$. Moyenne $V_{BE}$.
\item \textbf{Calculs} : $C_{\text{diluée}} = \frac{C_B V_{BE}}{V_A}$ ; $C_{\text{commercial}} = 10 \times C_{\text{diluée}}$ ; Degré = $C_{\text{commercial}} \times M(\ce{CH3COOH})$ (en g/100mL).
\end{enumerate}
\end{experience}

\begin{exoresolu}[Vérification de l'étiquette d'un déboucheur]
L'étiquette d'une solution commerciale d'acide chlorhydrique (déboucheur) indique « 23 \% en masse de \ce{HCl} », densité $d = 1{,}12$. On dilue cette solution 100 fois, puis on dose $V_A = \SI{20,0}{mL}$ de la solution diluée par la soude à $C_B = \SI{0,100}{mol.L^{-1}}$ : $V_{BE} = \SI{15,0}{mL}$. $M(\ce{HCl}) = \SI{36,5}{g.mol^{-1}}$. L'étiquette est-elle exacte ?
\tcblower
\textbf{Étape 1 : concentration de la solution diluée.}
\[ C_d = \frac{C_B V_{BE}}{V_A} = \frac{0{,}100 \times 15{,}0}{20{,}0} = \SI{7,5e-2}{mol.L^{-1}}. \]

\textbf{Étape 2 : concentration commerciale.} Dilution 100 fois :
\[ C_0 = 100 \times C_d = \SI{7,5}{mol.L^{-1}}. \]

\textbf{Étape 3 : concentration massique.} $C_m = C_0 \times M = 7{,}5 \times 36{,}5 = \SI{273,75}{g.L^{-1}} \approx \SI{274}{g.L^{-1}}$.

\textbf{Étape 4 : pourcentage massique.} Un litre de solution commerciale a une masse $m = d \times \rho_{\text{eau}} \times V = 1{,}12 \times 1000 = \SI{1120}{g}$. D'où :
\[ \%\,m = \frac{273{,}75}{1120}\times 100 = 24{,}4\,\% \approx 23\,\%\ \text{(à l'incertitude expérimentale près)}. \]

\textbf{Conclusion :} l'écart relatif $\dfrac{24{,}4 - 23}{23} \approx 6\,\%$ est de l'ordre des erreurs de manipulation ; \textbf{l'étiquette est compatible} avec le résultat du dosage.
\end{exoresolu}

\begin{exoresolu}[Degré d'acidité d'un vinaigre du marché de Mokolo]
Un vinaigre est dilué 10 fois. On dose $V_A = \SI{10,0}{mL}$ de la solution diluée par la soude à $C_B = \SI{0,100}{mol.L^{-1}}$ en présence de phénolphtaléine. Le virage persistant est obtenu pour $V_{BE} = \SI{12,0}{mL}$. $M(\ce{CH3COOH}) = \SI{60}{g.mol^{-1}}$ ; on assimile la densité du vinaigre à 1.
\begin{enumerate}
\item Écrire l'équation de la réaction de dosage.
\item Calculer la concentration en acide éthanoïque du vinaigre commercial.
\item En déduire son degré d'acidité et conclure (un vinaigre doit titrer au moins 5 degrés).
\end{enumerate}
\tcblower
\textbf{1. Équation de dosage :}
\[ \ce{CH3COOH + OH- -> CH3COO- + H2O} \]
(réaction totale, rapide ; la phénolphtaléine convient car $\mathrm{pH}_E > 7$.)

\textbf{2. Concentration du vinaigre.} Solution diluée :
\[ C_d = \frac{C_B V_{BE}}{V_A} = \frac{0{,}100 \times 12{,}0}{10{,}0} = \SI{0,120}{mol.L^{-1}}. \]
Vinaigre commercial (dilution 10 fois) : $C_0 = 10 \times 0{,}120 = \SI{1,20}{mol.L^{-1}}$.

\textbf{3. Degré d'acidité.} Concentration massique :
\[ C_m = C_0 \times M = 1{,}20 \times 60 = \SI{72}{g.L^{-1}}, \]
soit \SI{7,2}{g} d'acide éthanoïque pour \SI{100}{mL} de vinaigre, c'est-à-dire pour environ \SI{100}{g} ($d \approx 1$).
\[ \boxed{\text{Degré d'acidité} \approx 7{,}2^\circ} \]

\textbf{Conclusion :} $7{,}2^\circ > 5^\circ$ : ce vinaigre est \textbf{conforme} à la réglementation ; son étiquette peut annoncer un vinaigre de qualité.
\end{exoresolu}

\courssub{Exercices d'entraînement}

\begin{exercice}[Dosage d'une solution de soude commerciale (lessive artisanale)]
On souhaite vérifier la concentration en soude (\ce{NaOH}) d'une solution utilisée pour la fabrication artisanale de savon (saponification de l'huile de palme). On prélève $V_0 = \SI{10,0}{mL}$ de cette solution que l'on dilue à \SI{250,0}{mL} (fiole jaugée). On dose ensuite $V_A = \SI{20,0}{mL}$ de cette solution diluée par une solution d'acide chlorhydrique de concentration $C_A = \SI{0,100}{mol.L^{-1}}$.
\begin{enumerate}
\item Écrire l'équation-bilan de la réaction de dosage.
\item L'équivalence est atteinte pour $V_{AE} = \SI{18,5}{mL}$. Calculer la concentration de la solution de soude commerciale.
\item Quel indicateur coloré choisir ? Justifier.
\end{enumerate}
\end{exercice}

\begin{exercice}[Dosage de l'ammoniac dans un produit de nettoyage]
Un produit de nettoyage « Eau de Javel parfumée » contient de l'ammoniac \ce{NH3}. On dose $V_B = \SI{10,0}{mL}$ de ce produit (dilué 10 fois au préalable) par de l'acide chlorhydrique $C_A = \SI{0,100}{mol.L^{-1}}$.
\begin{enumerate}
\item Écrire l'équation-bilan. Préciser la nature de la réaction.
\item On obtient $V_{AE} = \SI{12,0}{mL}$ au virage de l'hélianthine. Calculer la concentration en ammoniac du produit commercial.
\item Justifier le choix de l'hélianthine. Pourquoi la phénolphtaléine est-elle inadaptée ?
\end{enumerate}
\end{exercice}

\begin{corrige}[Exercice 1 : Dosage soude commerciale]
\textbf{1. Équation-bilan :} \ce{H3O+ + OH- -> 2H2O} (totale, rapide, exothermique).
\textbf{2. Calcul :} Solution diluée : $C_{\text{dil}} V_A = C_A V_{AE} \Rightarrow C_{\text{dil}} = \frac{0,100 \times 18,5}{20,0} = \SI{9,25e-2}{mol.L^{-1}}$.
Solution commerciale (dilution 25 fois : $250/10$) : $C_0 = 25 \times 9,25\times10^{-2} = \SI{2,31}{mol.L^{-1}}$.
\textbf{3. Indicateur :} Acide fort / Base forte $\Rightarrow \mathrm{pH}_E = 7$. Choix : \textbf{Bleu de Bromothymol (BBT)} (zone 6,0--7,6).
\end{corrige}

\begin{corrige}[Exercice 2 : Dosage ammoniac]
\textbf{1. Équation-bilan :} \ce{H3O+ + NH3 -> NH4+ + H2O}. Réaction totale (acide fort + base faible), rapide.
\textbf{2. Calcul :} Solution diluée : $C_{\text{dil}} V_B = C_A V_{AE} \Rightarrow C_{\text{dil}} = \frac{0,100 \times 12,0}{10,0} = \SI{0,120}{mol.L^{-1}}$.
Produit commercial (dilution 10 fois) : $C_0 = 10 \times 0,120 = \SI{1,20}{mol.L^{-1}}$.
\textbf{3. Choix indicateur :} À l'équivalence, solution de \ce{NH4+} (acide faible) $\Rightarrow \mathrm{pH}_E < 7$ (environ 5,5). L'hélianthine (3,1--4,4) a sa zone dans le saut de pH. La phénolphtaléine (8,2--10) vire en zone basique, bien avant l'équivalence $\rightarrow$ erreur massive (volume mesuré trop petit).
\end{corrige}

\begin{attention}[Les 5 erreurs qui coûtent des points au Bac]
\begin{enumerate}
\item \textbf{Écrire \ce{HCl + NaOH -> NaCl + H2O} comme équation de dosage.} Au Bac, on exige l'équation ionique réelle : \ce{H3O+ + OH- -> 2H2O} (ou \ce{HA + OH- -> A- + H2O} pour un acide faible). \ce{Na+} et \ce{Cl-} sont spectateurs.
\item \textbf{Oublier la dilution préalable.} La relation $C_A V_A = C_B V_{BE}$ donne la concentration de la solution \emph{dosée} ; il faut multiplier par le facteur de dilution pour retrouver la concentration commerciale. C'est l'erreur n° 1 des exercices « vérification d'étiquette ».
\item \textbf{Affirmer que le pH à l'équivalence vaut toujours 7.} Faux : $\mathrm{pH}_E = 7$ uniquement pour acide fort/base forte ; $\mathrm{pH}_E > 7$ (acide faible/base forte, ion \ce{A-} base faible) et $\mathrm{pH}_E < 7$ (acide fort/base faible, ion \ce{BH+} acide faible). Justifier par la nature de l'ion formé.
\item \textbf{Mal choisir l'indicateur coloré.} Le critère est : zone de virage \emph{contenant} le pH à l'équivalence et incluse dans le saut de pH. Utiliser la phénolphtaléine pour doser l'ammoniac par \ce{HCl} est une faute de raisonnement sanctionnée.
\item \textbf{Négliger les unités et les chiffres significatifs.} Les volumes peuvent rester en mL dans $C_A V_A = C_B V_{BE}$ (le rapport suffit), mais toute quantité de matière exige des litres. Donner le résultat avec 2 ou 3 chiffres significatifs et une unité (\si{mol.L^{-1}}), et conclure par une phrase.
\end{enumerate}
\end{attention}

\newpage

\coursec{Exercices}

\courssub{Je vérifie mes connaissances}

\begin{aretenir}[Rappels essentiels : équivalence et demi-équivalence]
L'\cle{équivalence acido-basique} est l'état du système pour lequel les réactifs ont été introduits dans les proportions stœchiométriques exactes. Pour un dosage d'un acide (A) par une base (B) : $n_A(\text{initial}) = n_B(\text{versé à l'équivalence})$, soit $C_A V_A = C_B V_{BE}$ (si coefficients stœchiométriques 1:1).
Le \cle{point de demi-équivalence} (uniquement pour les couples faible/fort) correspond au volume $V_{B1/2} = V_{BE}/2$. En ce point, $\mathrm{pH} = \mathrm{p}K_a$ de l'acide dosé (ou $\mathrm{pH} = \mathrm{p}K_a$ de l'acide conjugué de la base dosée).
\end{aretenir}

\begin{exercice}[QCM : l'équivalence acido-basique]
Pour chaque question, une seule réponse est exacte. Entoure la bonne lettre.

\textbf{1.} À l'équivalence d'un dosage acido-basique :
\begin{itemize}
\item[a)] le pH est toujours égal à 7 ;
\item[b)] les réactifs ont été mélangés dans les proportions stœchiométriques ;
\item[c)] l'indicateur coloré est toujours incolore.
\end{itemize}

\textbf{2.} Le dosage d'un acide fort par une base forte donne à l'équivalence un pH :
\begin{itemize}
\item[a)] inférieur à 7 ; \quad b) égal à 7 ; \quad c) supérieur à 7.
\end{itemize}

\textbf{3.} Au point de demi-équivalence du dosage d'un acide faible par une base forte :
\begin{itemize}
\item[a)] $\mathrm{pH} = 7$ ; \quad b) $\mathrm{pH} = \mathrm{p}K_a$ ; \quad c) $\mathrm{pH} = 14$.
\end{itemize}

\textbf{4.} Pour doser l'acide éthanoïque par la soude ($\mathrm{pH}_E \approx 8,7$), l'indicateur le mieux adapté est :
\begin{itemize}
\item[a)] l'hélianthine (zone de virage 3,1 -- 4,4) ;
\item[b)] le bleu de bromothymol (zone de virage 6,0 -- 7,6) ;
\item[c)] la phénolphtaléine (zone de virage 8,2 -- 10,0).
\end{itemize}
\end{exercice}

\begin{exercice}[Vrai ou faux, en justifiant]
Réponds par \textbf{vrai} ou \textbf{faux}, puis justifie ta réponse en une ou deux phrases.
\begin{enumerate}
\item La réaction entre les ions \ce{H3O+} et \ce{OH-} est lente et endothermique.
\item Lors du dosage d'un acide fort par une base faible, le pH à l'équivalence est inférieur à 7.
\item La méthode des tangentes parallèles permet de déterminer graphiquement le volume équivalent.
\item Un indicateur coloré est bien choisi si sa zone de virage contient le pH initial de la solution.
\item À la demi-équivalence du dosage de \ce{CH3COOH} par \ce{NaOH}, on a $[\ce{CH3COOH}] = [\ce{CH3COO-}]$.
\end{enumerate}
\end{exercice}

\begin{exercice}[Texte à trous]
Complète le texte suivant avec les mots ou expressions : \emph{exothermique}, \emph{totale}, \emph{équivalence}, \emph{tangentes}, \emph{stœchiométriques}, \emph{dérivée}, \emph{zone de virage}.

La réaction entre un acide fort et une base forte est \ldots\ldots\ldots\ldots, rapide et \ldots\ldots\ldots\ldots : elle dégage de la chaleur. L'\ldots\ldots\ldots\ldots est atteinte lorsque les réactifs ont été mélangés dans les proportions \ldots\ldots\ldots\ldots. Sur la courbe $\mathrm{pH} = f(V_B)$, le point d'équivalence se détermine par la méthode des \ldots\ldots\ldots\ldots parallèles ou par la méthode de la \ldots\ldots\ldots\ldots, qui présente un extremum à l'équivalence. Pour un dosage colorimétrique, on choisit un indicateur dont la \ldots\ldots\ldots\ldots contient le pH à l'équivalence.
\end{exercice}

\begin{exercice}[Définitions et équations]
\begin{enumerate}
\item Définis l'\cle{équivalence acido-basique} et écris la relation entre les quantités de matière qui la traduit pour le dosage d'un acide par une base.
\item Écris l'équation-bilan de la réaction de dosage : a) entre l'acide chlorhydrique et la soude ; b) entre l'acide éthanoïque et la soude ; c) entre l'ammoniac et l'acide chlorhydrique.
\item Donne la définition du \cle{point de demi-équivalence} et la propriété du pH en ce point.
\end{enumerate}
\end{exercice}

\courssub{Je m'entraîne}

\begin{methode}[Détermination graphique du volume équivalent $V_E$ (Méthode des tangentes parallèles)]
\begin{enumerate}
\item Tracer la courbe $\mathrm{pH} = f(V_{\text{titrant}})$ point par point.
\item Tracer la tangente à la branche avant le saut (zone tampon ou acide/base fort) et la tangente à la branche après le saut (excès de titrant).
\item Tracer la tangente à la partie verticale (au point d'inflexion), parallèle aux deux précédentes.
\item Les intersections de cette tangente centrale avec les deux tangentes externes se projettent sur l'axe des abscisses au même point : $V_E$.
\end{enumerate}
\end{methode}

\begin{exercice}[Volume équivalent d'un dosage simple]
On dose un volume $V_A = \SI{15,0}{mL}$ d'une solution d'acide éthanoïque de concentration $C_A = \SI{0,080}{mol.L^{-1}}$ par une solution de soude de concentration $C_B = \SI{0,10}{mol.L^{-1}}$.
\begin{enumerate}
\item Écris l'équation-bilan de la réaction de dosage.
\item Calcule le volume $V_{BE}$ de soude nécessaire pour atteindre l'équivalence.
\item Sachant que $\mathrm{p}K_a(\ce{CH3COOH}/\ce{CH3COO-}) = 4,8$, donne la valeur du pH à la demi-équivalence, puis précise si le pH à l'équivalence est inférieur, égal ou supérieur à 7. Justifie.
\end{enumerate}
\end{exercice}

\begin{exercice}[Prévision d'une courbe de dosage]
On dose $V_B = \SI{25,0}{mL}$ d'une solution d'ammoniac de concentration $C_B = \SI{0,10}{mol.L^{-1}}$ par une solution d'acide chlorhydrique de concentration $C_A = \SI{0,10}{mol.L^{-1}}$. On donne $\mathrm{p}K_a(\ce{NH4+}/\ce{NH3}) = 9,2$.
\begin{enumerate}
\item Écris l'équation-bilan de la réaction de dosage.
\item Calcule le volume équivalent $V_{AE}$.
\item Prévois le sens de variation du pH au cours du dosage et indique, en justifiant, si le pH à l'équivalence est inférieur, égal ou supérieur à 7.
\item Choisis l'indicateur coloré le mieux adapté parmi : hélianthine (3,1 -- 4,4), BBT (6,0 -- 7,6), phénolphtaléine (8,2 -- 10,0).
\end{enumerate}
\end{exercice}

\begin{attention}[Piège classique : Choix de l'indicateur pour base faible / acide fort]
Pour un dosage \textbf{base faible / acide fort}, $\mathrm{pH}_E < 7$. L'indicateur doit virer en domaine \textbf{acide}. Parmi la liste officielle (Hélianthine, BBT, Phénolphtaléine), seule l'hélianthine est acide. \textbf{Toutefois}, sa zone (3,1--4,4) est souvent trop acide par rapport au $\mathrm{pH}_E$ réel (souvent 5--6). Au Bac, si l'hélianthine est le seul indicateur acide proposé, on le choisit en notant la limite. Vérifie toujours $\mathrm{pH}_E$ par le calcul si les concentrations sont données.
\end{attention}

\begin{exercice}[Exploitation d'un tableau de mesures]
Au cours du dosage de $V_A = \SI{20,0}{mL}$ d'acide chlorhydrique par la soude à $C_B = \SI{0,10}{mol.L^{-1}}$, on obtient les mesures suivantes :

\begin{center}
\begin{tabularx}{\linewidth}{|l|X|X|X|X|X|X|X|X|}
\hline
\rowcolor{popblueL} $V_B$ (mL) & 0 & 5 & 10 & 14 & 15 & 16 & 18 & 20 \\
\hline pH & 1,5 & 1,8 & 2,2 & 2,9 & 7,0 & 11,0 & 11,7 & 12,0 \\
\hline
\end{tabularx}
\end{center}

\begin{enumerate}
\item Trace la courbe $\mathrm{pH} = f(V_B)$ sur papier millimétré (échelles : \SI{1}{cm} pour \SI{2}{mL} et \SI{1}{cm} pour 1 unité de pH).
\item Détermine graphiquement le volume équivalent $V_{BE}$ et le pH à l'équivalence.
\item En déduis la concentration $C_A$ de la solution d'acide chlorhydrique.
\item Le saut de pH est-il compatible avec un dosage colorimétrique au BBT ? Justifie.
\end{enumerate}
\end{exercice}

\begin{exercice}[Les trois flacons sans étiquette]
Au laboratoire du lycée, trois flacons ont perdu leur étiquette. Ils contiennent respectivement une solution d'acide chlorhydrique, une solution d'acide éthanoïque et une solution d'ammoniac, toutes trois de concentration \SI{0,10}{mol.L^{-1}}. On donne $\mathrm{p}K_a(\ce{CH3COOH}/\ce{CH3COO-}) = 4,8$ et $\mathrm{p}K_a(\ce{NH4+}/\ce{NH3}) = 9,2$.
\begin{enumerate}
\item Propose un protocole expérimental de dosage pH-métrique permettant d'identifier le contenu de chaque flacon (réactif titrant, verrerie, grandeur suivie).
\item Pour chaque flacon, précise l'allure attendue de la courbe (sens de variation du pH, position approximative du saut) et la valeur attendue du pH à l'équivalence (inférieure, égale ou supérieure à 7).
\item Explique comment la valeur du pH à la demi-équivalence permet de confirmer l'identification des deux espèces faibles.
\end{enumerate}
\end{exercice}

\courssub{Je me prépare au Bac}

\begin{methode}[Vérification d'étiquette par dilution et dosage]
\begin{enumerate}
\item \textbf{Dilution connue :} $S_0 \xrightarrow{\text{facteur } f} S_1$ (pipette + fiole jaugée). $C_1 = C_0 / f$.
\item \textbf{Dosage de $S_1$ :} Prélèvement $V_A$ de $S_1$ + indicateur ou pH-mètre. Titrage par solution titrée $C_B$. Mesure $V_{BE}$.
\item \textbf{Calcul remonté :} $C_A (S_1) = \frac{C_B V_{BE}}{V_A}$ (stoechio 1:1). Puis $C_0 = f \times C_A$.
\item \textbf{Conclusion :} Comparer $C_0$ trouvé à la valeur étiquette. Préciser l'incertitude si demandée.
\end{enumerate}
\end{methode}

\begin{exercice}[Type Bac : vérification de l'étiquette d'une solution d'acide chlorhydrique]
Au laboratoire d'un lycée de Douala, on dispose d'une solution commerciale $S_0$ d'acide chlorhydrique dont l'étiquette indique « concentration : \SI{10}{mol.L^{-1}} ». Afin de vérifier cette indication, on dilue $S_0$ cent fois : on obtient une solution $S_1$. On prélève $V_A = \SI{20,0}{mL}$ de $S_1$ que l'on dose par une solution d'hydroxyde de sodium de concentration $C_B = \SI{0,10}{mol.L^{-1}}$, en suivant le pH après chaque addition. La courbe $\mathrm{pH} = f(V_B)$ obtenue présente un saut de pH important, et la méthode des tangentes parallèles donne un point d'équivalence de coordonnées $V_{BE} = \SI{16,0}{mL}$ et $\mathrm{pH}_E = 7,0$.
\begin{enumerate}
\item Fais le schéma annoté du dispositif de dosage pH-métrique.
\item Écris l'équation-bilan de la réaction de dosage et donne deux caractéristiques de cette réaction.
\item Définis l'équivalence acido-basique et décris brièvement le principe de la méthode des tangentes parallèles utilisée pour repérer le point d'équivalence.
\item Calcule la concentration $C_A$ de la solution diluée $S_1$, puis la concentration $C_0$ de la solution commerciale $S_0$.
\item L'indication de l'étiquette est-elle exacte ? Justifie ta réponse.
\item Parmi l'hélianthine (3,1 -- 4,4), le BBT (6,0 -- 7,6) et la phénolphtaléine (8,2 -- 10,0), quel(s) indicateur(s) coloré(s) pourrai(en)t remplacer le pH-mètre pour ce dosage ? Justifie.
\end{enumerate}
\end{exercice}

\begin{exercice}[Type Bac : degré d'acidité d'un vinaigre du marché de Mokolo]
Le \cle{degré d'acidité} d'un vinaigre est la masse d'acide éthanoïque, en grammes, contenue dans \SI{100}{g} de vinaigre. Une coopérative de Yaoundé commercialise un vinaigre étiqueté « 5° ». Pour vérifier cette étiquette, on dilue le vinaigre dix fois (solution $S$). On dose ensuite $V_A = \SI{10,0}{mL}$ de $S$ par une solution de soude de concentration $C_B = \SI{0,020}{mol.L^{-1}}$. L'équivalence est atteinte pour $V_{BE} = \SI{12,0}{mL}$.

\emph{Données :} $M(\ce{CH3COOH}) = \SI{60}{g.mol^{-1}}$ ; masse volumique du vinaigre : $\rho = \SI{1,0}{g.mL^{-1}}$ ; $\mathrm{p}K_a(\ce{CH3COOH}/\ce{CH3COO-}) = 4,8$.
\begin{enumerate}
\item Pourquoi dilue-t-on le vinaigre avant le dosage ? Décris le protocole de cette dilution (verrerie utilisée).
\item Écris l'équation-bilan de la réaction de dosage.
\item Calcule la concentration $C_A$ en acide éthanoïque de la solution diluée $S$, puis celle $C$ du vinaigre commercial.
\item Calcule la masse d'acide éthanoïque contenue dans \SI{1}{L} de vinaigre, puis le degré d'acidité du vinaigre.
\item L'étiquette « 5° » est-elle exacte ? Conclus.
\item Le suivi pH-métrique du dosage donne $\mathrm{pH} = 8,6$ à l'équivalence. Justifie que ce pH soit supérieur à 7, donne la valeur du pH à la demi-équivalence, et choisis l'indicateur coloré approprié parmi hélianthine (3,1 -- 4,4), BBT (6,0 -- 7,6), phénolphtaléine (8,2 -- 10,0).
\end{enumerate}
\end{exercice}

\begin{exercice}[Type Bac : dosage de l'acide éthanoïque et détermination du pKa]
On dose $V_A = \SI{20,0}{mL}$ d'une solution d'acide éthanoïque de concentration inconnue par une solution de soude de concentration $C_B = \SI{0,10}{mol.L^{-1}}$. La courbe $\mathrm{pH} = f(V_B)$ est donnée ci-dessous.

\begin{popfigure}
\begin{center}
\begin{tikzpicture}
\begin{axis}[
    width=\linewidth, height=0.55\linewidth,
    axis lines=middle,
    xlabel={$V_B$ (mL)}, ylabel={pH},
    xmin=0, xmax=26, ymin=0, ymax=14,
    xtick={0,5,10,15,20,25}, ytick={0,2,4,6,8,10,12,14},
    grid=both, grid style={popline, thin},
    major tick style={popdark, thin},
    tick label style={font=\small, popdark},
    x label style={at={(axis description cs:1,0)}, anchor=north, font=\small},
    y label style={at={(axis description cs:0,1)}, anchor=south, font=\small},
    clip=false
]
% Courbe théorique : 20 mL AcOH 0,075 M + NaOH 0,10 M
% Points clés : V0=0 pH=2,96 ; V1/2=7,5 pH=4,80 ; VE=15 pH=8,72 ; V=25 pH=12,20
\addplot[popcurve, very thick, smooth] coordinates {
    (0, 2.96)
    (1, 3.35)
    (2, 3.65)
    (3, 3.88)
    (4, 4.07)
    (5, 4.23)
    (6, 4.38)
    (7, 4.52)
    (7.5, 4.80) % Demi-équivalence
    (8, 5.05)
    (9, 5.35)
    (10, 5.65)
    (11, 6.00)
    (12, 6.40)
    (13, 6.90)
    (14, 7.55)
    (14.5, 8.05)
    (15, 8.72) % Équivalence
    (15.5, 9.45)
    (16, 10.05)
    (17, 10.70)
    (18, 11.15)
    (20, 11.65)
    (25, 12.20)
};
% Point E (équivalence)
\addplot[poporange, only marks, mark=*, mark size=2.5] coordinates {(15, 8.72)};
\node[poporange, above right, font=\small] at (axis cs:15, 8.72) {E};
\draw[poporange, dashed, thick] (axis cs:15, 0) -- (axis cs:15, 8.72) -- (axis cs:0, 8.72);
% Point D (demi-équivalence)
\addplot[popgreen, only marks, mark=*, mark size=2.5] coordinates {(7.5, 4.80)};
\node[popgreen, below right, font=\small] at (axis cs:7.5, 4.80) {D};
\draw[popgreen, dashed, thick] (axis cs:7.5, 0) -- (axis cs:7.5, 4.80) -- (axis cs:0, 4.80);
\end{axis}
\end{tikzpicture}
\end{center}
\legende{Courbe de dosage de l'acide éthanoïque ($V_A=20,0$ mL, $C_A$ inconnue) par la soude ($C_B=0,10$ mol.L$^{-1}$). E : équivalence ; D : demi-équivalence.}
\end{popfigure}

\begin{enumerate}
\item Écris l'équation-bilan de la réaction de dosage.
\item Détermine graphiquement les coordonnées du point d'équivalence E. Le pH à l'équivalence est-il cohérent avec la nature de l'acide dosé ? Justifie.
\item Calcule la concentration $C_A$ de la solution d'acide éthanoïque.
\item Définis le point de demi-équivalence D, détermine graphiquement ses coordonnées et en déduis la valeur du $\mathrm{p}K_a$ du couple $\ce{CH3COOH}/\ce{CH3COO-}$. Compare avec la valeur tabulée $\mathrm{p}K_a = 4,8$.
\item Justifie le choix de la phénolphtaléine (zone de virage 8,2 -- 10,0) comme indicateur coloré pour ce dosage, et précise le changement de teinte observé à l'équivalence.
\item Calcule le volume de soude versé à la demi-équivalence et vérifie la cohérence avec le graphique.
\end{enumerate}
\end{exercice}

\begin{savaistu}[Le vinaigre au Cameroun : tradition et contrôle]
Au Cameroun, le vinaigre est souvent produit artisanalement (vin de palme acétifié, bili-bili tourné, jus de fruits fermentés). Le « degré d'acidité » (ex: 5°, 6°) garantit la capacité de conservation (pickles, condiments) et la saveur. Un vinaigre à 5° contient 5 g d'acide éthanoïque pur pour 100 g de solution (soit $\approx 0,83\ \text{mol.L}^{-1}$). Le contrôle par dosage (norme CODEX/ANOR) protège le consommateur contre les fraudes (eau ajoutée, acide sulfurique frauduleux). La dilution avant dosage est une pratique standard au laboratoire (SONARA, laboratoires de contrôle qualité, lycées) pour adapter la concentration à la burette et à la précision du pH-mètre.
\end{savaistu}

\newpage

\coursec{Corrigés des exercices}

\coursec{Leçon 09 : Réactions acide/base et dosages}

\courssub{Je vérifie mes connaissances}

\begin{corrige}[Exercice 1 --- QCM : l'équivalence acido-basique]
\textbf{1. b)} À l'équivalence, les réactifs ont été mélangés dans les \cle{proportions stœchiométriques} de l'équation-bilan : c'est la définition. La réponse a) est fausse : $\mathrm{pH}_E = 7$ uniquement pour un dosage acide fort / base forte. La réponse c) est fausse : l'indicateur \emph{change de couleur} à l'équivalence, il n'est pas nécessairement incolore.

\textbf{2. b)} La réaction $\ce{H3O+ + OH- -> 2H2O}$ est totale ; à l'équivalence il ne reste que de l'eau et des ions spectateurs (\ce{Na+}, \ce{Cl-}) sans action sur l'eau : $\mathrm{pH}_E = 7$ à \SI{25}{\celsius}.

\textbf{3. b)} À la demi-équivalence, $[\ce{AH}] = [\ce{A-}]$, donc $\mathrm{pH} = \mathrm{p}K_a + \log\dfrac{[\ce{A-}]}{[\ce{AH}]} = \mathrm{p}K_a$.

\textbf{4. c)} Le pH à l'équivalence ($\approx 8,7$) doit appartenir à la zone de virage de l'indicateur : seule la phénolphtaléine (8,2 -- 10,0) convient. L'hélianthine virerait bien avant l'équivalence, le BBT légèrement avant.
\end{corrige}

\begin{corrige}[Exercice 2 --- Vrai ou faux]
\begin{enumerate}
\item \textbf{Faux.} La réaction $\ce{H3O+ + OH- -> 2H2O}$ est \textbf{rapide} (quasi instantanée), \textbf{totale} et \textbf{exothermique} : la température du mélange augmente pendant le dosage.
\item \textbf{Vrai.} À l'équivalence, la solution contient l'acide conjugué de la base faible (par exemple \ce{NH4+} si l'on dose \ce{NH3}), qui réagit avec l'eau : $\ce{NH4+ + H2O <=> NH3 + H3O+}$. Le milieu est donc acide : $\mathrm{pH}_E < 7$.
\item \textbf{Vrai.} On trace deux tangentes parallèles aux branches de la courbe de part et d'autre du saut, puis une troisième droite parallèle et équidistante ; son intersection avec la courbe se projette sur l'axe des volumes en $V_E$.
\item \textbf{Faux.} Un indicateur est bien choisi si sa \textbf{zone de virage contient le pH à l'équivalence} $\mathrm{pH}_E$ (et non le pH initial), afin que le virage signale précisément l'équivalence.
\item \textbf{Vrai.} À la demi-équivalence, la moitié de l'acide a été consommée et transformée en \ce{CH3COO-} : $n(\ce{CH3COOH})_{\text{restant}} = n(\ce{CH3COO-})_{\text{formé}}$, donc, dans le même volume, $[\ce{CH3COOH}] = [\ce{CH3COO-}]$, et $\mathrm{pH} = \mathrm{p}K_a = 4,8$.
\end{enumerate}
\end{corrige}

\begin{corrige}[Exercice 3 --- Texte à trous]
La réaction entre un acide fort et une base forte est \textbf{totale}, rapide et \textbf{exothermique} : elle dégage de la chaleur. L'\textbf{équivalence} est atteinte lorsque les réactifs ont été mélangés dans les proportions \textbf{stœchiométriques}. Sur la courbe $\mathrm{pH} = f(V_B)$, le point d'équivalence se détermine par la méthode des \textbf{tangentes} parallèles ou par la méthode de la \textbf{dérivée}, qui présente un extremum à l'équivalence. Pour un dosage colorimétrique, on choisit un indicateur dont la \textbf{zone de virage} contient le pH à l'équivalence.
\end{corrige}

\begin{corrige}[Exercice 4 --- Définitions et équations]
\begin{enumerate}
\item L'\cle{équivalence acido-basique} est l'état du système pour lequel les réactifs (acide dosé et base titrante) ont été introduits dans les proportions stœchiométriques exactes de l'équation-bilan : aucun des deux n'est en excès. Pour une réaction de coefficients 1:1 :
\[ n_A(\text{initial}) = n_B(\text{versé à l'équivalence}) \quad \Longleftrightarrow \quad C_A \cdot V_A = C_B \cdot V_{BE} \]
\item a) Acide chlorhydrique + soude (fort/fort) :
\[ \ce{H3O+_{(aq)} + OH-_{(aq)} -> 2H2O_{(l)}} \]
b) Acide éthanoïque + soude (faible/fort) : l'acide faible s'écrit sous forme moléculaire :
\[ \ce{CH3COOH_{(aq)} + OH-_{(aq)} -> CH3COO-_{(aq)} + H2O_{(l)}} \]
c) Ammoniac + acide chlorhydrique (faible/fort) : la base faible s'écrit sous forme moléculaire :
\[ \ce{NH3_{(aq)} + H3O+_{(aq)} -> NH4+_{(aq)} + H2O_{(l)}} \]
\item Le \cle{point de demi-équivalence} est le point de la courbe correspondant au volume versé $V_{1/2} = V_E/2$ : la moitié de l'espèce dosée a alors réagi. En ce point, les concentrations de l'acide faible et de sa base conjuguée sont égales, donc $\mathrm{pH} = \mathrm{p}K_a$ du couple mis en jeu.
\end{enumerate}
\end{corrige}

\begin{attention}[Points de vigilance -- Écriture des équations-bilan]
Ne jamais écrire \ce{HCl + NaOH} dans une équation-bilan ionique : \ce{HCl} et \ce{NaOH} sont totalement dissociés en solution, il faut écrire les ions \ce{H3O+} et \ce{OH-}. En revanche, \ce{CH3COOH} et \ce{NH3} (espèces faibles, peu dissociées) s'écrivent sous \textbf{forme moléculaire}.
\end{attention}

\courssub{Je m'entraîne}

\begin{corrige}[Exercice 5 --- Volume équivalent d'un dosage simple : acide éthanoïque / soude]
\begin{enumerate}
\item Équation-bilan de la réaction de dosage :
\[ \ce{CH3COOH_{(aq)} + OH-_{(aq)} -> CH3COO-_{(aq)} + H2O_{(l)}} \]
Cette réaction est totale ($K = K_a/K_e = 10^{-4,8}/10^{-14} = 10^{9,2} \gg 1$), rapide : elle peut servir de réaction de dosage.

\item À l'équivalence : $n(\ce{CH3COOH})_{\text{initial}} = n(\ce{OH-})_{\text{versé}}$, soit $C_A V_A = C_B V_{BE}$ :
\[ V_{BE} = \frac{C_A \cdot V_A}{C_B} = \frac{0,080 \times 15,0}{0,10} = \SI{12,0}{mL} \]
\emph{Vérification :} $n_A = 0,080 \times 15,0 \times 10^{-3} = \SI{1,20e-3}{mol}$ ; $n_B = 0,10 \times 12,0 \times 10^{-3} = \SI{1,20e-3}{mol}$. Les quantités sont bien égales.

\item \textbf{Demi-équivalence :} $V_{B,1/2} = \dfrac{V_{BE}}{2} = \SI{6,0}{mL}$, et $\mathrm{pH}_{1/2} = \mathrm{p}K_a = \textbf{4,8}$.

\textbf{Équivalence :} $\mathrm{pH}_E > 7$. \emph{Justification :} à l'équivalence, tout l'acide a été transformé en ions éthanoate \ce{CH3COO-}, base faible qui réagit avec l'eau :
\[ \ce{CH3COO- + H2O <=> CH3COOH + OH-} \]
La présence d'ions \ce{OH-} rend la solution basique.
\end{enumerate}
\end{corrige}

\begin{aretenir}[Relations clés à l'équivalence]
Pour une réaction de dosage 1:1 : $\boxed{C_A V_A = C_B V_{BE}}$.
Demi-équivalence : $\boxed{V_{1/2} = V_E/2 \quad \text{et} \quad \mathrm{pH}_{1/2} = \mathrm{p}K_a}$.
pH à l'équivalence : fort/fort $\to$ 7 ; faible/fort $\to$ $>7$ ; fort/faible $\to$ $<7$.
\end{aretenir}

\begin{corrige}[Exercice 6 --- Prévision d'une courbe de dosage : ammoniac / acide chlorhydrique]
\begin{enumerate}
\item Équation-bilan (l'ammoniac, base faible, s'écrit sous forme moléculaire) :
\[ \ce{NH3_{(aq)} + H3O+_{(aq)} -> NH4+_{(aq)} + H2O_{(l)}} \]

\item À l'équivalence : $C_B V_B = C_A V_{AE}$ :
\[ V_{AE} = \frac{C_B \cdot V_B}{C_A} = \frac{0,10 \times 25,0}{0,10} = \SI{25,0}{mL} \]

\item On verse un \textbf{acide} dans une \textbf{base} : le pH \textbf{décroît} au cours du dosage (courbe décroissante, avec un saut de pH vers le bas autour de $V_{AE} = \SI{25,0}{mL}$).

À l'équivalence, la solution ne contient pratiquement que des ions ammonium \ce{NH4+}, acide faible conjugué de \ce{NH3} :
\[ \ce{NH4+ + H2O <=> NH3 + H3O+} \]
Le milieu est donc acide : $\mathrm{pH}_E < 7$.

\emph{Estimation quantitative :} à l'équivalence, $[\ce{NH4+}] = \dfrac{C_B V_B}{V_B + V_{AE}} = \dfrac{0,10 \times 25,0}{50,0} = \SI{0,050}{mol.L^{-1}}$.
Pour un acide faible : $[\ce{H3O+}] \approx \sqrt{K_a \cdot C} = \sqrt{10^{-9,2} \times 0,050} = \sqrt{3,16 \times 10^{-11}} \approx \SI{5,6e-6}{mol.L^{-1}}$, d'où $\mathrm{pH}_E \approx 5,3 < 7$. Cohérent.

\item $\mathrm{pH}_E \approx 5,3$ : le saut de pH se situe en zone acide. Parmi les indicateurs usuels, le \textbf{rouge de méthyle} (zone 4,4 -- 6,2) est le choix classique car sa zone de virage contient $\mathrm{pH}_E = 5,3$.
\begin{itemize}
\item \textbf{Hélianthine (3,1 -- 4,4) :} vire \emph{après} l'équivalence (pH trop bas), non convenable.
\item \textbf{BBT (6,0 -- 7,6) :} vire \emph{avant} l'équivalence, non convenable.
\item \textbf{Phénolphtaléine (8,2 -- 10,0) :} vire très en amont, non convenable.
\end{itemize}
\end{enumerate}
\end{corrige}

\begin{attention}[Piège classique : indicateur pour base faible / acide fort]
Pour un dosage base faible / acide fort, le pH \textbf{diminue} et $\mathrm{pH}_E < 7$ (souvent vers 5--6). Il faut un indicateur virant en zone \textbf{acide} (rouge de méthyle, vert de bromocrésol). Choisir la phénolphtaléine « par habitude » ou l'hélianthine « car c'est acide » sont les erreurs les plus fréquentes au Bac : leurs zones de virage ne contiennent pas $\mathrm{pH}_E$.
\end{attention}

\begin{corrige}[Exercice 7 --- Exploitation d'un tableau de mesures : dosage fort/fort]
\begin{enumerate}
\item \textbf{Tracé :} sur papier millimétré, abscisses $V_B$ de 0 à \SI{20}{mL} (1 cm $\leftrightarrow$ 2 mL), ordonnées pH de 0 à 13 (1 cm $\leftrightarrow$ 1 unité). Les points (0 ; 1,5), (5 ; 1,8), (10 ; 2,2), (14 ; 2,9) s'alignent sur une branche presque horizontale ; entre 14 et \SI{16}{mL} le pH bondit de 2,9 à 11,0 (branche quasi verticale) ; les points (18 ; 11,7) et (20 ; 12,0) forment la branche finale horizontale. Relier les points par une courbe lisse et régulière (pas de segments de droite).

\item \textbf{Méthode des tangentes parallèles :} on trace deux tangentes parallèles aux branches horizontales, puis la droite parallèle équidistante qui coupe la courbe au point d'inflexion E. La projection de E donne :
\[ V_{BE} = \SI{15,0}{mL} \qquad \mathrm{pH}_E = 7,0 \]
\emph{Remarque :} le saut est symétrique autour de $V_B = \SI{15}{mL}$ (2,9 à 14 mL ; 11,0 à 16 mL), ce qui confirme $V_{BE} = \SI{15,0}{mL}$. La valeur $\mathrm{pH}_E = 7,0$ est caractéristique d'un dosage acide fort / base forte.

\item À l'équivalence : $C_A V_A = C_B V_{BE}$ :
\[ C_A = \frac{C_B \cdot V_{BE}}{V_A} = \frac{0,10 \times 15,0}{20,0} = \SI{0,075}{mol.L^{-1}} \]
\emph{Vérification :} $n_B = 0,10 \times 15,0 \times 10^{-3} = \SI{1,5e-3}{mol}$ ; $n_A = 0,075 \times 20,0 \times 10^{-3} = \SI{1,5e-3}{mol}$. Égalité parfaite.

\item \textbf{Oui}, le BBT convient : le saut de pH (de $\approx 3$ à $\approx 11$) est très étendu et \textbf{contient} la zone de virage du BBT (6,0 -- 7,6) ainsi que $\mathrm{pH}_E = 7,0$. Le virage jaune $\to$ bleu du BBT coïncidera avec l'équivalence, à une goutte près. (La phénolphtaléine conviendrait aussi ; l'hélianthine virerait trop tôt.)
\end{enumerate}
\end{corrige}

\begin{methode}[Méthode des tangentes parallèles (détermination graphique de $V_E$)]
\begin{enumerate}
\item Tracer la courbe $\mathrm{pH} = f(V)$ point par point, en lissant le trait (pas de segments).
\item Tracer une tangente à la branche montante (ou descendante) avant le saut.
\item Tracer une tangente \textbf{parallèle} à la première, sur la branche après le saut.
\item Tracer la droite \textbf{parallèle} à ces deux tangentes et \textbf{équidistante} (au milieu).
\item Cette droite coupe la courbe au point d'inflexion E : lire $V_E$ par projection sur l'axe des abscisses et $\mathrm{pH}_E$ sur l'axe des ordonnées.
\end{enumerate}
\end{methode}

\begin{corrige}[Exercice 8 --- Les trois flacons sans étiquette : identification par pH-métrie]
\begin{enumerate}
\item \textbf{Protocole :} pour chaque flacon, prélever à la pipette jaugée $V = \SI{20,0}{mL}$ de solution, verser dans un bécher, ajouter de l'eau distillée pour immerger la sonde (sans influence sur $V_E$), placer le barreau magnétique et la sonde du pH-mètre étalonné.
\begin{itemize}
\item Si le pH initial est acide (flacons 1 et 2) : doser par la soude à \SI{0,10}{mol.L^{-1}} en burette.
\item Si le pH initial est basique (flacon 3) : doser par l'acide chlorhydrique à \SI{0,10}{mol.L^{-1}} en burette.
\end{itemize}
Relever le pH après chaque ajout (par \SI{1}{mL}, puis par \SI{0,2}{mL} autour du saut) et tracer $\mathrm{pH} = f(V_{\text{titrant}})$.

\item Allures attendues ($V_E = \SI{20,0}{mL}$ dans les trois cas, car concentrations et volumes identiques) :
\begin{center}
\begin{tabularx}{\linewidth}{|l|X|X|X|}
\hline
\rowcolor{popblueL} \textbf{Flacon} & \textbf{Titrant} & \textbf{Allure} & \textbf{pH à l'équivalence} \\
\hline \ce{HCl} (fort) & \ce{NaOH} & pH croissant, grand saut vertical & $\mathrm{pH}_E = 7$ \\
\hline \ce{CH3COOH} (faible) & \ce{NaOH} & pH croissant, saut plus court, départ moins acide & $\mathrm{pH}_E > 7$ ($\approx 8,7$) \\
\hline \ce{NH3} (faible) & \ce{HCl} & pH \textbf{décroissant}, saut vers le bas & $\mathrm{pH}_E < 7$ ($\approx 5,3$) \\
\hline
\end{tabularx}
\end{center}
Le flacon d'ammoniac est immédiatement repéré par son pH initial basique et la décroissance de la courbe. Les deux acides se distinguent par le pH à l'équivalence ($=7$ contre $>7$) et par la longueur du saut.

\item À la demi-équivalence ($V = V_E/2 = \SI{10,0}{mL}$), $\mathrm{pH} = \mathrm{p}K_a$ du couple présent :
\begin{itemize}
\item si $\mathrm{pH}_{1/2} \approx 4,8$ : c'est le couple \ce{CH3COOH}/\ce{CH3COO-}, donc le flacon contient l'\textbf{acide éthanoïque} ;
\item si $\mathrm{pH}_{1/2} \approx 9,2$ : c'est le couple \ce{NH4+}/\ce{NH3}, donc le flacon contient l'\textbf{ammoniac}.
\end{itemize}
La mesure de $\mathrm{pH}_{1/2}$ constitue ainsi une identification sans ambiguïté des deux espèces faibles, et confirme celle de l'acide fort par élimination.
\end{enumerate}
\end{corrige}

\courssub{Je me prépare au Bac}

\begin{corrige}[Exercice 9 --- Type Bac : vérification de l'étiquette d'une solution d'acide chlorhydrique]
\begin{enumerate}
\item \textbf{Schéma annoté du dispositif :}

\begin{popfigure}
\begin{center}
\begin{tikzpicture}[scale=0.95, every node/.style={font=\small}]
% Support et burette
\draw[popdark, thick] (0,0) -- (0,6.2);
\draw[popdark, thick] (0,5.8) -- (2.2,5.8);
\draw[popblue, thick] (1.6,5.8) -- (1.6,2.6);
\draw[popblue, thick] (2.0,5.8) -- (2.0,2.6);
\draw[popblue, thick] (1.6,2.6) -- (1.8,2.2) -- (2.0,2.6);
\draw[popblue!60] (1.6,5.2) -- (2.0,5.2);
\node[popblue, right, align=left] at (2.3,4.6) {Burette graduée :\\ soude $C_B = \SI{0,10}{mol.L^{-1}}$};
% Bécher
\draw[popdark, thick] (0.7,0) -- (0.7,2.2);
\draw[popdark, thick] (2.9,0) -- (2.9,2.2);
\draw[popdark, thick] (0.7,0) -- (2.9,0);
\fill[popblue!15] (0.72,0.02) rectangle (2.88,1.3);
\draw[popblue!60] (0.72,1.3) -- (2.88,1.3);
\node[popdark, below, align=center] at (1.8,-0.15) {Bécher : $V_A = \SI{20,0}{mL}$ de $S_1$\\ + barreau aimanté};
% Sonde pH
\draw[poppurple, thick] (2.6,1.0) -- (2.6,3.4);
\draw[poppurple, thick, fill=poppurpleL] (2.45,3.4) rectangle (2.75,3.9);
\draw[poppurple, thick] (2.6,3.9) -- (4.3,3.9);
\draw[poppurple, thick, fill=poppurpleL] (4.3,3.5) rectangle (5.9,4.3);
\node[poppurple] at (5.1,3.9) {pH-mètre};
\node[poppurple, left] at (2.4,2.6) {sonde};
\end{tikzpicture}
\end{center}
\legende{Dispositif de dosage pH-métrique de la solution diluée $S_1$ par la soude.}
\end{popfigure}

\item \textbf{Équation-bilan :}
\[ \ce{H3O+_{(aq)} + OH-_{(aq)} -> 2H2O_{(l)}} \]
\textbf{Caractéristiques :} réaction \textbf{totale} et \textbf{rapide} (et exothermique) : elle convient à un dosage.

\item \textbf{Équivalence :} état du mélange pour lequel les réactifs ont été introduits dans les proportions stœchiométriques : $n(\ce{H3O+})_{\text{initial}} = n(\ce{OH-})_{\text{versé}}$, soit $C_A V_A = C_B V_{BE}$.

\textbf{Méthode des tangentes parallèles :} on trace deux tangentes à la courbe, parallèles entre elles, de part et d'autre du saut de pH ; on trace ensuite la droite parallèle à ces deux tangentes et équidistante : elle coupe la courbe au point d'équivalence E, dont on lit les coordonnées $(V_{BE}\,;\,\mathrm{pH}_E)$.

\item \textbf{Concentration de $S_1$ :}
\[ C_A = \frac{C_B \cdot V_{BE}}{V_A} = \frac{0,10 \times 16,0}{20,0} = \SI{0,080}{mol.L^{-1}} \]
\textbf{Concentration de $S_0$} (dilution au centième, facteur $f = 100$) :
\[ C_0 = f \times C_A = 100 \times 0,080 = \SI{8,0}{mol.L^{-1}} \]

\item \textbf{Conclusion :} l'étiquette annonce \SI{10}{mol.L^{-1}} ; le dosage donne \SI{8,0}{mol.L^{-1}}. L'écart relatif est $\dfrac{10 - 8,0}{10} = 20\ \%$, très supérieur aux incertitudes expérimentales ($\approx 1\ \%$). \textbf{L'indication de l'étiquette est donc fausse} (la solution commerciale s'est peut-être dégradée, \ce{HCl} étant volatil).

\item Le saut de pH d'un dosage fort/fort est très étendu (environ de 3 à 11) et $\mathrm{pH}_E = 7,0$ :
\begin{itemize}
\item \textbf{BBT (6,0 -- 7,6) : convient}, sa zone contient $\mathrm{pH}_E = 7,0$ ; virage jaune $\to$ bleu.
\item \textbf{Phénolphtaléine (8,2 -- 10,0) : convient}, sa zone est incluse dans le saut ; virage incolore $\to$ rose.
\item \textbf{Hélianthine (3,1 -- 4,4) : à éviter}, sa zone ne contient pas $\mathrm{pH}_E$ : elle virerait légèrement avant l'équivalence.
\end{itemize}
\end{enumerate}

\textbf{Barème indicatif (sur 10) :} schéma annoté 1,5 ; équation-bilan 1 ; caractéristiques 1 ; définition équivalence 1 ; méthode des tangentes 1 ; $C_A$ 1,5 ; $C_0$ 1 ; conclusion justifiée 1 ; choix des indicateurs 1.
\end{corrige}

\begin{attention}[Points de vigilance -- Dosage par dilution]
Ne pas oublier de \textbf{remonter la dilution} ($C_0 = 100\,C_A$) : oublier le facteur 100 est l'erreur la plus fréquente. Respecter les chiffres significatifs et convertir les volumes en litres dans les calculs de quantités de matière.
\end{attention}

\begin{corrige}[Exercice 10 --- Type Bac : degré d'acidité d'un vinaigre du marché de Mokolo]
\begin{enumerate}
\item \textbf{Pourquoi diluer :} le vinaigre commercial est trop concentré ($\approx \SI{0,8}{mol.L^{-1}}$) ; le doser directement conduirait à un volume équivalent trop grand (supérieur à la capacité de la burette) ou obligerait à utiliser une soude très concentrée, source d'imprécision. La dilution au dixième ramène le volume équivalent à une valeur commode ($\approx \SI{12}{mL}$).

\textbf{Protocole de dilution (facteur 10) :} prélever \SI{10,0}{mL} de vinaigre à la \textbf{pipette jaugée}, les introduire dans une \textbf{fiole jaugée de \SI{100,0}{mL}}, compléter avec de l'eau distillée jusqu'au trait de jauge, boucher et agiter pour homogénéiser. On obtient la solution $S$ : $C_S = C/10$.

\item Équation-bilan de la réaction de dosage :
\[ \ce{CH3COOH_{(aq)} + OH-_{(aq)} -> CH3COO-_{(aq)} + H2O_{(l)}} \]

\item \textbf{Tableau d'avancement à l'équivalence} (quantités en mol) :
\begin{center}
\begin{tabular}{|l|c|c|c|c|}
\hline
\rowcolor{popblueL} Équation & \ce{CH3COOH} & \ce{OH-} & \ce{CH3COO-} & \ce{H2O} \\
\hline État initial & $C_A V_A$ & $C_B V_{BE}$ & 0 & excès \\
\hline Équivalence & 0 & 0 & $C_A V_A$ & excès \\
\hline
\end{tabular}
\end{center}
À l'équivalence : $C_A V_A = C_B V_{BE}$ :
\[ C_A = \frac{C_B \cdot V_{BE}}{V_A} = \frac{0,020 \times 12,0}{10,0} = \SI{0,024}{mol.L^{-1}} \]
Concentration du vinaigre commercial :
\[ C = 10 \times C_A = \SI{0,24}{mol.L^{-1}} \]

\item Masse d'acide éthanoïque dans \SI{1}{L} de vinaigre :
\[ m = C \times M = 0,24 \times 60 = \SI{14,4}{g.L^{-1}} \]
\SI{1}{L} de vinaigre a une masse $m_V = \rho \times V = 1,0 \times 1000 = \SI{1000}{g}$.
Par proportionnalité, la masse d'acide dans \SI{100}{g} de vinaigre est :
\[ m_{100} = 14,4 \times \frac{100}{1000} = \SI{1,44}{g} \]
\textbf{Degré d'acidité : $1,44^\circ \approx 1,4^\circ$.}

\item \textbf{Conclusion :} l'étiquette annonce $5^\circ$ ; le dosage donne $1,4^\circ$, soit plus de trois fois moins. \textbf{L'étiquette est inexacte} : le vinaigre est beaucoup trop pauvre en acide éthanoïque (dilution frauduleuse ou acétification incomplète).

\item $\mathrm{pH}_E = 8,6 > 7$ : à l'équivalence, la solution contient des ions éthanoate \ce{CH3COO-}, base faible qui capte des protons de l'eau : $\ce{CH3COO- + H2O <=> CH3COOH + OH-}$ ; le milieu est donc basique.

\textbf{Demi-équivalence :} $V_{B,1/2} = \dfrac{12,0}{2} = \SI{6,0}{mL}$ et $\mathrm{pH}_{1/2} = \mathrm{p}K_a = \textbf{4,8}$.

\textbf{Indicateur :} la \textbf{phénolphtaléine} (zone 8,2 -- 10,0), car sa zone de virage contient $\mathrm{pH}_E = 8,6$ ; virage incolore $\to$ rose. Le BBT (6,0 -- 7,6) virerait avant l'équivalence et l'hélianthine bien trop tôt.
\end{enumerate}

\textbf{Barème indicatif (sur 12) :} dilution justifiée + protocole 2 ; équation-bilan 1 ; tableau d'avancement 1 ; $C_A$ 1,5 ; $C$ 1 ; masse dans 1 L 1 ; degré d'acidité 1,5 ; conclusion 1 ; $\mathrm{pH}_E > 7$ justifié 1 ; $\mathrm{pH}_{1/2}$ et indicateur 1.
\end{corrige}

\begin{attention}[Points de vigilance -- Degré d'acidité]
Le degré d'acidité est une \textbf{masse d'acide pour \SI{100}{g} de solution}, pas une concentration molaire : il faut impérativement utiliser la masse volumique $\rho$ pour passer du litre aux 100 g. Ne pas confondre $C_A$ (solution diluée) et $C$ (vinaigre commercial) : le facteur de dilution 10 doit apparaître explicitement.
\end{attention}

\begin{corrige}[Exercice 11 --- Type Bac : dosage de l'acide éthanoïque et détermination du pKa]
\begin{enumerate}
\item Équation-bilan de la réaction de dosage :
\[ \ce{CH3COOH_{(aq)} + OH-_{(aq)} -> CH3COO-_{(aq)} + H2O_{(l)}} \]

\item \textbf{Lecture graphique :} la méthode des tangentes parallèles (ou le repérage du point d'inflexion au milieu du saut) donne le point E :
\[ \boxed{V_{BE} = \SI{15,0}{mL} \quad ; \quad \mathrm{pH}_E = 8,7} \]
\textbf{Cohérence :} oui. L'acide éthanoïque est un \textbf{acide faible} ; à l'équivalence, la solution contient sa base conjuguée \ce{CH3COO-}, qui réagit avec l'eau ($\ce{CH3COO- + H2O <=> CH3COOH + OH-}$) : le pH à l'équivalence doit être \textbf{supérieur à 7}, ce que confirme $\mathrm{pH}_E = 8,7$.

\item À l'équivalence : $C_A V_A = C_B V_{BE}$ :
\[ C_A = \frac{C_B \cdot V_{BE}}{V_A} = \frac{0,10 \times 15,0}{20,0} = \SI{0,075}{mol.L^{-1}} \]
\emph{Vérification :} $n_B = 0,10 \times 15,0 \times 10^{-3} = \SI{1,5e-3}{mol}$ ; $n_A = 0,075 \times 20,0 \times 10^{-3} = \SI{1,5e-3}{mol}$. Égalité exacte.

\item \textbf{Définition :} le point de demi-équivalence D correspond au volume versé $V_{B,1/2} = V_{BE}/2$ ; la moitié de l'acide initial a alors réagi et $[\ce{CH3COOH}] = [\ce{CH3COO-}]$.

\textbf{Lecture graphique :} D a pour coordonnées :
\[ V_{B,1/2} = \SI{7,5}{mL} \quad ; \quad \mathrm{pH}_{1/2} = 4,8 \]
\textbf{Déduction :} à la demi-équivalence, $\mathrm{pH} = \mathrm{p}K_a$, donc :
\[ \mathrm{p}K_a(\ce{CH3COOH}/\ce{CH3COO-}) = 4,8 \]
\textbf{Comparaison :} la valeur déterminée graphiquement (4,8) coïncide avec la valeur tabulée (4,8) : la détermination expérimentale est excellente.

\item $\mathrm{pH}_E = 8,7$ appartient à la zone de virage de la \textbf{phénolphtaléine} (8,2 -- 10,0) : cet indicateur est donc bien choisi, son virage signale l'équivalence à une goutte près.

\textbf{Changement de teinte :} la solution, \textbf{incolore} avant l'équivalence (milieu acide), devient \textbf{rose} (rose persistant) dès le léger excès de soude, à l'équivalence.

\item Volume de soude à la demi-équivalence :
\[ V_{B,1/2} = \frac{V_{BE}}{2} = \frac{15,0}{2} = \SI{7,5}{mL} \]
\textbf{Cohérence graphique :} le point D de la courbe a bien pour abscisse \SI{7,5}{mL} et pour ordonnée 4,8 : calcul et graphique sont parfaitement cohérents.
\end{enumerate}

\textbf{Barème indicatif (sur 10) :} équation-bilan 1 ; lecture de E 1,5 ; cohérence $\mathrm{pH}_E > 7$ 1 ; $C_A$ 1,5 ; définition de D 1 ; lecture de D et $\mathrm{p}K_a$ 1,5 ; comparaison 0,5 ; indicateur + teinte 1,5 ; $V_{B,1/2}$ et vérification 0,5.
\end{corrige}

\begin{attention}[Trois pièges classiques au Bac]
\begin{enumerate}
\item Lire $V_{BE}$ au \textbf{point d'inflexion} (milieu du saut) et non au sommet de la courbe.
\item Ne pas confondre demi-équivalence ($\mathrm{pH} = \mathrm{p}K_a$) et équivalence ($\mathrm{pH} > 7$ pour un acide faible).
\item Justifier le choix de l'indicateur par l'\textbf{appartenance de $\mathrm{pH}_E$ à la zone de virage}, et non par la couleur.
\end{enumerate}
\end{attention}

\newpage

\coursec{Fiche bilan}

\begin{popfigure}
\begin{tikzpicture}[
  mindmap/.style={text=white, font=\bfseries\small, align=center},
  central/.style={circle, fill=popdark, text=white, font=\bfseries, align=center, minimum size=2.8cm},
  branche/.style={rounded corners=6pt, fill=#1, text=white, font=\small\bfseries, align=center, inner sep=4pt, minimum width=3.2cm},
  lien/.style={line width=1.2pt, draw=popmuted},
  annot/.style={font=\scriptsize, text=popdark, align=center}
]
\node[central] (C) {Réactions\\acide/base\\et dosages};
\node[branche=popblue, above left=1.2cm and 2.5cm of C] (AFBF) {Acide fort\\-- base forte};
\node[branche=popgreen, left=2.7cm of C] (AFaBF) {Acide faible\\-- base forte};
\node[branche=poporange, below left=1.2cm and 2.5cm of C] (AFBFa) {Acide fort\\-- base faible};
\node[branche=poppurple, above right=1.2cm and 2.5cm of C] (EQ) {Équivalence\\$C_A V_A = C_B V_{BE}$};
\node[branche=popgold, right=2.7cm of C] (GR) {Courbes\\$\mathrm{pH}=f(V)$};
\node[branche=poppink, below right=1.2cm and 2.5cm of C] (IC) {Indicateurs\\colorés};
\node[branche=popblue!70, below=1.8cm of C] (AP) {Applications\\vinaigre, étiquettes};
\draw[lien] (C) -- (AFBF); \draw[lien] (C) -- (AFaBF); \draw[lien] (C) -- (AFBFa);
\draw[lien] (C) -- (EQ); \draw[lien] (C) -- (GR); \draw[lien] (C) -- (IC); \draw[lien] (C) -- (AP);
\node[annot, above=0.15cm of AFBF] {$\mathrm{pH}_E=7$};
\node[annot, above=0.15cm of AFaBF, xshift=-0.4cm] {$\mathrm{pH}_E>7$ ; demi-éq. : $\mathrm{pH}=\mathrm{p}K_a$};
\node[annot, below=0.15cm of AFBFa] {$\mathrm{pH}_E<7$};
\node[annot, above=0.15cm of GR, xshift=0.3cm] {tangentes parallèles};
\node[annot, below=0.15cm of IC] {zone de virage $\ni \mathrm{pH}_E$};
\node[annot, below=0.15cm of AP] {dilution, degré d'acidité};
\end{tikzpicture}
\legende{Carte mentale de la leçon : les trois types de dosage, l'équivalence, les courbes, les indicateurs et les applications.}
\end{popfigure}

\begin{aretenir}[L'essentiel de la leçon]
\textbf{Définitions clés}
\begin{itemize}
\item \cle{Dosage} (ou titrage) : détermination de la concentration inconnue d'une solution (\cle{solution dosée}) grâce à une réaction \textbf{rapide, totale et unique} avec un réactif de concentration connue (\cle{solution titrante}).
\item \cle{Équivalence acido-basique} : état du mélange où l'acide et la base ont été introduits dans les proportions stœchiométriques exactes. Pour une réaction 1/1 : $n_{\text{acide}} = n_{\text{base}}$, soit $C_A V_A = C_B V_{BE}$.
\item \cle{Point de demi-équivalence} (dosage d'un acide faible par une base forte) : volume $V_B = \dfrac{V_{BE}}{2}$ ; on y a $\mathrm{pH} = \mathrm{p}K_a$ du couple acide/base.
\item \cle{Indicateur coloré} : couple acide/base $\ce{AH}/\ce{A-}$ dont les deux formes ont des couleurs distinctes ; sa \cle{zone de virage} (intervalle de pH où le changement de couleur est visible) doit contenir le $\mathrm{pH}$ à l'équivalence.
\end{itemize}

\textbf{Équations-bilan des réactions de dosage (à connaître par c\oe ur)}
\begin{itemize}
\item Acide fort + base forte : \[ \ce{H3O+ + HO- -> 2 H2O} \quad \text{(totale, rapide, exothermique)} \]
\item Acide faible + base forte (ex. acide éthanoïque) : \[ \ce{CH3COOH + HO- -> CH3COO- + H2O} \]
\item Acide fort + base faible (ex. ammoniac) : \[ \ce{H3O+ + NH3 -> H2O + NH4+} \]
\end{itemize}

\textbf{Formules et résultats fondamentaux}
\begin{center}
\begin{tabularx}{\linewidth}{|l|X|X|}
\hline
\rowcolor{popblueL} \textbf{Type de dosage} & \textbf{Relation à l'équivalence} & \textbf{pH à l'équivalence \& justification} \\
\hline
Acide fort / Base forte & $C_A V_A = C_B V_{BE}$ & $\mathrm{pH}_E = 7$ (formation d'eau seulement) \\
\hline
Acide faible / Base forte & $C_A V_A = C_B V_{BE}$ & $\mathrm{pH}_E > 7$ (solution basique : anion \ce{CH3COO-} hydrolysé) \\
\hline
Acide fort / Base faible & $C_A V_A = C_B V_{AE}$ & $\mathrm{pH}_E < 7$ (solution acide : cation \ce{NH4+} hydrolysé) \\
\hline
\end{tabularx}
\end{center}
\begin{itemize}
\item \textbf{Demi-équivalence (acide faible)} : $\mathrm{pH}_{1/2} = \mathrm{p}K_a$ ; lecture graphique directe du $\mathrm{p}K_a$ sur la courbe $\mathrm{pH}=f(V_B)$.
\item \textbf{Dilution commerciale} : $C_{\text{commerce}} = C_{\text{dosée}} \times f$ avec $f = \dfrac{V_{\text{flacon jaugé}}}{V_{\text{pipette}}}$ (facteur de dilution).
\item \textbf{Degré d'acidité du vinaigre} : masse d'acide éthanoïque (en g) pour \SI{100}{mL} de vinaigre (ou pour \SI{100}{g} selon étiquette). $t = \dfrac{C_A \times M_{\ce{CH3COOH}} \times 10}{d}$ (si $C_A$ en \si{mol.L^{-1}}, $M=\SI{60,05}{g.mol^{-1}}$, $d$ masse volumique en \si{g.mL^{-1}}).
\end{itemize}

\textbf{Méthodes express}
\begin{itemize}
\item \cle{Méthode des tangentes parallèles} : tracer deux tangentes parallèles aux branches montantes/descendantes encadrant le saut de pH, puis la droite équidistante ; son intersection avec la courbe donne le point d'équivalence $E(V_E\,;\,\mathrm{pH}_E)$.
\item \cle{Méthode de la dérivée} (notion) : le maximum de $\dfrac{\mathrm{d}\mathrm{pH}}{\mathrm{d}V}$ correspond à l'équivalence.
\item \cle{Choix de l'indicateur coloré} :
  \begin{center}
  \begin{tabular}{|c|c|c|}\hline
  \rowcolor{popblueL} \textbf{Indicateur} & \textbf{Zone de virage} & \textbf{Cas d'utilisation} \\ \hline
  Hélianthine (orange de méthyle) & $3,1$--$4,4$ & $\mathrm{pH}_E < 7$ (fort/faible) \\ \hline
  BBT (bleu de bromothymol) & $6,0$--$7,6$ & $\mathrm{pH}_E \approx 7$ (fort/fort) \\ \hline
  Phénolphtaléine & $8,2$--$10,0$ & $\mathrm{pH}_E > 7$ (faible/forte) \\ \hline
  \end{tabular}
  \end{center}
\item \cle{Protocole opératoire} : pipette jaugée (prise d'essai) $\rightarrow$ erlenmeyer + quelques gouttes d'indicateur (ou électrode pH-métrique) ; burette graduée (solution titrante) ; agitation magnétique ; approche du virage goutte à goutte ; lecture du volume à l'équivalence (virage persistant $\approx \SI{30}{\second}$ ou $\mathrm{pH}_E$).
\end{itemize}
\end{aretenir}

\begin{aretenir}[Checklist avant le Bac]
\begin{itemize}
\item[$\square$] J'écris sans erreur les trois équations-bilan : \ce{H3O+ + HO-}, \ce{CH3COOH + HO-}, \ce{H3O+ + NH3}.
\item[$\square$] Je justifie qu'une réaction de dosage doit être \textbf{totale, rapide et unique}.
\item[$\square$] Je définis l'équivalence et j'applique la relation $C_A V_A = C_B V_{BE}$ (volumes en \textbf{même unité}, concentrations en \si{mol.L^{-1}}).
\item[$\square$] Je connais et \textbf{justifie} les pH à l'équivalence : $=7$ (fort/fort), $>7$ (faible/forte $\rightarrow$ base conjuguée), $<7$ (fort/faible $\rightarrow$ acide conjugué).
\item[$\square$] Je repère le point de demi-équivalence ($V_{BE}/2$) et j'en déduis $\mathrm{p}K_a = \mathrm{pH}_{1/2}$.
\item[$\square$] Je détermine $V_E$ graphiquement (tangentes parallèles) ou par la dérivée (maximum).
\item[$\square$] Je choisis l'indicateur dont la zone de virage contient $\mathrm{pH}_E$ (hélianthine, BBT, phénolphtaléine).
\item[$\square$] Je construis un tableau d'avancement pour calculer le pH en tout point (avant, à, après équivalence).
\item[$\square$] Je gère une dilution avant dosage : $C_0 = C_1 \times f$ et je remonte à la concentration commerciale.
\item[$\square$] Je calcule le degré d'acidité (g/100 mL ou g/100 g) et je vérifie la conformité d'une étiquette (vinaigre, lessive).
\item[$\square$] Je décris le protocole complet : matériel (pipette jaugée, burette, erlenmeyer, agitateur), manipulation (rinçage, prise d'essai, titrage goutte à goutte près de l'équivalence), lecture.
\item[$\square$] Je vérifie la cohérence de mes résultats (ordre de grandeur, pH final logique, unités).
\end{itemize}
\end{aretenir}

\findecours

\end{document}
